% Primitives d'une fonction continue — Mathématiques Tle D — Cours signé OnBuch+
% Programme officiel MINESEC. Compiler avec : tectonic M05-primitives.tex
\newcommand{\DOCMATIERE}{Mathématiques}
\newcommand{\DOCNIVEAU}{Tle D}
\newcommand{\DOCMODULE}{Relations et opérations fondamentales dans l'ensemble des nombres réels}
\newcommand{\DOCLECON}{05}
\newcommand{\DOCTITRE}{Primitives d'une fonction continue}
% OnBuch+ — préambule « cours signés » ENRICHI (compilé avec tectonic / XeLaTeX)
% Système visuel « Pop » d'OnBuch+ : fond crème (jamais blanc), bordures encre
% épaisses, ombres « dures » décalées, titres Archivo Black en capitales.
% Version MATHÉMATIQUES : graphes (pgfplots/tikz), boîtes pédagogiques
% (définition, propriété, méthode, attention, le savais-tu, exercice résolu,
% exercice, TP, bilan), page de garde et signature OnBuch+ ancrée en bas.
%
% Macros attendues dans main.tex AVANT \input{preamble} :
%   \DOCMATIERE  \DOCNIVEAU  \DOCMODULE  \DOCLECON (numéro)  \DOCTITRE
\documentclass[11pt]{article}

\usepackage{fontspec}
\usepackage[french]{babel}
\usepackage[a4paper,margin=2cm,top=3.4cm,bottom=2.8cm,headheight=1.4cm,footskip=1.3cm]{geometry}
\usepackage{amsmath,amssymb,amsfonts}
\usepackage{mathtools} % \xmapsto, \coloneqq... souvent utilisés par les modèles
\usepackage{cancel} % simplifications barrées
% \abs{z} (module, valeur absolue) et \norm{u} (norme), utilisés par les modèles
\providecommand{\abs}{}\let\abs\relax\DeclarePairedDelimiter{\abs}{\lvert}{\rvert}
\providecommand{\norm}{}\let\norm\relax\DeclarePairedDelimiter{\norm}{\lVert}{\rVert}
\usepackage[table]{xcolor} % [table] : fournit aussi \rowcolors (alternance de lignes)
\usepackage{pagecolor}
\usepackage{tikz}
\usetikzlibrary{arrows.meta,positioning,calc,shapes.geometric,shapes.misc,shapes.arrows,decorations.pathmorphing,decorations.markings,patterns,shadows,mindmap,trees,fit,backgrounds,matrix,intersections}
\usepackage{pgfplots}
\usepackage{pgf-pie} % diagrammes circulaires \pie (statistiques)
\pgfplotsset{compat=1.18}
\usepgfplotslibrary{fillbetween,groupplots} % aires entre courbes (calcul intégral)
\usepackage{enumitem}
\usepackage{fancyhdr}
% [most] charge toutes les bibliothèques tcolorbox (listings, minted, raster,
% documentation...) et gonfle inutilement la mémoire TeX sur un document long
% (~300 boîtes) : on ne charge que ce qui est réellement utilisé.
\usepackage[skins,breakable]{tcolorbox}
\usepackage{graphicx}
\usepackage{colortbl}
\usepackage{booktabs}
\usepackage{tabularx}
\usepackage{diagbox} % case d'en-tête diagonale des tableaux à double entrée
% Colonnes extensibles centrée / gauche / droite (C, L, R), souvent utilisées par les modèles
\newcolumntype{C}{>{\centering\arraybackslash}X}
\newcolumntype{L}{>{\raggedright\arraybackslash}X}
\newcolumntype{R}{>{\raggedleft\arraybackslash}X}
\usepackage{array}
\usepackage{multicol}
\usepackage{siunitx}
% parse-numbers=false : accepte aussi \SI{2,5 \times 10^{-3}}{...}
\sisetup{locale=FR,output-decimal-marker={,},group-minimum-digits=4,parse-numbers=false}
\DeclareSIUnit\ohm{\ensuremath{\symup{\Omega}}} % U+2126 absent de la police
\usepackage{qrcode}
% \degree : commande fréquemment utilisée par les modèles pour un angle ou une
% température en dehors de \SI{}{} (non fournie par les paquets chargés).
\providecommand{\degree}{^{\circ}}
% \qed (fin de démonstration), utilisé par les modèles sans amsthm
\providecommand{\qed}{\hfill$\square$}
% \barre : double barre des valeurs interdites dans un tableau de variations (tkz-tab)
\providecommand{\barre}{\ensuremath{\|}}
% environnement proof (amsthm non chargé) utilisé par les modèles
\ifdefined\proof\else\newenvironment{proof}[1][Démonstration]{\par\noindent\textit{#1.}\ }{\qed\par}\fi
\usepackage{lastpage}
\usepackage{hyperref}
\hypersetup{hidelinks}

% ============================================================
% Palette « Pop » OnBuch+ — mêmes hex que la classe P (plus_ui.dart)
% ============================================================
\definecolor{poporange}{HTML}{F59321}
\definecolor{popdark}{HTML}{E07A0C}
\definecolor{popcream}{HTML}{FFF6E8}
\definecolor{popink}{HTML}{1C1714}
\definecolor{poppurple}{HTML}{7A5AE0}
\definecolor{popgreen}{HTML}{2E9E6A}
\definecolor{poppink}{HTML}{E0457B}
\definecolor{popblue}{HTML}{2F7DE1}
\definecolor{popgold}{HTML}{C98A12}
\definecolor{popmuted}{HTML}{8A7F76}
\definecolor{popline}{HTML}{EADFD2}
\definecolor{popgray}{HTML}{8A7F76}
\colorlet{popgrey}{popgray}
% Alias tolérants (le modèle nomme parfois les couleurs différemment)
\colorlet{popwhite}{white}
\colorlet{popblack}{popink}
\colorlet{popred}{poppink}
\colorlet{popyellow}{popgold}
% Teintes claires pour fonds de tableaux / zones
\colorlet{popblueL}{popblue!10!white}
\colorlet{poporangeL}{poporange!14!white}
\colorlet{popgreenL}{popgreen!12!white}
\colorlet{poppurpleL}{poppurple!12!white}
\colorlet{poppinkL}{poppink!12!white}
\colorlet{popgoldL}{popgold!14!white}

\pagecolor{popcream}

% ============================================================
% Typographie
% ============================================================
\defaultfontfeatures{Ligatures=TeX}
\setmainfont{PlusJakartaSans-Medium.ttf}[Path=../../../fonts/,BoldFont=PlusJakartaSans-Bold.ttf]
% Maths en sans-serif (Fira Math) avec chiffres et lettres droites de Plus
% Jakarta, pour que les formules se fondent dans le texte.
\usepackage[math-style=ISO,bold-style=ISO]{unicode-math}
\setmathfont{FiraMath-Regular.otf}[Path=../../../fonts/]
\setmathfont{PlusJakartaSans-Medium.ttf}[Path=../../../fonts/,range={up/num,up/{Latin,latin},\mathup}]
\usepackage{icomma} % virgule décimale sans espace en mode math
% Caractères mathématiques Unicode absents des polices de texte (Plus Jakarta,
% Archivo Black) : rendus par leur équivalent mathématique, en texte comme en
% formule — sinon ils s'affichent en carré vide (ex. « Arithmétique dans ℤ »).
\usepackage{newunicodechar}
\newunicodechar{ℝ}{\ensuremath{\mathbb{R}}}
\newunicodechar{ℤ}{\ensuremath{\mathbb{Z}}}
\newunicodechar{ℂ}{\ensuremath{\mathbb{C}}}
\newunicodechar{ℕ}{\ensuremath{\mathbb{N}}}
\newunicodechar{ℚ}{\ensuremath{\mathbb{Q}}}
\newunicodechar{𝔻}{\ensuremath{\mathbb{D}}}
\newunicodechar{α}{\ensuremath{\alpha}}
\newunicodechar{β}{\ensuremath{\beta}}
\newunicodechar{γ}{\ensuremath{\gamma}}
\newunicodechar{δ}{\ensuremath{\delta}}
\newunicodechar{ε}{\ensuremath{\varepsilon}}
\newunicodechar{θ}{\ensuremath{\theta}}
\newunicodechar{λ}{\ensuremath{\lambda}}
\newunicodechar{μ}{\ensuremath{\mu}}
\newunicodechar{σ}{\ensuremath{\sigma}}
\newunicodechar{τ}{\ensuremath{\tau}}
\newunicodechar{φ}{\ensuremath{\varphi}}
\newunicodechar{ψ}{\ensuremath{\psi}}
\newunicodechar{ω}{\ensuremath{\omega}}
\newunicodechar{Σ}{\ensuremath{\Sigma}}
% majuscules produites par \MakeUppercase dans les titres (α-aminés -> Α)
\newunicodechar{Α}{\ensuremath{\alpha}}
\newunicodechar{Β}{\ensuremath{\beta}}
\newunicodechar{Γ}{\ensuremath{\Gamma}}
\newunicodechar{Δ}{\ensuremath{\Delta}}
\newunicodechar{Ε}{\ensuremath{\varepsilon}}
\newunicodechar{Θ}{\ensuremath{\Theta}}
\newunicodechar{Λ}{\ensuremath{\Lambda}}
\newunicodechar{Μ}{\ensuremath{\mu}}
\newunicodechar{Τ}{\ensuremath{\tau}}
\newunicodechar{Φ}{\ensuremath{\Phi}}
\newunicodechar{Ψ}{\ensuremath{\Psi}}
\newunicodechar{Ω}{\ensuremath{\Omega}}
\newunicodechar{↦}{\ensuremath{\mapsto}}
\newunicodechar{∈}{\ensuremath{\in}}
\newunicodechar{∉}{\ensuremath{\notin}}
\newunicodechar{∧}{\ensuremath{\wedge}}
\newunicodechar{⇔}{\ensuremath{\Leftrightarrow}}
\newunicodechar{⟺}{\ensuremath{\Longleftrightarrow}}
\newunicodechar{⇒}{\ensuremath{\Rightarrow}}
\newunicodechar{⟹}{\ensuremath{\Longrightarrow}}
\newunicodechar{∘}{\ensuremath{\circ}}
\newunicodechar{∪}{\ensuremath{\cup}}
\newunicodechar{∩}{\ensuremath{\cap}}
\newunicodechar{≡}{\ensuremath{\equiv}}
\newunicodechar{⊂}{\ensuremath{\subset}}
\newunicodechar{⊥}{\ensuremath{\perp}}
\newunicodechar{∃}{\ensuremath{\exists}}
\newunicodechar{∀}{\ensuremath{\forall}}
\newunicodechar{∛}{\ensuremath{\sqrt[3]{\,}}}
\newunicodechar{∜}{\ensuremath{\sqrt[4]{\,}}}
\newunicodechar{⃗}{\ensuremath{{}^{\to}}}
\newunicodechar{ⁿ}{\ifmmode^{n}\else\textsuperscript{\ensuremath{n}}\fi}
\newunicodechar{ˣ}{\ifmmode^{x}\else\textsuperscript{\ensuremath{x}}\fi}
\newunicodechar{ᵇ}{\ifmmode^{b}\else\textsuperscript{\ensuremath{b}}\fi}
\newunicodechar{ʳ}{\ifmmode^{r}\else\textsuperscript{\ensuremath{r}}\fi}
\newunicodechar{ᵘ}{\ifmmode^{u}\else\textsuperscript{\ensuremath{u}}\fi}
\newunicodechar{ʸ}{\ifmmode^{y}\else\textsuperscript{\ensuremath{y}}\fi}
\newunicodechar{ᵃ}{\ifmmode^{a}\else\textsuperscript{\ensuremath{a}}\fi}
\newunicodechar{ᵉ}{\ifmmode^{e}\else\textsuperscript{\ensuremath{e}}\fi}
\newunicodechar{ᵏ}{\ifmmode^{k}\else\textsuperscript{\ensuremath{k}}\fi}
\newunicodechar{ᵖ}{\ifmmode^{p}\else\textsuperscript{\ensuremath{p}}\fi}
\newunicodechar{ᴺ}{\ifmmode^{N}\else\textsuperscript{\ensuremath{N}}\fi}
\newunicodechar{ᶜ}{\ifmmode^{c}\else\textsuperscript{\ensuremath{c}}\fi}
\newunicodechar{ʰ}{\ifmmode^{h}\else\textsuperscript{\ensuremath{h}}\fi}
\newunicodechar{ᵠ}{\ifmmode^{q}\else\textsuperscript{\ensuremath{q}}\fi}
\newunicodechar{ₙ}{\ifmmode_{n}\else\textsubscript{\ensuremath{n}}\fi}
\newunicodechar{ₐ}{\ifmmode_{a}\else\textsubscript{\ensuremath{a}}\fi}
\newunicodechar{ᵢ}{\ifmmode_{i}\else\textsubscript{\ensuremath{i}}\fi}
\newunicodechar{ᵣ}{\ifmmode_{r}\else\textsubscript{\ensuremath{r}}\fi}
\newunicodechar{ₖ}{\ifmmode_{k}\else\textsubscript{\ensuremath{k}}\fi}
\newunicodechar{ᵦ}{\ifmmode_{\beta}\else\textsubscript{\ensuremath{\beta}}\fi}
\newunicodechar{ₓ}{\ifmmode_{x}\else\textsubscript{\ensuremath{x}}\fi}
\newfontfamily\popdisplay{ArchivoBlack-Regular.ttf}[Path=../../../fonts/]
\newfontfamily\poplogo{SpaceGrotesk-Bold.ttf}[Path=../../../fonts/]
\newfontfamily\popsemi{PlusJakartaSans-SemiBold.ttf}[Path=../../../fonts/]
\newfontfamily\popxbold{PlusJakartaSans-ExtraBold.ttf}[Path=../../../fonts/]
\newfontfamily\popxboldls{PlusJakartaSans-ExtraBold.ttf}[Path=../../../fonts/,LetterSpace=3.0]

\setlist[enumerate]{leftmargin=1.7em,itemsep=5pt,topsep=4pt}
\setlist[itemize]{leftmargin=1.7em,itemsep=4pt,topsep=4pt,label={\color{poporange}\textbullet}}
\setlength{\parindent}{0pt}
\setlength{\parskip}{5pt}
\renewcommand{\arraystretch}{1.25}

% pgfplots : style Pop par défaut
\pgfplotsset{
  every axis/.append style={axis line style={popink,line width=1pt},tick style={popink},
    grid style={popline,line width=0.5pt},label style={font=\small},tick label style={font=\footnotesize}},
  popcurve/.style={poporange,line width=1.8pt,no markers,smooth},
}
% Même style disponible dans un \draw TikZ simple (hors pgfplots)
\tikzset{popcurve/.style={poporange,line width=1.8pt}}

% ============================================================
% Pill / squiggle
% ============================================================
\newcommand{\poppill}[3]{%
  \tikz[baseline=-0.55ex]\node[draw=popink,line width=1.2pt,rounded corners=2.6pt,%
    fill=#2,inner xsep=6.5pt,inner ysep=3pt]{%
    \popxboldls\fontsize{7}{7}\selectfont\color{#3}\MakeUppercase{#1}};%
}
\newcommand{\popsquiggle}[1]{%
  \tikz[baseline=-0.4ex]{
    \draw[#1,line width=2.6pt,line cap=round]
      plot [smooth] coordinates {(0,0.09) (0.34,0.01) (0.68,0.21) (1.02,0.01) (1.36,0.10)};
  }%
}
% Mot-clé surligné (marqueur orange sous le texte)
\newcommand{\cle}[1]{{\popxbold\color{popdark}#1}}

% ============================================================
% En-tête / pied
% ============================================================
\pagestyle{fancy}
\fancyhf{}
\renewcommand{\headrulewidth}{0pt}
\renewcommand{\footrulewidth}{0pt}
\lhead{\raisebox{-0.15ex}{\poplogo\fontsize{13}{13}\selectfont\color{popink}OnBuch\color{poporange}+}}
\rhead{\poppill{\DOCMATIERE\ \textbullet\ \DOCNIVEAU\ \textbullet\ Leçon \DOCLECON}{white}{popink}}
\lfoot{\popxbold\fontsize{7.6}{7.6}\selectfont\color{popmuted}\MakeUppercase{Cours signé OnBuch+}\ \textbullet\ {\popsemi\DOCTITRE}}
\rfoot{\tikz[baseline=-0.6ex]\node[rounded corners=8.5pt,fill=popink,minimum height=17pt,minimum width=17pt,inner xsep=5pt,inner sep=0]{\popxbold\fontsize{8}{8}\selectfont\color{popcream}\thepage\,/\,\pageref*{LastPage}};}

% ============================================================
% Titres
% ============================================================
\newcommand{\chaptitre}[2]{%
  \begin{tcolorbox}[enhanced,colback=poporange,colframe=popink,boxrule=2.6pt,arc=11pt,
      drop shadow=popink,left=15pt,right=15pt,top=12pt,bottom=11pt,before skip=3pt,after skip=18pt]
    \poppill{#2}{popink}{popcream}\par\vspace{9pt}
    {\raggedright\hyphenpenalty=10000\exhyphenpenalty=10000\popdisplay\fontsize{22}{24}\selectfont\color{popcream}\MakeUppercase{#1}\par}
  \end{tcolorbox}
}
% Section numérotée : pastille orange avec le numéro + titre Archivo
\newcounter{popsec}
\newcommand{\coursec}[1]{%
  \stepcounter{popsec}\setcounter{popsub}{0}%
  \par\vspace{16pt}\noindent
  \begin{minipage}{\linewidth}
  \tikz[baseline=-0.7ex]\node[draw=popink,line width=1.6pt,fill=poporange,rounded corners=4pt,minimum size=22pt,inner sep=0]{\popdisplay\fontsize{12}{12}\selectfont\color{popcream}\thepopsec};%
  \hspace{8pt}{\popdisplay\fontsize{15}{17}\selectfont\color{popink}\MakeUppercase{#1}}\par
  \vspace{4pt}\noindent{\color{poporange}\rule{36pt}{3.6pt}}
  \end{minipage}\par\nopagebreak\vspace{8pt}
}
\newcounter{popsub}
\newcommand{\courssub}[1]{%
  \stepcounter{popsub}%
  \par\vspace{9pt}\noindent{\popxbold\fontsize{11.5}{13}\selectfont\color{popink}{\color{poporange}\thepopsec.\thepopsub}\hspace{6pt}#1}\par\nopagebreak\vspace{4pt}
}

% ============================================================
% Boîtes pédagogiques — même recette : fond blanc, bordure épaisse colorée,
% ombre dure encre, badge plein. Titre optionnel [#1].
% ============================================================
\tcbset{popbase/.style={breakable,enhanced,colback=white,boxrule=2.3pt,arc=9pt,
  shadow={3pt}{-3pt}{0pt}{popink},left=14pt,right=14pt,top=11pt,bottom=11pt,before skip=11pt,after skip=15pt,
  fontupper=\popsemi\fontsize{10.6}{14.5}\selectfont\color{popink},
  fontlower=\popsemi\fontsize{10.6}{14.5}\selectfont\color{popink}}}
% Coche dessinée (le glyphe ✓ n'existe pas dans les polices du document)
\newcommand{\popcheck}[1]{\tikz[baseline=-0.5ex]\draw[#1,line width=1.6pt,line cap=round,line join=round](-0.13,0.0)--(-0.04,-0.09)--(0.13,0.1);}
\newcommand{\popboxhead}[3]{\poppill{#1}{#2}{white}\ifx\relax#3\relax\else\hspace{6pt}{\popxbold\fontsize{10.6}{12}\selectfont\color{popink}#3}\fi\par\vspace{7pt}}

\newtcolorbox{aretenir}[1][]{popbase,colframe=popgreen,before upper={\popboxhead{À retenir}{popgreen}{#1}}}
\newtcolorbox{exemplebox}[1][]{popbase,colframe=poporange,before upper={\popboxhead{Exemple}{poporange}{#1}}}
\newtcolorbox{definition}[1][]{popbase,colframe=popblue,before upper={\popboxhead{Définition}{popblue}{#1}}}
\newtcolorbox{propriete}[1][]{popbase,colframe=poppurple,before upper={\popboxhead{Propriété}{poppurple}{#1}}}
\newtcolorbox{methode}[1][]{popbase,colframe=popgold,before upper={\popboxhead{Méthode}{popgold}{#1}}}
\newtcolorbox{attention}[1][]{popbase,colframe=poppink,before upper={\popboxhead{Attention}{poppink}{#1}}}
\newtcolorbox{experience}[1][]{popbase,colframe=popblue,colback=popblueL,before upper={\popboxhead{Expérience}{popblue}{#1}}}
% « Le savais-tu ? » : carte violette pleine (contexte, culture, Cameroun)
\newtcolorbox{savaistu}[1][]{popbase,colback=poppurple,colframe=popink,
  fontupper=\popsemi\fontsize{10.4}{14.2}\selectfont\color{white},
  before upper={\poppill{Le savais-tu\,?}{white}{poppurple}\ifx\relax#1\relax\else\hspace{6pt}{\popxbold\color{white}#1}\fi\par\vspace{7pt}}}
% Exercice résolu : énoncé en haut, \tcblower puis la solution sur fond crème
\newcounter{popexo}\newcounter{popexr}
\newtcolorbox{exoresolu}[1][]{popbase,colframe=poporange,colbacklower=popcream,
  segmentation style={popink,line width=1.2pt,dashed},
  before upper={\stepcounter{popexr}\popboxhead{Exercice résolu \thepopexr}{poporange}{#1}},
  before lower={\poppill{Solution}{popink}{popcream}\par\vspace{6pt}}}
% Exercice (non résolu en place) — la correction est dans la partie Corrigés
\newtcolorbox{exercice}[1][]{popbase,colframe=popink,
  before upper={\stepcounter{popexo}\popboxhead{Exercice \thepopexo}{popink}{#1}}}
% Corrigé
\newtcolorbox{corrige}[1][]{popbase,colframe=popgreen,colback=popgreenL,
  before upper={\popboxhead{Corrigé}{popgreen}{#1}}}
% Objectifs / prérequis / bilan
\newtcolorbox{objectifs}{popbase,colframe=popink,colback=poporangeL,
  before upper={\popboxhead{Objectifs de la leçon}{popink}{}}}
\newtcolorbox{prerequis}{popbase,colframe=popink,colback=popblueL,
  before upper={\popboxhead{Prérequis}{popblue}{}}}
\newtcolorbox{bilan}{popbase,colframe=popink,colback=white,boxrule=2.8pt,
  before upper={\popboxhead{Fiche bilan}{poporange}{L'essentiel à connaître pour l'examen}}}
% Figure encadrée avec légende
\newtcolorbox{popfigure}[1][]{enhanced,breakable=false,colback=white,colframe=popline,boxrule=1.4pt,arc=8pt,
  left=10pt,right=10pt,top=10pt,bottom=8pt,before skip=10pt,after skip=12pt,halign=center,
  lower separated=false,halign lower=center,
  fontlower=\popsemi\fontsize{9}{11.5}\selectfont\color{popmuted},
  #1}
\newcommand{\legende}[1]{\par\vspace{6pt}{\popsemi\fontsize{9}{11.5}\selectfont\color{popmuted}#1}}

% ============================================================
% Signature OnBuch+
% ============================================================
\newcommand{\poplogoinline}[1]{{\poplogo\fontsize{#1}{#1}\selectfont\color{popink}OnBuch\color{poporange}+}}
\newcommand{\signatureonbuch}{%
  \begin{tcolorbox}[enhanced,colback=popink,colframe=popink,boxrule=2.2pt,arc=10pt,
      drop shadow=poporange,left=14pt,right=14pt,top=10pt,bottom=10pt,before skip=0pt,after skip=0pt]
    \tikz[baseline=-0.7ex]\node[circle,fill=popgreen,draw=popcream,line width=1.3pt,minimum size=22pt,inner sep=0]{\popcheck{white}};%
    \hspace{9pt}\begin{minipage}[c]{0.62\linewidth}
      {\popxbold\fontsize{12}{13}\selectfont\color{popcream}Signé\ \textbullet\ {\poplogo OnBuch\color{poporange}+}}\par\vspace{2pt}
      {\raggedright\popsemi\fontsize{8}{10.5}\selectfont\color{popline}\MakeUppercase{Équipe pédagogique OnBuch+}\\\MakeUppercase{Conforme au programme officiel MINESEC}\par}
    \end{minipage}\hfill
    {\poplogo\fontsize{22}{22}\selectfont\color{popcream}OnBuch\color{poporange}+}
  \end{tcolorbox}%
}

% ============================================================
% Page de garde
% #1 titre · #2 matière · #3 niveau · #4 module · #5 n° leçon · #6 durée ·
% #7 description / objectifs (texte ou itemize)
% ============================================================
\newcommand{\pagegarde}[7]{%
  \thispagestyle{empty}
  \newgeometry{a4paper,margin=1.7cm,top=1.6cm,bottom=1.4cm}%
  \begin{tikzpicture}[remember picture,overlay]
    \fill[poporange!18] ([xshift=-3.2cm,yshift=-3.6cm]current page.north east) circle (4.6cm);
    \fill[poppurple!14] ([xshift=2.4cm,yshift=7.6cm]current page.south west) circle (3.8cm);
  \end{tikzpicture}%
  {\poplogo\fontsize{30}{30}\selectfont\color{popink}OnBuch\color{poporange}+}\hfill
  \raisebox{0.6ex}{\poppill{Cours signé \textbullet\ Premium}{popink}{popcream}}\par
  \vspace{1pt}\popsquiggle{poppurple}\par
  \vspace{8pt}
  \begin{tcolorbox}[enhanced,colback=poporange,colframe=popink,boxrule=2.8pt,arc=13pt,
      drop shadow=popink,left=18pt,right=18pt,top=12pt,bottom=13pt,after skip=10pt]
    \poppill{#2 \textbullet\ #3}{popink}{popcream}\hspace{5pt}\poppill{#4}{white}{popink}\par\vspace{8pt}
    {\popdisplay\fontsize{14}{14}\selectfont\color{popink}LEÇON #5}\par\vspace{4pt}
    {\raggedright\hyphenpenalty=10000\exhyphenpenalty=10000\popdisplay\fontsize{25}{27}\selectfont\color{popcream}\MakeUppercase{#1}\par}
    \vspace{8pt}
    {\popxbold\fontsize{9.5}{9.5}\selectfont\color{popink}\MakeUppercase{Durée conseillée : #6}}
  \end{tcolorbox}
  \begin{tcolorbox}[enhanced,colback=white,colframe=popink,boxrule=2.2pt,arc=10pt,
      drop shadow=popink,left=16pt,right=16pt,top=10pt,bottom=10pt,after skip=8pt]
    \poppill{Dans cette leçon}{poppurple}{white}\par\vspace{5pt}
    \popsemi\fontsize{9.3}{12.6}\selectfont\color{popink}#7
  \end{tcolorbox}
  \noindent\begin{minipage}[t]{0.487\linewidth}
    \begin{tcolorbox}[enhanced,colback=popgreen,colframe=popink,boxrule=2.2pt,arc=10pt,
        drop shadow=popink,left=11pt,right=11pt,top=10pt,bottom=10pt,height=3.1cm,valign=center]
      \tcbox[colback=white,colframe=popink,boxrule=1.3pt,arc=5pt,boxsep=0pt,nobeforeafter,
        left=3pt,right=3pt,top=3pt,bottom=3pt,valign=center]{\qrcode[height=1.9cm]{https://play.google.com/store/apps/details?id=com.luvvix.onbuch&hl=fr}}%
      \hspace{8pt}\begin{minipage}[c]{0.5\linewidth}
        {\popdisplay\fontsize{11}{12.5}\selectfont\color{white}\MakeUppercase{Télécharge OnBuch}}\par\vspace{4pt}
        {\raggedright\popsemi\fontsize{8.4}{11}\selectfont\color{white}Gratuit \textbullet\ Google Play\\Tuteur IA, annales, concours, résultats\par}
      \end{minipage}
    \end{tcolorbox}
  \end{minipage}\hfill
  \begin{minipage}[t]{0.487\linewidth}
    \begin{tcolorbox}[enhanced,colback=poppurple,colframe=popink,boxrule=2.2pt,arc=10pt,
        drop shadow=popink,left=11pt,right=11pt,top=10pt,bottom=10pt,height=3.1cm,valign=center]
      \tikz[baseline=-0.7ex]\node[circle,fill=white,draw=popink,line width=1.3pt,minimum size=1.6cm,inner sep=0]{\popdisplay\fontsize{24}{24}\selectfont\color{poppurple}+};%
      \hspace{8pt}\begin{minipage}[c]{0.6\linewidth}
        {\popdisplay\fontsize{11}{12.5}\selectfont\color{white}\MakeUppercase{Passe à OnBuch+}}\par\vspace{4pt}
        {\raggedright\popsemi\fontsize{8.4}{11}\selectfont\color{white}Tous les cours signés, Tuteur IA illimité, fiches PDF — onglet OnBuch+ de l'app\par}
      \end{minipage}
    \end{tcolorbox}
  \end{minipage}
  \vspace{8pt}
  \popcontenu
  \vfill
  \signatureonbuch
  \restoregeometry
  \newpage
}

% Rangée « Ce cours contient » (page de garde)
\newcommand{\poptile}[3]{% #1 couleur #2 gros texte #3 légende
  \begin{tcolorbox}[enhanced,colback=#1,colframe=popink,boxrule=2pt,arc=9pt,shadow={2.5pt}{-2.5pt}{0pt}{popink},
      left=6pt,right=6pt,top=9pt,bottom=9pt,halign=center,valign=center,height=1.75cm,width=\linewidth,nobeforeafter]
    {\popdisplay\fontsize{12.5}{13}\selectfont\color{white}\MakeUppercase{#2}}\par\vspace{3pt}
    {\centering\popsemi\fontsize{7.8}{9.5}\selectfont\color{white}#3\par}
  \end{tcolorbox}}
\newcommand{\popcontenu}{%
  \par\noindent{\popxboldls\fontsize{7.5}{7.5}\selectfont\color{popmuted}CE COURS SIGNÉ CONTIENT}\par\vspace{7pt}
  \noindent\begin{minipage}[t]{0.235\linewidth}\poptile{popblue}{Cours}{complet, illustré, pas à pas}\end{minipage}\hfill
  \begin{minipage}[t]{0.235\linewidth}\poptile{poporange}{Méthodes}{et exercices résolus}\end{minipage}\hfill
  \begin{minipage}[t]{0.235\linewidth}\poptile{poppink}{Exercices}{type examen + corrigés}\end{minipage}\hfill
  \begin{minipage}[t]{0.235\linewidth}\poptile{popgreen}{Fiche bilan}{l'essentiel à retenir}\end{minipage}\par}

% Sommaire
\newtcolorbox{sommaire}{enhanced,colback=white,colframe=popink,boxrule=2.2pt,arc=10pt,
  drop shadow=popink,left=18pt,right=18pt,top=13pt,bottom=13pt,before skip=6pt,after skip=20pt,
  before upper={\poppill{Sommaire}{poppurple}{white}\par\vspace{9pt}\popsemi\fontsize{10.6}{14.5}\selectfont\color{popink}}}

% Signature de fin de cours — poussée tout en bas de la dernière page
\newcommand{\findecours}{%
  \par\vfill\nopagebreak
  \begin{center}\popsquiggle{poporange}\end{center}\vspace{4pt}
  \signatureonbuch
}
\AtBeginDocument{\def\llbracket{\mathopen{[\mkern-3.5mu[}}\def\rrbracket{\mathclose{]\mkern-3.5mu]}}} % intervalles d'entiers (glyphe ⟦ absent de la police)
% Glyphes absents de Fira Math : redéfinis à partir de glyphes disponibles
\AtBeginDocument{%
  \def\setminus{\mathbin{\backslash}}\let\smallsetminus\setminus
  \def\vdots{\mathord{\vbox{\baselineskip3.2pt\lineskiplimit0pt\kern4pt\hbox{.}\hbox{.}\hbox{.}}}}%
  \def\ddots{\mathinner{\mkern1mu\raise6.5pt\hbox{.}\mkern2mu\raise3.6pt\hbox{.}\mkern2mu\raise0.7pt\hbox{.}\mkern1mu}}%
  \def\triangle{\mathord{\scalebox{0.9}{$\symup{\Delta}$}}}%
  \def\nabla{\mathord{\raisebox{\depth}{\scalebox{1}[-1]{$\symup{\Delta}$}}}}%
  \def\checkmark{\popcheck{popdark}}%
  \def\star{\mathbin{\ast}}\def\bigstar{\mathord{\ast}}%
  \def\diamond{\mathbin{\scalebox{0.6}{$\Diamond$}}}%
}

%% FIGFIX-BEGIN v3 — correctifs globaux des figures (docs/onbuch-generation/tools/figfix)
\usepackage{adjustbox}\usepackage{etoolbox}
% toute figure TikZ/pgfplots plus large que la ligne est réduite à la largeur disponible
\BeforeBeginEnvironment{tikzpicture}{\begin{adjustbox}{max width=\linewidth}}
\AfterEndEnvironment{tikzpicture}{\end{adjustbox}}
% légendes pgfplots lisibles par-dessus les courbes
\pgfplotsset{every axis/.append style={legend style={fill=white,fill opacity=0.92,text opacity=1,draw=black!15,inner sep=3pt}}}
% étiquettes d'arêtes (syntaxe edge["..."]) avec fond blanc
\tikzset{every edge quotes/.append style={fill=white,inner sep=1.2pt}}
% bulles de cartes mentales : texte à la ligne dans la bulle, bulle un peu plus grande
\tikzset{every concept/.append style={text width=2.0cm,align=center,minimum size=2.3cm,inner sep=2pt}}
%% FIGFIX-END
\begin{document}

\pagegarde{Primitives d'une fonction continue}{Mathématiques}{Tle D}{Relations et opérations fondamentales dans l'ensemble des nombres réels}{05}{5 h}{Cette leçon introduit la notion de primitive d'une fonction continue sur un intervalle, opération inverse de la dérivation. Elle fournit les tableaux de primitives usuelles et les techniques de reconnaissance des formes composées (u′·uʳ, u′/u, u′eᵘ) indispensables pour le calcul intégral et les équations différentielles du programme de Terminale.
\vspace{4pt}
\begin{multicols}{2}\small
\begin{itemize}
\item À la fin de cette leçon, je sais définir une primitive d'une fonction continue sur un intervalle et vérifier qu'une fonction F est une primitive d'une fonction f
\item À la fin de cette leçon, je sais justifier que deux primitives d'une même fonction sur un intervalle diffèrent d'une constante
\item À la fin de cette leçon, je sais déterminer l'ensemble des primitives d'une fonction usuelle (puissances, cosinus, sinus)
\item À la fin de cette leçon, je sais déterminer l'unique primitive d'une fonction prenant une valeur donnée en un point
\item À la fin de cette leçon, je sais utiliser la linéarité : primitive d'une somme et d'un produit par un réel
\item À la fin de cette leçon, je sais reconnaître la forme u′·uʳ pour primitiver une fonction composée
\end{itemize}\end{multicols}}

\begin{sommaire}
\begin{multicols}{2}
\begin{enumerate}
\item Notion de primitive et existence
\item Ensemble des primitives ; condition initiale
\item Primitives des fonctions usuelles
\item Primitives des formes composées : u′·uʳ et a/(cx+d)ʳ
\item Primitives de cosⁿx et sinⁿx
\item Compléments utiles au Bac : u′/u et u′eᵘ ; synthèse méthodologique
\item Activité d'intégration
\item Méthodes et exercices résolus
\item Exercices
\item Corrigés des exercices
\item Fiche bilan
\end{enumerate}
\end{multicols}
\end{sommaire}

\begin{objectifs}
\begin{itemize}
\item À la fin de cette leçon, je sais définir une primitive d'une fonction continue sur un intervalle et vérifier qu'une fonction F est une primitive d'une fonction f
\item À la fin de cette leçon, je sais justifier que deux primitives d'une même fonction sur un intervalle diffèrent d'une constante
\item À la fin de cette leçon, je sais déterminer l'ensemble des primitives d'une fonction usuelle (puissances, cosinus, sinus)
\item À la fin de cette leçon, je sais déterminer l'unique primitive d'une fonction prenant une valeur donnée en un point
\item À la fin de cette leçon, je sais utiliser la linéarité : primitive d'une somme et d'un produit par un réel
\item À la fin de cette leçon, je sais reconnaître la forme u′·uʳ pour primitiver une fonction composée
\item À la fin de cette leçon, je sais reconnaître les formes u′/u et u′eᵘ (complément utile au Bac)
\item À la fin de cette leçon, je sais linéariser ou utiliser les formules trigonométriques pour primitiver cosⁿx et sinⁿx (n = 2, 3)
\item À la fin de cette leçon, je sais rédiger rigoureusement la recherche d'une primitive en précisant l'intervalle de validité
\end{itemize}
\end{objectifs}

\begin{prerequis}
\begin{itemize}
\item Dérivation : dérivées des fonctions usuelles et dérivée d'une fonction composée (Première)
\item Formules de trigonométrie : cos²x + sin²x = 1, cos(2x) = 2cos²x − 1, duplication
\item Fonctions puissances à exposant rationnel et règles de calcul sur les exposants
\item Notion de fonction continue sur un intervalle (début d'année de Terminale)
\item Fonctions ln et exponentielle et leurs dérivées (leçons précédentes de Terminale)
\end{itemize}
\end{prerequis}

\newpage

\coursec{Notion de primitive et existence}

\courssub{Situation introductive : le problème inverse de la dérivation}

Imagine un chauffeur de taxi-brousse sur l'axe Douala--Yaoundé. À chaque instant du trajet, son compteur de vitesse fonctionne parfaitement : il connaît sa \emph{vitesse} $v(t)$. Malheureusement, son GPS est en panne et le compteur kilométrique est cassé. Question : peut-il retrouver la \emph{distance parcourue} $d(t)$ depuis le départ ?

De même, à Bafoussam, une entreprise d'eau minérale surveille le \emph{rythme de remplissage} de son réservoir, c'est-à-dire le débit instantané en litres par minute. Mais le responsable de production veut connaître la \emph{quantité totale d'eau stockée} à un instant donné.

Dans les deux situations, on connaît le \emph{taux de variation} d'une grandeur, c'est-à-dire sa \cle{dérivée}, et on cherche la grandeur elle-même. Or la vitesse est la dérivée de la distance parcourue, et le débit est la dérivée du volume stocké. Autrement dit :

\begin{center}
\emph{On connaît $f$, et on cherche une fonction $F$ telle que $F' = f$.}
\end{center}

C'est exactement le \textbf{problème inverse de la dérivation}. Ce problème, d'une importance capitale en mathématiques et dans toutes les sciences (physique, économie, biologie), est résolu par la notion de \cle{primitive}, que nous allons définir rigoureusement dans cette leçon.

\begin{popfigure}
\begin{center}
\begin{tikzpicture}[scale=1.1, every node/.style={font=\small}]
\node[draw=popblue, fill=popblueL, rounded corners=4pt, minimum width=3.2cm, minimum height=1cm] (F) at (0,0) {$F$ \ (la fonction)};
\node[draw=poporange, fill=poporangeL, rounded corners=4pt, minimum width=3.2cm, minimum height=1cm] (f) at (7,0) {$f$ \ (sa dérivée)};
\draw[-{Stealth[length=3mm]}, very thick, popdark] (F) to[bend left=25] node[midway, above]{\textbf{dérivation} : $F \longmapsto F'$} (f);
\draw[-{Stealth[length=3mm]}, very thick, popgreen] (f) to[bend left=25] node[midway, below]{\textbf{primitivation} : $f \longmapsto F$} (F);
\end{tikzpicture}
\end{center}
\legende{La dérivation et la primitivation sont deux opérations réciproques : dériver $F$ donne $f$ ; primitiver $f$ redonne $F$ (à une constante près).}
\end{popfigure}

\courssub{Définition d'une primitive}

Reprenons l'exemple du taxi-brousse. Si la distance parcourue est modélisée par une fonction $d$ du temps, alors la vitesse affichée au compteur est $v(t) = d'(t)$. Retrouver la distance à partir de la vitesse, c'est donc retrouver une fonction \emph{dérivable} dont la dérivée est connue. Cette intuition conduit naturellement à la définition suivante.

\begin{definition}[Primitive d'une fonction]
Soit $f$ une fonction définie sur un intervalle $I$. On appelle \cle{primitive} de $f$ sur $I$ toute fonction $F$ \textbf{dérivable sur $I$} telle que :
\[ \text{pour tout } x \in I, \quad F'(x) = f(x). \]
On dit aussi que $F$ est une \emph{antidérivée} de $f$.
\end{definition}

Deux points de cette définition méritent d'être soulignés :

\begin{itemize}
\item La fonction $F$ doit être \textbf{dérivable} sur $I$ tout entier : c'est une condition préalable, sans laquelle parler de $F'$ n'a pas de sens.
\item L'égalité $F'(x) = f(x)$ doit être vraie \textbf{pour tout} $x$ de l'intervalle $I$, et pas seulement en quelques points.
\end{itemize}

\begin{attention}[Primitive ou dérivée : ne pas inverser les rôles]
L'erreur la plus fréquente est de confondre les deux directions. Dire \og $F$ est une primitive de $f$ \fg{} signifie $F' = f$, et \textbf{pas} $f' = F$. Ainsi, primitiver la fonction $x \mapsto x^2$ donne $x \mapsto \dfrac{x^3}{3} + c$, et \textbf{surtout pas} $2x$ (qui est sa \emph{dérivée}). En cas de doute, dérive toujours ton résultat : tu dois retomber sur la fonction de départ.
\end{attention}

\courssub{Vérifier qu'une fonction est une primitive d'une autre}

La définition fournit immédiatement une méthode de vérification, très souvent demandée au Baccalauréat sous la forme \og montrer que $F$ est une primitive de $f$ sur $I$ \fg.

\begin{methode}[Vérifier que $F$ est une primitive de $f$ sur $I$]
\begin{enumerate}
\item Justifier que $F$ est \textbf{dérivable} sur $I$ (somme, produit, quotient ou composée de fonctions dérivables).
\item \textbf{Calculer} $F'(x)$ pour tout $x$ de $I$.
\item \textbf{Conclure} : si $F'(x) = f(x)$ pour tout $x \in I$, alors $F$ est une primitive de $f$ sur $I$.
\end{enumerate}
\end{methode}

\begin{exemplebox}[Vérification directe]
Vérifions que $F(x) = \dfrac{x^3}{3}$ est une primitive de $f(x) = x^2$ sur $\mathbb{R}$.

La fonction $F$ est une fonction polynôme, donc dérivable sur $\mathbb{R}$, et pour tout réel $x$ :
\[ F'(x) = \frac{1}{3} \times 3x^2 = x^2 = f(x). \]
Donc $F$ est bien une primitive de $f$ sur $\mathbb{R}$.
\end{exemplebox}

\begin{exoresolu}[Une primitive d'une fonction composée]
\textbf{Énoncé.} Vérifier que la fonction $F$ définie par $F(x) = \dfrac{(2x+1)^3}{6}$ est une primitive de la fonction $f$ définie par $f(x) = (2x+1)^2$ sur $\mathbb{R}$.

\tcblower
\textbf{Solution détaillée.}

\textbf{Étape 1 : dérivabilité.} La fonction $x \mapsto 2x+1$ est dérivable sur $\mathbb{R}$, donc son cube aussi, et $F$ est dérivable sur $\mathbb{R}$.

\textbf{Étape 2 : calcul de la dérivée.} On utilise la formule de dérivation d'une composée $(u^3)' = 3u'u^2$ avec $u(x) = 2x+1$, donc $u'(x) = 2$ :
\begin{align*}
F'(x) &= \frac{1}{6} \times 3 \times 2 \times (2x+1)^2 \\
&= \frac{6}{6}(2x+1)^2 \\
&= (2x+1)^2.
\end{align*}

\textbf{Étape 3 : conclusion.} Pour tout réel $x$, $F'(x) = (2x+1)^2 = f(x)$. Donc $F$ est bien une primitive de $f$ sur $\mathbb{R}$.
\end{exoresolu}

\begin{exoresolu}[Un contre-exemple : attention à la dérivabilité]
\textbf{Énoncé.} La fonction $F$ définie par $F(x) = |x|$ vérifie-t-elle $F'(x) = f(x)$ avec $f(x) = 1$ si $x > 0$ et $f(x) = -1$ si $x < 0$ ? $F$ est-elle une primitive de $f$ sur $\mathbb{R}$ ?

\tcblower
\textbf{Solution détaillée.} Pour $x > 0$, $F(x) = x$ donc $F'(x) = 1 = f(x)$ ; pour $x < 0$, $F(x) = -x$ donc $F'(x) = -1 = f(x)$. L'égalité $F' = f$ est vraie sur $\mathbb{R}^*$.

Cependant, $F$ n'est \textbf{pas dérivable en $0$} (la courbe présente un point anguleux). La première condition de la définition n'est donc pas satisfaite sur $\mathbb{R}$ tout entier : $F$ n'est \textbf{pas} une primitive de $f$ sur $\mathbb{R}$. En revanche, $F$ est une primitive de $f$ sur $]0\,;+\infty[$ et sur $]-\infty\,;0[$ séparément.

Ce contre-exemple montre que l'hypothèse de dérivabilité sur l'intervalle \emph{entier} est indispensable.
\end{exoresolu}

\courssub{Existence des primitives : le théorème fondamental}

Nous savons maintenant \emph{vérifier} qu'une fonction est une primitive d'une autre. Mais une question plus profonde se pose : étant donnée une fonction $f$, existe-t-il \emph{toujours} une fonction $F$ dont elle est la dérivée ?

La réponse est non en général : la fonction $f$ de l'exercice résolu précédent (qui vaut $-1$ puis $1$) n'admet aucune primitive sur $\mathbb{R}$. En revanche, une hypothèse naturelle suffit à garantir l'existence : la \cle{continuité}.

\begin{propriete}[Théorème d'existence des primitives (admis)]
Toute fonction \textbf{continue} sur un intervalle $I$ admet des primitives sur $I$.
\end{propriete}

Ce théorème est \textbf{admis} en Terminale : sa démonstration repose sur la notion d'intégrale et sera donnée dans la leçon sur l'intégration, où l'on montrera que si $f$ est continue sur $I$ et $a \in I$, alors la fonction $x \mapsto \displaystyle\int_a^x f(t)\,\mathrm{d}t$ est une primitive de $f$ sur $I$. C'est le fameux \emph{théorème fondamental du calcul intégral}.

\begin{attention}[Continuité : une condition suffisante, à ne pas oublier]
Retiens bien le sens du théorème : \og continue $\Rightarrow$ admet des primitives \fg. Toutes les fonctions usuelles du programme (polynômes, fonctions rationnelles sur leur domaine, cosinus, sinus, exponentielle, logarithme) sont continues sur chaque intervalle de leur ensemble de définition, donc admettent des primitives sur ces intervalles. En pratique, avant de chercher une primitive, vérifie toujours sur quel intervalle la fonction est continue.
\end{attention}

\courssub{Premiers exemples de primitives}

Cherchons des primitives de quelques fonctions de référence, en \og remontant \fg{} les formules de dérivation connues.

\begin{exemplebox}[Trois exemples fondamentaux]
\textbf{1. Primitive de $f(x) = x^2$ sur $\mathbb{R}$.} On sait que la dérivée de $x^3$ est $3x^2$. En divisant par $3$, on obtient : $F(x) = \dfrac{x^3}{3}$ est une primitive de $f$ sur $\mathbb{R}$. Vérification : $F'(x) = \dfrac{3x^2}{3} = x^2$.

\textbf{2. Primitive de $f(x) = \dfrac{1}{x^2}$ sur $]0\,;+\infty[$.} On sait que la dérivée de $\dfrac{1}{x}$ est $-\dfrac{1}{x^2}$. Donc $F(x) = -\dfrac{1}{x}$ est une primitive de $f$ sur $]0\,;+\infty[$ (et aussi sur $]-\infty\,;0[$, mais \emph{pas} sur $\mathbb{R}^*$ considéré comme un seul ensemble, car ce n'est pas un intervalle). Vérification : $F'(x) = -\left(-\dfrac{1}{x^2}\right) = \dfrac{1}{x^2}$.

\textbf{3. Primitive de $f(x) = \cos x$ sur $\mathbb{R}$.} On sait que $(\sin x)' = \cos x$. Donc $F(x) = \sin x$ est une primitive de $\cos$ sur $\mathbb{R}$.
\end{exemplebox}

\begin{attention}[Le domaine doit être un intervalle]
La fonction $x \mapsto \dfrac{1}{x^2}$ est continue sur $\mathbb{R}^* = ]-\infty\,;0[ \cup ]0\,;+\infty[$, mais $\mathbb{R}^*$ n'est \textbf{pas un intervalle}. Le théorème d'existence s'applique séparément sur $]-\infty\,;0[$ et sur $]0\,;+\infty[$. Cette subtilité aura son importance dans la section suivante, lorsque nous parlerons de la constante d'intégration.
\end{attention}

\courssub{Non-unicité d'une primitive : une famille de courbes parallèles}

Revenons à notre chauffeur de taxi-brousse. Connaître sa vitesse à chaque instant ne lui dit \emph{pas} où il se trouve de façon absolue : deux véhicules roulant exactement au même rythme, mais partis l'un de Douala et l'autre d'Akonolinga, ont la même fonction vitesse, mais des fonctions distance qui diffèrent d'une constante (l'écart initial).

Mathématiquement, cela se traduit ainsi : si $F$ est une primitive de $f$ sur $I$, alors pour toute constante réelle $c$, la fonction $G = F + c$ est aussi dérivable sur $I$ et
\[ G'(x) = F'(x) + 0 = f(x), \]
donc $G$ est \emph{aussi} une primitive de $f$ sur $I$. Une fonction qui admet une primitive en admet donc \textbf{une infinité}, toutes obtenues en ajoutant une constante.

\begin{aretenir}[Une primitive n'est jamais unique]
Si $F$ est une primitive de $f$ sur un intervalle $I$, alors pour tout réel $c$, la fonction $F + c$ est encore une primitive de $f$ sur $I$. Nous montrerons dans la section suivante que, réciproquement, \emph{toutes} les primitives de $f$ sur $I$ sont de la forme $F + c$.
\end{aretenir}

Géométriquement, ajouter une constante revient à translater verticalement la courbe. Toutes les courbes des primitives d'une même fonction sont donc des translatées les unes des autres : en un point d'abscisse donnée, leurs tangentes ont le même coefficient directeur $f(x)$, donc sont \textbf{parallèles}.

\begin{popfigure}
\begin{center}
\begin{tikzpicture}
\begin{axis}[
  width=0.85\linewidth, height=7cm,
  axis lines=middle,
  xlabel={$x$}, ylabel={$y$},
  xmin=-2.6, xmax=2.6, ymin=-4.5, ymax=7.5,
  xtick={-2,-1,1,2}, ytick={-3,-2,2,4,6},
  legend style={at={(0.02,0.98)}, anchor=north west, font=\scriptsize},
  clip=false
]
\addplot[popcurve, popblue, domain=-2.4:2.4, samples=100]{x^2};
\addlegendentry{$F_1(x)=x^2$}
\addplot[popcurve, poporange, domain=-2.2:2.2, samples=100]{x^2+2};
\addlegendentry{$F_2(x)=x^2+2$}
\addplot[popcurve, popgreen, domain=-1.7:1.7, samples=100]{x^2-3};
\addlegendentry{$F_3(x)=x^2-3$}
% tangentes paralleles en x = 1, pente 2
\addplot[popdark, thick, dashed, domain=0.3:1.7, samples=2]{2*x-1};
\addplot[popdark, thick, dashed, domain=0.3:1.7, samples=2]{2*x+1};
\addplot[popdark, thick, dashed, domain=0.6:1.4, samples=2]{2*x-4};
\addplot[only marks, mark=*, popink, mark size=1.6pt] coordinates {(1,1) (1,3) (1,-2)};
\end{axis}
\end{tikzpicture}
\end{center}
\legende{Trois primitives de $f(x)=2x$ : les courbes de $F_1$, $F_2$, $F_3$ se déduisent l'une de l'autre par translation verticale. En $x=1$, les trois tangentes (en pointillés) ont le même coefficient directeur $f(1)=2$ : elles sont parallèles.}
\end{popfigure}

\begin{savaistu}[D'où vient le mot \og primitive \fg{} ?]
Le mot \emph{primitive} vient du latin \emph{primitivus}, qui signifie \og premier \fg{} ou \og originel \fg. Primitiver une fonction, c'est remonter à la fonction \og première \fg, celle qui existait \emph{avant} l'opération de dérivation. C'est exactement ce que fait notre chauffeur de taxi-brousse : la distance parcourue est la grandeur \og première \fg, la vitesse n'en est que l'image dérivée. La notion a été développée au XVII\textsuperscript{e} siècle par Newton et Leibniz, dans le cadre de la naissance du calcul intégral : Newton cherchait précisément à retrouver la position d'un astre à partir de sa vitesse, un problème de... primitive !
\end{savaistu}

\begin{exercice}[À toi de jouer]
\begin{enumerate}
\item Vérifier que $F(x) = \dfrac{x^4}{4} - 2x$ est une primitive de $f(x) = x^3 - 2$ sur $\mathbb{R}$.
\item Déterminer une primitive de $g(x) = \sin x$ sur $\mathbb{R}$, puis une primitive de $h(x) = 3x^2$ sur $\mathbb{R}$.
\item Le réservoir de l'entreprise de Bafoussam se remplit à un débit constant de $50$ litres par minute. Donner une fonction $V(t)$ exprimant le volume d'eau stocké en fonction du temps $t$ (en minutes), sachant que le réservoir contenait $200$ litres à l'instant $t = 0$.
\end{enumerate}
\end{exercice}

\begin{corrige}[Exercice : À toi de jouer]
\textbf{1.} $F$ est une fonction polynôme, dérivable sur $\mathbb{R}$, et $F'(x) = \dfrac{4x^3}{4} - 2 = x^3 - 2 = f(x)$ pour tout réel $x$. Donc $F$ est une primitive de $f$ sur $\mathbb{R}$.

\textbf{2.} Comme $(-\cos x)' = \sin x$, une primitive de $g$ est $G(x) = -\cos x$. Comme $(x^3)' = 3x^2$, une primitive de $h$ est $H(x) = x^3$.

\textbf{3.} Le débit est la dérivée du volume : $V'(t) = 50$. Une primitive de la fonction constante $50$ est $V(t) = 50t + c$. La condition $V(0) = 200$ donne $c = 200$, d'où $V(t) = 50t + 200$ litres. On retrouve l'idée clé de la leçon : connaître la dérivée ne suffit pas, il faut une information supplémentaire (la valeur initiale) pour déterminer la fonction de façon unique.
\end{corrige}

Dans cette section, nous avons posé les fondations : une primitive de $f$ est une fonction $F$ dérivable telle que $F' = f$ ; toute fonction continue sur un intervalle en admet ; et dès qu'il en existe une, il en existe une infinité, différant d'une constante. La section suivante précisera ce dernier point et montrera comment une \emph{condition initiale} permet de sélectionner une unique primitive.

\coursec{Ensemble des primitives ; condition initiale}

Dans la section précédente, nous avons découvert ce qu'est une \cle{primitive} d'une fonction continue $f$ sur un intervalle $I$ : c'est une fonction $F$ dérivable sur $I$ dont la dérivée redonne $f$, c'est-à-dire $F' = f$. Nous avons aussi admis que toute fonction continue sur un intervalle admet au moins une primitive. Mais une question surgit immédiatement : cette primitive est-elle \emph{unique} ? Prenons un exemple simple. La fonction $F(x) = x^2$ est une primitive de $f(x) = 2x$ sur $\mathbb{R}$, puisque $F'(x) = 2x$. Mais la fonction $G(x) = x^2 + 5$ vérifie aussi $G'(x) = 2x$ ! Et $H(x) = x^2 - 17$ également. Il semble donc qu'une fonction possède une \emph{infinité} de primitives. Cette section répond à deux questions fondamentales : comment décrire \textbf{toutes} les primitives d'une fonction, et comment en sélectionner une \emph{seule} grâce à une information supplémentaire appelée \cle{condition initiale}.

\courssub{De la constatation à la propriété : les translatées d'une primitive sont des primitives}

Commençons par une observation géométrique. Si l'on ajoute une constante $c$ à une fonction $F$, le graphe de $F + c$ s'obtient en translatant verticalement celui de $F$. Or une translation verticale ne modifie ni la pente ni l'inclinaison des tangentes : en chaque point d'abscisse $x$, les courbes de $F$ et de $F + c$ ont des tangentes \emph{parallèles}. Comme la dérivée mesure précisément la pente de la tangente, on devine que $F$ et $F + c$ ont la même dérivée. Formalisons cette intuition.

\begin{propriete}[Translatées d'une primitive]
Soit $f$ une fonction continue sur un intervalle $I$ et $F$ une primitive de $f$ sur $I$. Alors, pour tout réel $c$, la fonction $F + c$ définie par $(F+c)(x) = F(x) + c$ est également une primitive de $f$ sur $I$.
\end{propriete}

\begin{experience}[Démonstration guidée]
Montrons que $F + c$ est une primitive de $f$.
\begin{enumerate}
\item $F$ est une primitive de $f$ sur $I$, donc $F$ est dérivable sur $I$ et $F' = f$.
\item La fonction constante $x \mapsto c$ est dérivable sur $I$, de dérivée nulle.
\item La somme de deux fonctions dérivables est dérivable, donc $F + c$ est dérivable sur $I$ et :
\[ (F + c)'(x) = F'(x) + 0 = f(x) \quad \text{pour tout } x \in I. \]
\item Conclusion : $F + c$ est bien une primitive de $f$ sur $I$. \hfill $\square$
\end{enumerate}
\end{experience}

Ainsi, dès qu'on connaît \emph{une} primitive, on en connaît une infinité : toutes ses translatées verticales. Reste la question réciproque, bien plus profonde : obtient-on ainsi \textbf{toutes} les primitives ? Autrement dit, deux primitives de la même fonction diffèrent-elles \emph{toujours} d'une constante ?

\courssub{Le théorème fondamental : deux primitives diffèrent d'une constante}

\begin{propriete}[Théorème — deux primitives d'une même fonction]
Soit $f$ une fonction continue sur un \textbf{intervalle} $I$. Si $F$ et $G$ sont deux primitives de $f$ sur $I$, alors il existe un réel $c$ tel que :
\[ \text{pour tout } x \in I, \quad G(x) = F(x) + c. \]
Autrement dit, deux primitives d'une même fonction sur un intervalle diffèrent d'une constante. Cette démonstration est exigible au Baccalauréat.
\end{propriete}

\begin{experience}[Démonstration guidée du théorème]
L'idée est lumineuse : pour montrer que $G$ et $F$ diffèrent d'une constante, on étudie leur \emph{différence} $G - F$ et on montre qu'elle est constante. Or comment prouver qu'une fonction est constante ? En montrant que sa dérivée est nulle !
\begin{enumerate}
\item \textbf{Hypothèses.} $F$ et $G$ sont des primitives de $f$ sur $I$, donc $F' = f$ et $G' = f$ sur $I$.
\item \textbf{On dérive la différence.} La fonction $G - F$ est dérivable sur $I$ (différence de deux fonctions dérivables) et, pour tout $x \in I$ :
\[ (G - F)'(x) = G'(x) - F'(x) = f(x) - f(x) = 0. \]
\item \textbf{On applique le théorème-clé.} Nous savons (chapitre sur la dérivation) que : \emph{si une fonction a une dérivée nulle sur un \textbf{intervalle}, alors elle est constante sur cet intervalle}. Ici $(G-F)' = 0$ sur l'intervalle $I$, donc $G - F$ est constante sur $I$.
\item \textbf{Conclusion.} Il existe un réel $c$ tel que $G(x) - F(x) = c$ pour tout $x \in I$, c'est-à-dire $G(x) = F(x) + c$. \hfill $\square$
\end{enumerate}
\end{experience}

\begin{aretenir}[Ensemble des primitives d'une fonction]
Soit $f$ une fonction continue sur un intervalle $I$ et $F_0$ une primitive \emph{particulière} de $f$ sur $I$. Alors l'ensemble de \textbf{toutes} les primitives de $f$ sur $I$ est :
\[ \left\lbrace F_0 + c \; ; \; c \in \mathbb{R} \right\rbrace. \]
On dit que les primitives de $f$ sont définies \og à une constante additive près \fg{}. Géométriquement, les courbes représentatives des primitives de $f$ forment une famille de courbes toutes déduites de l'une d'elles par des translations verticales.
\end{aretenir}

\begin{attention}[L'hypothèse \og intervalle \fg{} est cruciale]
Le théorème repose entièrement sur le fait que $I$ est un \textbf{intervalle} : c'est ce qui permet d'affirmer qu'une fonction de dérivée nulle est constante. Sur une \emph{réunion d'intervalles disjoints}, le résultat devient faux !

Contre-exemple classique : soit $f(x) = \dfrac{1}{x}$ sur $\mathbb{R}^* = \left]-\infty ; 0\right[ \cup \left]0 ; +\infty\right[$ (qui n'est \textbf{pas} un intervalle). Définissons :
\[ F(x) = \ln|x| \qquad \text{et} \qquad G(x) = \begin{cases} \ln|x| + 1 & \text{si } x > 0 \\ \ln|x| - 3 & \text{si } x < 0 \end{cases} \]
On vérifie que $F' = G' = f$ sur $\mathbb{R}^*$, pourtant $G - F$ vaut $1$ sur $\left]0 ; +\infty\right[$ et $-3$ sur $\left]-\infty ; 0\right[$ : la différence n'est \emph{pas} constante ! Sur une réunion d'intervalles, chaque morceau peut avoir sa propre constante. Retiens bien : \textbf{toujours vérifier que l'on travaille sur un intervalle} avant d'écrire \og les primitives sont $F + c$ \fg{}.
\end{attention}

\courssub{Illustration graphique : une famille de courbes parallèles}

Visualisons ce résultat avec la fonction $f(x) = 3x^2 - 4x + 1$. Une primitive est $F(x) = x^3 - 2x^2 + x$ (vérifie en dérivant !). Toutes les primitives s'écrivent $F_c(x) = x^3 - 2x^2 + x + c$ avec $c \in \mathbb{R}$. La figure ci-dessous représente trois de ces courbes, pour $c \in \{-2 \, ; \, 0 \, ; \, 2\}$.

\begin{popfigure}
\begin{center}
\begin{tikzpicture}
\begin{axis}[
    width=0.85\linewidth,
    height=8cm,
    axis lines=middle,
    xlabel={$x$}, ylabel={$y$},
    xmin=-1.2, xmax=2.6,
    ymin=-4.5, ymax=4.5,
    grid=both,
    grid style={popline!40},
    legend style={at={(0.02,0.98)}, anchor=north west, font=\small},
    samples=200
]
\addplot[popcurve, popblue, domain=-1:2.4] {x^3 - 2*x^2 + x + 2};
\addlegendentry{$c = 2$}
\addplot[popcurve, poporange, domain=-1:2.4] {x^3 - 2*x^2 + x};
\addlegendentry{$c = 0$}
\addplot[popcurve, popgreen, domain=-1:2.4] {x^3 - 2*x^2 + x - 2};
\addlegendentry{$c = -2$}
\addplot[only marks, mark=*, mark size=2.5pt, poppink] coordinates {(1,0)};
\node[popdark, anchor=south west, font=\small] at (axis cs:1.02,0.05) {$(1\,;\,0)$};
\end{axis}
\end{tikzpicture}
\end{center}
\legende{Trois primitives de $f(x) = 3x^2 - 4x + 1$ : les courbes $y = x^3 - 2x^2 + x + c$ se déduisent l'une de l'autre par translation verticale. En chaque abscisse, leurs tangentes sont parallèles (même dérivée $f$). Une seule courbe de la famille passe par le point $(1\,;\,0)$ : celle correspondant à $c = 0$.}
\end{popfigure}

Observe bien : les trois courbes ont exactement la même \emph{forme} ; elles sont simplement décalées verticalement. En tout point de même abscisse, les tangentes aux trois courbes sont parallèles, car elles ont toutes pour coefficient directeur $f(x)$. C'est la traduction géométrique du théorème.

\courssub{La condition initiale : une et une seule primitive}

Regarde à nouveau la figure : parmi toutes les courbes de la famille, une \emph{seule} passe par le point $(1\,;\,0)$. C'est un phénomène général et capital : imposer le passage par un point donné revient à fixer la constante $c$, et donc à sélectionner \textbf{une unique} primitive. C'est ce qu'on appelle une \cle{condition initiale}.

\begin{propriete}[Théorème — existence et unicité sous condition initiale]
Soit $f$ une fonction continue sur un intervalle $I$, $x_0$ un point de $I$ et $y_0$ un réel. Alors il existe une \textbf{unique} primitive $F$ de $f$ sur $I$ telle que :
\[ F(x_0) = y_0. \]
\end{propriete}

\begin{experience}[Démonstration guidée]
\textbf{Existence.} $f$ est continue sur $I$, donc elle admet une primitive $F_0$ sur $I$ (théorème d'existence admis en section 1). Posons $c = y_0 - F_0(x_0)$ et $F = F_0 + c$. Alors $F$ est une primitive de $f$ (translatée d'une primitive) et :
\[ F(x_0) = F_0(x_0) + c = F_0(x_0) + y_0 - F_0(x_0) = y_0. \]
\textbf{Unicité.} Supposons que $F$ et $G$ soient deux primitives de $f$ vérifiant toutes deux $F(x_0) = y_0$ et $G(x_0) = y_0$. D'après le théorème fondamental, il existe $k \in \mathbb{R}$ tel que $G = F + k$ sur $I$. En évaluant en $x_0$ : $G(x_0) = F(x_0) + k$, soit $y_0 = y_0 + k$, donc $k = 0$ et $G = F$. La primitive est unique. \hfill $\square$
\end{experience}

\begin{methode}[Déterminer la primitive vérifiant $F(x_0) = y_0$]
Pour trouver l'unique primitive $F$ de $f$ telle que $F(x_0) = y_0$ :
\begin{enumerate}
\item \textbf{Trouver une primitive particulière} $F_0$ de $f$ (en lisant les dérivées usuelles \og à l'envers \fg{}).
\item \textbf{Écrire la forme générale} : $F(x) = F_0(x) + c$, où $c$ est un réel à déterminer.
\item \textbf{Exploiter la condition initiale} : résoudre l'équation $F_0(x_0) + c = y_0$ d'inconnue $c$.
\item \textbf{Conclure} en écrivant explicitement $F$, puis \emph{vérifier} mentalement que $F' = f$ et $F(x_0) = y_0$.
\end{enumerate}
\end{methode}

\begin{exoresolu}[Primitive de $f(x) = 3x^2 - 4x + 1$ qui s'annule en $1$]
\textbf{Énoncé.} Déterminer la primitive $F$ de la fonction $f$ définie sur $\mathbb{R}$ par $f(x) = 3x^2 - 4x + 1$ qui s'annule en $x = 1$, c'est-à-dire telle que $F(1) = 0$.
\tcblower
\textbf{Solution détaillée.}

\emph{Étape 1 : recherche d'une primitive particulière.} On primitivise terme à terme, en lisant les formules de dérivation à l'envers :
\begin{itemize}
\item une primitive de $3x^2$ est $x^3$ (car $(x^3)' = 3x^2$) ;
\item une primitive de $-4x$ est $-2x^2$ (car $(-2x^2)' = -4x$) ;
\item une primitive de $1$ est $x$.
\end{itemize}
Donc $F_0(x) = x^3 - 2x^2 + x$ est une primitive de $f$ sur $\mathbb{R}$.

\emph{Étape 2 : forme générale.} Comme $\mathbb{R}$ est un intervalle, toutes les primitives de $f$ s'écrivent :
\[ F(x) = x^3 - 2x^2 + x + c, \quad c \in \mathbb{R}. \]

\emph{Étape 3 : condition initiale.} On impose $F(1) = 0$ :
\begin{align*}
F(1) &= 1^3 - 2 \times 1^2 + 1 + c \\
&= 1 - 2 + 1 + c \\
&= c.
\end{align*}
La condition $F(1) = 0$ donne donc $c = 0$.

\emph{Étape 4 : conclusion.} L'unique primitive cherchée est :
\[ \boxed{F(x) = x^3 - 2x^2 + x} \]
\emph{Vérification} : $F'(x) = 3x^2 - 4x + 1 = f(x)$ \checkmark{} et $F(1) = 1 - 2 + 1 = 0$ \checkmark. C'est exactement la courbe orange de la figure, la seule de la famille passant par $(1\,;\,0)$.
\end{exoresolu}

\begin{exemplebox}[Un deuxième exemple avec une condition initiale non nulle]
Déterminons la primitive $G$ de $f(x) = 3x^2 - 4x + 1$ telle que $G(2) = 5$.
\begin{enumerate}
\item Forme générale : $G(x) = x^3 - 2x^2 + x + c$.
\item Condition : $G(2) = 8 - 8 + 2 + c = 2 + c$.
\item On résout $2 + c = 5$, d'où $c = 3$.
\item Conclusion : $G(x) = x^3 - 2x^2 + x + 3$.
\end{enumerate}
Vérification : $G(2) = 8 - 8 + 2 + 3 = 5$ \checkmark.
\end{exemplebox}

\courssub{Les pièges à éviter absolument}

\begin{attention}[Erreurs fréquentes au Baccalauréat]
\begin{itemize}
\item \textbf{Oublier le \og $+c$ \fg{}.} Quand l'énoncé demande \emph{l'ensemble} des primitives (ou \emph{une} primitive sans condition initiale), écrire seulement $F(x) = x^3 - 2x^2 + x$ est incomplet : il faut écrire $F(x) = x^3 - 2x^2 + x + c$, $c \in \mathbb{R}$. Cette omission coûte des points.
\item \textbf{Croire que deux primitives peuvent différer autrement que par une constante.} Par exemple, affirmer que $x^3 - 2x^2 + x$ et $x^3 - 2x^2 + 2x$ sont deux primitives de $3x^2 - 4x + 1$ est faux : la dérivée de la seconde est $3x^2 - 4x + 2 \neq f(x)$. Deux vraies primitives diffèrent \emph{uniquement} d'une constante additive.
\item \textbf{Oublier l'hypothèse d'intervalle.} Sur $\mathbb{R}^*$, écrire que les primitives de $\frac{1}{x}$ sont $\ln|x| + c$ avec \emph{une seule} constante $c$ est incorrect : il faut une constante par intervalle (une sur $\left]-\infty;0\right[$, une autre sur $\left]0;+\infty\right[$).
\item \textbf{Mal utiliser la condition initiale.} La condition $F(x_0) = y_0$ s'applique à la forme générale $F_0(x) + c$, pas seulement à $F_0$. On résout $F_0(x_0) + c = y_0$ pour trouver $c = y_0 - F_0(x_0)$.
\end{itemize}
\end{attention}

\begin{savaistu}[Pourquoi \og condition initiale \fg{} ?]
Ce vocabulaire vient de la physique. Si $f(t)$ est la \emph{vitesse} d'un mobile à l'instant $t$ (par exemple un taxi entre Douala et Yaoundé), alors une primitive $F(t)$ représente sa \emph{position} au cours du temps. Connaître la vitesse à chaque instant ne suffit pas à connaître la position : il faut aussi savoir \emph{d'où} le mobile est parti, c'est-à-dire sa position à l'instant initial $t_0$. Cette donnée $F(t_0) = y_0$ est la condition initiale, et elle détermine le mouvement de façon unique. Ce lien entre primitives et équations différentielles, développé par Newton et Leibniz au XVII\textsuperscript{e} siècle, est au cœur de toute la physique moderne — et tu le retrouveras dans le chapitre sur les équations différentielles !
\end{savaistu}

\begin{exercice}[À toi de jouer]
\begin{enumerate}
\item Déterminer l'ensemble des primitives de $f(x) = 6x^2 + 2x - 5$ sur $\mathbb{R}$.
\item Déterminer la primitive $F$ de $f(x) = 6x^2 + 2x - 5$ sur $\mathbb{R}$ telle que $F(0) = 3$.
\item Déterminer la primitive $G$ de $g(x) = \dfrac{1}{x^2}$ sur $\left]0 ; +\infty\right[$ telle que $G(1) = 2$.
\end{enumerate}
\end{exercice}

\begin{corrige}[Exercice]
\begin{enumerate}
\item Une primitive de $6x^2$ est $2x^3$, une primitive de $2x$ est $x^2$, une primitive de $-5$ est $-5x$. Comme $\mathbb{R}$ est un intervalle, l'ensemble des primitives est :
\[ \left\lbrace x \mapsto 2x^3 + x^2 - 5x + c \; ; \; c \in \mathbb{R} \right\rbrace. \]
\item Forme générale : $F(x) = 2x^3 + x^2 - 5x + c$. Condition : $F(0) = c = 3$. Donc $F(x) = 2x^3 + x^2 - 5x + 3$.
\item Sur $\left]0 ; +\infty\right[$, une primitive de $\dfrac{1}{x^2}$ est $-\dfrac{1}{x}$ (car $\left(-\dfrac{1}{x}\right)' = \dfrac{1}{x^2}$). Forme générale : $G(x) = -\dfrac{1}{x} + c$. Condition : $G(1) = -1 + c = 2$, d'où $c = 3$. Donc $G(x) = -\dfrac{1}{x} + 3$.
\end{enumerate}
\end{corrige}

\begin{aretenir}[L'essentiel de la section]
\begin{itemize}
\item Si $F$ est une primitive de $f$ sur $I$, alors $F + c$ aussi, pour tout réel $c$.
\item Réciproquement, sur un \textbf{intervalle}, deux primitives de $f$ diffèrent d'une constante : l'ensemble des primitives est $\{F_0 + c \, ; \, c \in \mathbb{R}\}$.
\item La condition initiale $F(x_0) = y_0$ (avec $x_0 \in I$) sélectionne une \textbf{unique} primitive : on trouve $c$ en résolvant $F_0(x_0) + c = y_0$.
\item Attention : sans l'hypothèse \og intervalle \fg{}, tout s'effondre (contre-exemple de $\frac{1}{x}$ sur $\mathbb{R}^*$).
\end{itemize}
\end{aretenir}

Maintenant que tu sais décrire toutes les primitives d'une fonction et en choisir une grâce à une condition initiale, il reste à apprendre à les \emph{calculer} efficacement. C'est l'objet de la section suivante : les primitives des fonctions usuelles.

\coursec{Primitives des fonctions usuelles}

Dans la section précédente, nous avons établi un fait fondamental : toute fonction continue sur un intervalle $I$ admet des primitives, et si $F$ en est une, alors toutes les autres s'écrivent $F + c$ où $c$ est une constante réelle. Autrement dit, \emph{il suffit d'en trouver une seule} pour les connaître toutes. Mais comment en trouver une ? C'est toute la question à laquelle cette section répond. La bonne nouvelle, c'est que tu possèdes déjà l'outil essentiel : ton tableau des dérivées, appris en Première. Nous allons simplement apprendre à le lire \emph{à l'envers}.

\courssub{Le principe : lire le tableau des dérivées à l'envers}

\textbf{Dériver et primitiver : deux gestes inverses}\par

Chercher une primitive de $f$, c'est répondre à la question : \og quelle fonction $F$ a pour dérivée $f$ ? \fg{}. Or tu sais déjà dériver de nombreuses fonctions. Chaque ligne de ton tableau des dérivées, lue de droite à gauche, donne une ligne du tableau des primitives.

Par exemple, tu sais que $(x^3)' = 3x^2$. En lisant cette égalité à l'envers, tu obtiens : une primitive de $3x^2$ est $x^3$, donc une primitive de $x^2$ est $\dfrac{x^3}{3}$. De même, $(\sin x)' = \cos x$ se relit : une primitive de $\cos x$ est $\sin x$. C'est aussi simple que cela : \cle{primitiver}, c'est \cle{dériver à rebours}.

\begin{methode}[Construire le tableau des primitives]
\begin{enumerate}
\item Partir d'une formule de dérivation connue : $F' = f$.
\item La relire à l'envers : \og une primitive de $f$ est $F$ \fg{}.
\item Ajouter la constante $c \in \mathbb{R}$ pour obtenir \emph{toutes} les primitives : $F(x) + c$.
\item Préciser l'intervalle de validité (celui où $F$ est dérivable et où $F' = f$).
\end{enumerate}
\end{methode}

\begin{attention}[Le réflexe de vérification, obligatoire au Bac]
Primitiver est une opération délicate car il n'existe pas de méthode systématique pour toutes les fonctions. En revanche, \emph{vérifier} est toujours facile : il suffit de dériver le résultat obtenu. Prends l'habitude, pour chaque primitive calculée, de dériver mentalement ta réponse et de contrôler que tu retombes sur la fonction de départ. Ce réflexe te coûte dix secondes et te sauve des points au Baccalauréat.
\end{attention}

\courssub{Primitives des fonctions puissances}

\textbf{Primitive d'une constante}\par

Commençons par le cas le plus simple. Si $f(x) = a$ est une fonction constante, quelle fonction a pour dérivée $a$ ? On sait que $(ax)' = a$. Donc une primitive de $x \mapsto a$ est $x \mapsto ax$, et l'ensemble des primitives est :
\[ F(x) = ax + c, \quad c \in \mathbb{R}. \]
Par exemple, les primitives de la fonction constante $f(x) = 7$ sont les fonctions $F(x) = 7x + c$.

\textbf{Primitive de $x^r$ ($r$ rationnel, $r \neq -1$)}\par

Tu sais que pour tout entier $n \geq 1$, $(x^n)' = n x^{n-1}$ : dériver fait \emph{descendre} l'exposant de $1$. Primitiver doit donc faire le contraire : \emph{monter} l'exposant de $1$, puis compenser le facteur parasite en divisant par le nouvel exposant. Vérifions : si l'on pose $F(x) = \dfrac{x^{r+1}}{r+1}$, alors
\[ F'(x) = \frac{1}{r+1} \cdot (r+1)\, x^{r} = x^r. \]
Le calcul fonctionne dès que $r + 1 \neq 0$, c'est-à-dire $r \neq -1$, et il reste valable pour tout exposant rationnel $r$ (sur un intervalle où $x^r$ est défini, typiquement $]0\,;+\infty[$ pour $r \notin \mathbb{Z}$).

\begin{propriete}[Primitives usuelles : puissances et trigonométrie]
Soit $c$ une constante réelle arbitraire.
\begin{itemize}
\item Pour tout rationnel $r \neq -1$, les primitives de $f(x) = x^r$ sur $]0\,;+\infty[$ (ou $\mathbb{R}$ si $r \in \mathbb{N}$, ou $\mathbb{R}^*$ si $r \in \mathbb{Z}_{<0}$) sont
\[ F(x) = \frac{x^{r+1}}{r+1} + c. \]
\item Pour tous réels $a \neq 0$ et $b$, les primitives de $f(x) = \cos(ax+b)$ sur $\mathbb{R}$ sont
\[ F(x) = \frac{\sin(ax+b)}{a} + c. \]
\item Pour tous réels $a \neq 0$ et $b$, les primitives de $f(x) = \sin(ax+b)$ sur $\mathbb{R}$ sont
\[ F(x) = -\frac{\cos(ax+b)}{a} + c. \]
\end{itemize}
\end{propriete}

\begin{experience}[Démonstration guidée de la formule trigonométrique]
Démontrons que $F(x) = \dfrac{\sin(ax+b)}{a}$ est bien une primitive de $f(x) = \cos(ax+b)$.
\begin{enumerate}
\item Rappel : la dérivée de $\sin u$ est $u' \cos u$ (dérivation des fonctions composées).
\item Ici $u(x) = ax + b$, donc $u'(x) = a$.
\item Alors $F'(x) = \dfrac{1}{a} \cdot a \cos(ax+b) = \cos(ax+b) = f(x)$. \quad $\square$
\end{enumerate}
Le facteur $\dfrac{1}{a}$ sert exactement à \emph{compenser} le facteur $a$ qui apparaît quand on dérive la composition. Retiens cette idée : elle reviendra pour toutes les formes composées.
\end{experience}

\begin{attention}[Deux pièges classiques]
\begin{itemize}
\item \textbf{Le cas $r = -1$.} La formule $x^r \to \dfrac{x^{r+1}}{r+1}$ ne s'applique \emph{pas} pour $r = -1$ : on obtiendrait $\dfrac{x^0}{0}$, une division par zéro ! La fonction $x \mapsto \dfrac{1}{x}$ est un cas à part : ses primitives sont $F(x) = \ln|x| + c$ (sur $]0\,;+\infty[$ ou $]-\infty\,;0[$). Ce cas sera étudié en détail avec la fonction logarithme.
\item \textbf{Le facteur $\dfrac{1}{a}$ oublié.} L'erreur la plus fréquente au Bac : écrire que la primitive de $\cos(3x)$ est $\sin(3x)$. Faux ! En dérivant $\sin(3x)$ on obtient $3\cos(3x)$, pas $\cos(3x)$. La bonne réponse est $\dfrac{\sin(3x)}{3}$. Vérifie toujours en dérivant.
\end{itemize}
\end{attention}

\textbf{Trois cas particuliers à connaître par cœur}\par

Certaines puissances reviennent sans cesse dans les sujets. Écrivons-les sous forme d'exposants rationnels pour appliquer la formule :

\begin{align*}
\frac{1}{x^2} &= x^{-2} \quad\Longrightarrow\quad F(x) = \frac{x^{-1}}{-1} + c = -\frac{1}{x} + c,\\[2pt]
\sqrt{x} &= x^{1/2} \quad\Longrightarrow\quad F(x) = \frac{x^{3/2}}{3/2} + c = \frac{2}{3}\,x\sqrt{x} + c,\\[2pt]
\frac{1}{\sqrt{x}} &= x^{-1/2} \quad\Longrightarrow\quad F(x) = \frac{x^{1/2}}{1/2} + c = 2\sqrt{x} + c.
\end{align*}

\begin{aretenir}[À mémoriser]
\[ \frac{1}{x^2} \xrightarrow{\ \text{primitive}\ } -\frac{1}{x} \qquad ; \qquad \sqrt{x} \xrightarrow{\ \text{primitive}\ } \frac{2}{3}\,x\sqrt{x} \qquad ; \qquad \frac{1}{\sqrt{x}} \xrightarrow{\ \text{primitive}\ } 2\sqrt{x}. \]
Ces trois résultats se retrouvent aussi en lisant à l'envers les dérivées classiques $\left(\dfrac{1}{x}\right)' = -\dfrac{1}{x^2}$ et $(\sqrt{x})' = \dfrac{1}{2\sqrt{x}}$.
\end{aretenir}

\courssub{Le tableau complet des primitives usuelles}

Voici le tableau de référence de cette leçon. Il est construit ligne par ligne par inversion du tableau des dérivées. Apprends-le dans les deux sens : de $f$ vers $F$ pour primitiver, de $F$ vers $f$ pour vérifier.

\begin{center}
\begin{tabularx}{\linewidth}{|>{\centering\arraybackslash}X|>{\centering\arraybackslash}X|>{\centering\arraybackslash}X|}
\hline
\rowcolor{popblueL}
\textbf{Fonction $f(x)$} & \textbf{Primitive $F(x)$ ($c \in \mathbb{R}$)} & \textbf{Intervalle de validité} \\
\hline
$a$ (constante) & $ax + c$ & $\mathbb{R}$ \\
\hline
$x^r \ (r \in \mathbb{Q},\ r \neq -1)$ & $\dfrac{x^{r+1}}{r+1} + c$ & $]0\,;+\infty[$ (général), $\mathbb{R}$ si $r\in\mathbb{N}$, $\mathbb{R}^*$ si $r\in\mathbb{Z}_{<0}$ \\
\hline
$\dfrac{1}{x^2}$ & $-\dfrac{1}{x} + c$ & $]0\,;+\infty[$ ou $]-\infty\,;0[$ \\
\hline
$\sqrt{x}$ & $\dfrac{2}{3}\,x\sqrt{x} + c$ & $]0\,;+\infty[$ \\
\hline
$\dfrac{1}{\sqrt{x}}$ & $2\sqrt{x} + c$ & $]0\,;+\infty[$ \\
\hline
$\cos(ax+b)$ ($a \neq 0$) & $\dfrac{\sin(ax+b)}{a} + c$ & $\mathbb{R}$ \\
\hline
$\sin(ax+b)$ ($a \neq 0$) & $-\dfrac{\cos(ax+b)}{a} + c$ & $\mathbb{R}$ \\
\hline
\end{tabularx}
\end{center}
\begin{center}
\emph{Tableau des primitives usuelles, obtenu par lecture inversée du tableau des dérivées.}
\end{center}

\courssub{Linéarité : primitiver terme à terme}

En pratique, les fonctions qu'on te donne au Bac sont rarement des fonctions usuelles isolées : ce sont des sommes, des différences, des produits par des constantes. Heureusement, la primitivation se comporte bien vis-à-vis de ces opérations.

\begin{propriete}[Linéarité de la primitivation]
Soient $f$ et $g$ deux fonctions continues sur un intervalle $I$, admettant respectivement pour primitives $F$ et $G$, et soit $k$ un réel. Alors :
\begin{itemize}
\item $F + G$ est une primitive de $f + g$ sur $I$ ;
\item $kF$ est une primitive de $kf$ sur $I$.
\end{itemize}
\emph{Démonstration (immédiate, par dérivation).} $(F+G)' = F' + G' = f + g$, car la dérivation est linéaire. De même $(kF)' = kF' = kf$. \quad $\square$
\end{propriete}

\begin{methode}[Primitiver un polynôme]
Pour primitiver un polynôme, on primitive \cle{terme à terme} :
\begin{enumerate}
\item on augmente chaque exposant de $1$ ;
\item on divise chaque terme par son \emph{nouvel} exposant ;
\item on ajoute la constante $c$.
\end{enumerate}
Exemple modèle : $f(x) = 6x^2 - 4x + 3$ a pour primitives
\[ F(x) = 6\cdot\frac{x^3}{3} - 4\cdot\frac{x^2}{2} + 3x + c = 2x^3 - 2x^2 + 3x + c. \]
Vérification : $F'(x) = 6x^2 - 4x + 3 = f(x)$. \checkmark
\end{methode}

\begin{exemplebox}[Fonction trigonométrique avec coefficient]
Déterminons les primitives de $f(x) = 5\cos(2x - 1)$ sur $\mathbb{R}$.

On applique la formule avec $A = 5$, $a = 2$, $b = -1$ : une primitive de $\cos(2x-1)$ est $\dfrac{\sin(2x-1)}{2}$, donc par linéarité :
\[ F(x) = \frac{5}{2}\,\sin(2x - 1) + c, \quad c \in \mathbb{R}. \]
Vérification : $F'(x) = \dfrac{5}{2} \times 2\cos(2x-1) = 5\cos(2x-1) = f(x)$. \checkmark
\end{exemplebox}

\courssub{Exemples progressifs et exercice résolu}

\begin{exoresolu}[Un exemple complet mêlant puissances et trigonométrie]
\textbf{Énoncé.} Déterminer l'ensemble des primitives de la fonction $f$ définie sur $]0\,;+\infty[$ par
\[ f(x) = 2x^3 - \frac{5}{x^2} + 3\cos(2x). \]
\tcblower
\textbf{Solution détaillée.} La fonction $f$ est continue sur $]0\,;+\infty[$ comme somme de fonctions continues, donc elle admet des primitives sur cet intervalle. On primitive terme à terme grâce à la linéarité.

\textbf{Premier terme :} $2x^3$. On augmente l'exposant et on divise par le nouvel exposant :
\[ 2x^3 \xrightarrow{\ \ } 2\cdot\frac{x^4}{4} = \frac{x^4}{2}. \]

\textbf{Deuxième terme :} $-\dfrac{5}{x^2} = -5x^{-2}$. Une primitive \textbf{de} $\dfrac{1}{x^2}$ est $-\dfrac{1}{x}$, donc :
\[ -\frac{5}{x^2} \xrightarrow{\ \ } -5 \times \left(-\frac{1}{x}\right) = \frac{5}{x}. \]

\textbf{Troisième terme :} $3\cos(2x)$. Une primitive de $\cos(2x)$ est $\dfrac{\sin(2x)}{2}$, donc :
\[ 3\cos(2x) \xrightarrow{\ \ } \frac{3}{2}\sin(2x). \]

\textbf{Conclusion.} L'ensemble des primitives de $f$ sur $]0\,;+\infty[$ est l'ensemble des fonctions
\[ \boxed{\,F(x) = \frac{x^4}{2} + \frac{5}{x} + \frac{3}{2}\sin(2x) + c, \quad c \in \mathbb{R}.\,} \]

\textbf{Vérification (le réflexe !).} Dérivons :
\[ F'(x) = \frac{4x^3}{2} - \frac{5}{x^2} + \frac{3}{2}\times 2\cos(2x) = 2x^3 - \frac{5}{x^2} + 3\cos(2x) = f(x). \]
Le résultat est confirmé.
\end{exoresolu}

\begin{exoresolu}[Polynôme : le prix d'une course en moto-taxi]
\textbf{Énoncé.} Le coût marginal (en centaines de francs CFA par kilomètre) d'une coopérative de moto-taxis à Bafoussam est modélisé par $f(x) = 3x^2 - 8x + 10$ pour $x \in [0\,;6]$. Déterminer l'ensemble des primitives de $f$, puis celle qui vaut $4$ en $0$.
\tcblower
\textbf{Solution.} On primitive terme à terme :
\[ F(x) = 3\cdot\frac{x^3}{3} - 8\cdot\frac{x^2}{2} + 10x + c = x^3 - 4x^2 + 10x + c. \]
Pour la condition $F(0) = 4$ : $F(0) = c = 4$. La primitive cherchée est donc
\[ F(x) = x^3 - 4x^2 + 10x + 4. \]
Vérification : $F'(x) = 3x^2 - 8x + 10 = f(x)$ et $F(0) = 4$. \checkmark
\end{exoresolu}

\courssub{Une lecture graphique pour bien comprendre}

Le lien entre $f$ et $F$ est visible sur un graphique : puisque $F' = f$, la fonction $F$ \emph{croît} exactement là où $f$ est \emph{positive}, et décroît là où $f$ est négative. Observons cela avec $f(x) = \cos(2x)$ et sa primitive $F(x) = \dfrac{\sin(2x)}{2}$.

\begin{popfigure}
\begin{tikzpicture}
\begin{axis}[
  width=\linewidth, height=6.5cm,
  axis lines=middle,
  xlabel={$x$}, ylabel={$y$},
  xmin=-0.3, xmax=3.5, ymin=-1.4, ymax=1.4,
  xtick={0, 0.785, 1.571, 2.356, 3.142},
  xticklabels={$0$, $\frac{\pi}{4}$, $\frac{\pi}{2}$, $\frac{3\pi}{4}$, $\pi$},
  legend style={at={(0.5,-0.22)}, anchor=north, draw=none},
  grid=both, grid style={popline!40},
]
\addplot[popcurve, popink, very thick, domain=0:3.2, samples=200]{cos(deg(2*x))};
\addlegendentry{$f(x) = \cos(2x)$}
\addplot[popcurve, poporange, very thick, domain=0:3.2, samples=200]{sin(deg(2*x))/2};
\addlegendentry{$F(x) = \dfrac{\sin(2x)}{2}$}
\addplot[popmuted, dashed, domain=0:3.2, samples=2]{0};
\end{axis}
\end{tikzpicture}
\legende{Sur $[0\,;\frac{\pi}{4}]$, $f(x) = \cos(2x) > 0$ et $F$ est croissante ; sur $[\frac{\pi}{4}\,;\frac{3\pi}{4}]$, $f < 0$ et $F$ décroît. Le signe de $f$ pilote les variations de $F$.}
\end{popfigure}

Cette observation n'est pas anodine : elle est le cœur du théorème fondamental de l'analyse que tu rencontreras avec le calcul intégral. Retiens l'idée intuitive : \emph{la primitive $F$ accumule ce que lui donne $f$} ; quand $f$ est positive, $F$ monte, quand $f$ est négative, $F$ redescend.

\begin{savaistu}[Pourquoi le facteur $\frac{1}{a}$ est-il si naturel ?]
En physique, si la vitesse d'un mobile est $v(t) = \cos(2t)$, sa position est $x(t) = \dfrac{\sin(2t)}{2}$ : le facteur $\frac{1}{2}$ traduit le fait que plus l'oscillation est rapide (grand $a$), plus le déplacement accumulé est petit. Les ingénieurs de Camtrack ou les techniciens télécom qui étudient des signaux périodiques manipulent en permanence ce type de primitives trigonométriques, souvent sans même s'en rendre compte !
\end{savaistu}

\begin{attention}[Bilan des pièges de la section]
\begin{itemize}
\item Ne jamais appliquer la formule des puissances à $r = -1$ : la primitive de $\dfrac{1}{x}$ est $\ln|x|$, pas $\dfrac{x^0}{0}$.
\item Ne jamais oublier le facteur $\dfrac{1}{a}$ devant la primitive de $\cos(ax+b)$ ou $\sin(ax+b)$.
\item Ne pas oublier le signe moins : la primitive de $\sin(ax+b)$ est $-\dfrac{\cos(ax+b)}{a}$, car $(\cos)' = -\sin$.
\item Toujours préciser l'intervalle et la constante $c \in \mathbb{R}$ quand on demande \emph{l'ensemble} des primitives.
\item Toujours vérifier le résultat en dérivant.
\end{itemize}
\end{attention}

Tu maîtrises maintenant les primitives des fonctions usuelles. Mais au Bac, les fonctions données sont souvent \emph{habillées} : elles cachent des formes $u' \cdot u^r$ ou $\dfrac{a}{(cx+d)^r}$. La section suivante t'apprendra à les reconnaître d'un coup d'œil.

\coursec{Primitives des formes composées : $u'\cdot u^{r}$ et $\dfrac{a}{(cx+d)^{r}}$}

Dans la section précédente, tu as appris à primitiver les fonctions usuelles : puissances, cosinus, sinus, fonctions de la forme $\dfrac{a}{(cx+d)^r}$ lues directement dans un tableau. Mais au Baccalauréat, les fonctions proposées sont rarement « toutes prêtes ». Le plus souvent, elles se présentent comme des \textbf{produits} ou des \textbf{quotients} dans lesquels se cache une structure bien précise : la forme $\cle{u'\cdot u^{r}}$. Toute cette section est consacrée à l'art de \emph{reconnaître} cette forme et à \emph{ajuster} les coefficients. C'est, sans exagération, la compétence de calcul de primitives la plus rentable de l'année de Terminale.

\courssub{Le point de départ : relire la dérivation des composées à l'envers}

Rappelons un résultat que tu connais bien depuis la Première : si $u$ est une fonction dérivable sur un intervalle $I$ et $r$ un nombre rationnel, alors (sous réserve que $u^r$ soit bien définie, par exemple $u>0$ si $r$ n'est pas entier) :

\[ \left(u^{r}\right)' = r\,u'\,u^{r-1}. \]

L'idée fondatrice de cette section est toute simple : \textbf{primitiver, c'est dériver à l'envers}. La formule ci-dessus, lue de droite à gauche, affirme que $u^{r}$ est une primitive de $r\,u'\,u^{r-1}$. En changeant de notation (remplaçons $r$ par $r+1$), on obtient le théorème central de cette section.

\begin{propriete}[Primitive de $u'\cdot u^{r}$]
Soit $u$ une fonction dérivable sur un intervalle $I$, et $r$ un nombre rationnel, $r \neq -1$. On suppose de plus que $u > 0$ sur $I$ lorsque $r$ n'est pas un entier (condition qui garantit que $u^r$ et $u^{r+1}$ sont bien définis). Alors la fonction
\[ F = \frac{u^{r+1}}{r+1} \]
est une primitive sur $I$ de la fonction $f = u'\cdot u^{r}$. L'ensemble des primitives de $f$ est donc :
\[ x \longmapsto \frac{u(x)^{r+1}}{r+1} + c, \qquad c \in \mathbb{R}. \]
\end{propriete}

\begin{experience}[Démonstration guidée (exigible au Bac)]
Il suffit de vérifier que la dérivée de $F$ redonne $f$. Décomposons le calcul pas à pas.
\begin{enumerate}
\item Posons $F = \dfrac{u^{r+1}}{r+1}$. Comme $r+1 \neq 0$ (car $r \neq -1$), $F$ est le produit de la constante $\dfrac{1}{r+1}$ par la fonction composée $u^{r+1}$.
\item La dérivée de $u^{r+1}$ est, par la formule de dérivation des composées : $(r+1)\,u'\,u^{r}$.
\item Donc $F' = \dfrac{1}{r+1}\times (r+1)\,u'\,u^{r} = u'\,u^{r} = f$. \quad $\square$
\end{enumerate}
Retiens la logique : \emph{le facteur $r+1$ apparaît par dérivation, puis se simplifie}. C'est exactement pour cela qu'on divise par $r+1$ dans la formule.
\end{experience}

\begin{attention}[Le cas $r = -1$ est exclu]
La formule ne s'applique pas pour $r = -1$, car on diviserait par $r+1 = 0$. La forme $u'/u$ relève d'une autre famille (le logarithme), que nous étudierons dans la section 6 et dans la leçon sur la fonction $\ln$. Ne mélange pas les deux situations !
\end{attention}

\courssub{La technique de reconnaissance : repérer $u$, vérifier $u'$, ajuster}

Face à un produit ou à un quotient, comment savoir si la forme $u'\,u^{r}$ s'applique ? Voici la démarche systématique.

\begin{methode}[Reconnaître et traiter la forme $u'\cdot u^{r}$]
\begin{enumerate}
\item \textbf{Repérer $u$} : chercher dans l'expression une « boîte » élevée à une puissance (sous un exposant, sous une racine, ou au dénominateur). Cette boîte est le candidat $u$.
\item \textbf{Calculer $u'$} et comparer avec le facteur qui accompagne la boîte : ce facteur doit être égal à $u'$ \emph{à une constante multiplicative près}.
\item \textbf{Ajuster le coefficient} : si le facteur présent vaut $k\,u'$, écrire $f = k \times u'\,u^{r}$ et conclure :
\[ F = k\,\frac{u^{r+1}}{r+1} + c. \]
\item \textbf{Vérifier} (réflexe de champion) : dériver mentalement le résultat pour contrôler qu'on retombe sur $f$.
\end{enumerate}
\end{methode}

\begin{exemplebox}[Exemple chiffré 1 : un produit]
Déterminons les primitives de $f(x) = (2x+1)(x^{2}+x)^{5}$ sur $\mathbb{R}$.
\begin{itemize}
\item \textbf{Boîte} : $u = x^{2}+x$, élevée à la puissance $r = 5$.
\item \textbf{Dérivée} : $u' = 2x+1$. C'est \emph{exactement} l'autre facteur !
\item Donc $f = u'\,u^{5}$, et une primitive est $\dfrac{u^{6}}{6}$ :
\[ F(x) = \frac{(x^{2}+x)^{6}}{6} + c, \qquad c \in \mathbb{R}. \]
\end{itemize}
Vérification : $F'(x) = \dfrac{6(2x+1)(x^2+x)^5}{6} = f(x)$. \checkmark
\end{exemplebox}

\begin{exemplebox}[Exemple chiffré 2 : ajustement du coefficient]
Déterminons les primitives de $f(x) = \dfrac{x}{(x^{2}+1)^{2}}$ sur $\mathbb{R}$.
\begin{itemize}
\item \textbf{Boîte} : $u = x^{2}+1$ au dénominateur, donc $f = x\,u^{-2}$ avec $r = -2$.
\item \textbf{Dérivée} : $u' = 2x$. Le facteur présent est $x = \dfrac{1}{2}\,u'$.
\item On écrit donc $f = \dfrac{1}{2}\,u'\,u^{-2}$, d'où :
\[ F(x) = \frac{1}{2}\times\frac{u^{-1}}{-1} + c = -\frac{1}{2(x^{2}+1)} + c, \qquad c \in \mathbb{R}. \]
\end{itemize}
Vérification : $F'(x) = -\dfrac{1}{2}\times\dfrac{-2x}{(x^2+1)^2} = \dfrac{x}{(x^2+1)^2}$. \checkmark
\end{exemplebox}

\begin{exemplebox}[Exemple chiffré 3 : avec une racine carrée]
Déterminons les primitives de $f(x) = (3x^{2}+1)\sqrt{x^{3}+x}$ sur $]0\,;\,+\infty[$ (intervalle où $x^3+x>0$).
\begin{itemize}
\item \textbf{Boîte} : $u = x^{3}+x$, avec $\sqrt{u} = u^{1/2}$, donc $r = \dfrac{1}{2}$.
\item \textbf{Dérivée} : $u' = 3x^{2}+1$ : c'est exactement l'autre facteur.
\item Donc $f = u'\,u^{1/2}$ et :
\[ F(x) = \frac{u^{3/2}}{3/2} + c = \frac{2}{3}\left(x^{3}+x\right)^{3/2} + c = \frac{2}{3}\left(x^{3}+x\right)\sqrt{x^{3}+x} + c. \]
\end{itemize}
\end{exemplebox}

\begin{attention}[Ajustement : uniquement des constantes !]
L'astuce « multiplier et diviser par ce qui manque » ne fonctionne \textbf{que pour une constante}. Par exemple, pour $f(x) = \dfrac{1}{(x^2+1)^2}$, on aurait envie d'écrire $1 = \dfrac{1}{2x}\times 2x$ pour faire apparaître $u' = 2x$ : c'est \textbf{interdit}, car $\dfrac{1}{2x}$ n'est pas une constante et ne peut pas « sortir » de la primitive. La forme $u'\,u^r$ exige que le facteur accompagnant soit proportionnel à $u'$ par un \emph{nombre}. Si ce n'est pas le cas, la méthode ne s'applique pas.
\end{attention}

\begin{attention}[Il n'existe pas de « primitive d'un produit »]
Contrairement à la dérivation, qui possède des formules pour $(uv)'$ et $\left(\frac{u}{v}\right)'$, il n'existe \textbf{aucune formule générale} donnant une primitive d'un produit $f\times g$ ou d'un quotient $\frac{f}{g}$. La seule chose qui fonctionne ici, c'est la \emph{reconnaissance d'une forme} : $u'\,u^r$, $u'/u$, $u'\mathrm{e}^{u}$... Écris-le dans ta tête : \emph{« je ne primitive pas un produit, je reconnais une forme »}.
\end{attention}

\begin{popfigure}
\begin{center}
\begin{tikzpicture}[
  node distance=0.55cm and 0.9cm,
  decision/.style={diamond, aspect=2.6, draw=poporange, fill=poporangeL, text width=4.6cm, align=center, inner sep=1pt, font=\small},
  action/.style={rectangle, rounded corners, draw=popblue, fill=popblueL, text width=4.4cm, align=center, font=\small},
  stopb/.style={rectangle, rounded corners, draw=popink, fill=poppinkL, text width=4.2cm, align=center, font=\small},
  fleche/.style={-{Stealth[length=2.5mm]}, thick, popdark}
]
\node[decision] (q1) {La fonction contient-elle une « boîte » $u$ élevée à une puissance $r$ ?};
\node[action, below=of q1] (a1) {Poser $u$ et calculer $u'$};
\node[decision, below=of a1] (q2) {Le facteur restant est-il $k\,u'$ ($k$ constante) ?};
\node[action, below=of q2] (a2) {Écrire $f = k\,u'\,u^{r}$};
\node[action, below=of a2, fill=popgreenL, draw=popgreen] (a3) {$F = k\,\dfrac{u^{r+1}}{r+1} + c$ \quad ($r \neq -1$)};
\node[stopb, right=1.6cm of q2] (stop) {Autre méthode (forme $u'/u$, $u'\mathrm{e}^{u}$, linéarisation...)};
\draw[fleche] (q1) -- node[right, font=\small]{oui} (a1);
\draw[fleche] (a1) -- (q2);
\draw[fleche] (q2) -- node[right, font=\small]{oui} (a2);
\draw[fleche] (a2) -- (a3);
\draw[fleche] (q2) -- node[above, font=\small]{non} (stop);
\end{tikzpicture}
\end{center}
\legende{Arbre de décision pour la reconnaissance de la forme $u'\cdot u^{r}$.}
\end{popfigure}

\courssub{Application : primitives de $\dfrac{a}{(cx+d)^{r}}$}

Un cas particulier très fréquent au Bac : les fonctions du type $\dfrac{a}{(cx+d)^{r}}$ avec $r$ rationnel, $r \neq 1$. Elles rentrent directement dans le moule précédent en posant $u = cx+d$ : alors $u' = c$, une constante, donc l'ajustement est toujours possible.

\begin{propriete}[Primitive de $\dfrac{a}{(cx+d)^{r}}$]
Soient $a$, $c$, $d$ des réels avec $c \neq 0$, et $r \in \mathbb{Q}$, $r \neq 1$. Sur tout intervalle où $cx+d > 0$ (ou $cx+d \neq 0$ si $r$ est entier), une primitive de
\[ f(x) = \frac{a}{(cx+d)^{r}} \qquad \text{est} \qquad F(x) = \frac{a}{c(1-r)\,(cx+d)^{r-1}}. \]
\end{propriete}

En effet, $f = \dfrac{a}{c}\times c \times (cx+d)^{-r} = \dfrac{a}{c}\,u'\,u^{-r}$, donc
\[ F = \frac{a}{c}\times\frac{u^{-r+1}}{-r+1} = \frac{a}{c(1-r)}\,u^{1-r} = \frac{a}{c(1-r)(cx+d)^{r-1}}. \]

\begin{exemplebox}[Exemple chiffré 4]
Déterminons les primitives de $g(x) = \dfrac{5}{(2x-3)^{3}}$ sur $\left]\dfrac{3}{2}\,;\,+\infty\right[$.
\begin{itemize}
\item $u = 2x-3$, $u' = 2$, donc $g = \dfrac{5}{2}\,u'\,u^{-3}$ (ici $r = -3$ dans la forme $u'\,u^r$).
\item On applique : $G(x) = \dfrac{5}{2}\times\dfrac{u^{-2}}{-2} + c = -\dfrac{5}{4(2x-3)^{2}} + c$, $c \in \mathbb{R}$.
\end{itemize}
Vérification : $G'(x) = -\dfrac{5}{4}\times\dfrac{-2\times 2}{(2x-3)^3} = \dfrac{5}{(2x-3)^3}$. \checkmark
\end{exemplebox}

\begin{attention}[Ne pas oublier de diviser par $c$]
L'erreur classique consiste à primitiver $(2x-3)^{-3}$ comme s'il s'agissait de $x^{-3}$, en oubliant que la dérivation de la « boîte » fait apparaître le facteur $2$. Résultat faux d'un facteur $2$ ! Le réflexe de vérification par dérivation te sauvera à chaque fois.
\end{attention}

\courssub{Lien avec la condition initiale}

Une fois la forme générale $F + c$ obtenue, une condition initiale permet de déterminer la constante $c$, exactement comme dans la section 2. La seule difficulté supplémentaire est le calcul de la forme composée.

\begin{exoresolu}[Primitive passant par un point donné]
\textbf{Énoncé.} Soit $f(x) = 2x\left(x^{2}+1\right)^{3}$ définie sur $\mathbb{R}$. Déterminer la primitive $F$ de $f$ sur $\mathbb{R}$ telle que $F(0) = \dfrac{5}{4}$.
\tcblower
\textbf{Solution détaillée.}

\textbf{Étape 1 : reconnaissance de la forme.} Posons $u = x^{2}+1$ ; alors $u' = 2x$, qui est exactement le facteur présent. Donc $f = u'\,u^{3}$, avec $r = 3 \neq -1$.

\textbf{Étape 2 : ensemble des primitives.}
\[ F(x) = \frac{u^{4}}{4} + c = \frac{\left(x^{2}+1\right)^{4}}{4} + c, \qquad c \in \mathbb{R}. \]

\textbf{Étape 3 : exploitation de la condition initiale.}
\[ F(0) = \frac{(0+1)^{4}}{4} + c = \frac{1}{4} + c. \]
La condition $F(0) = \dfrac{5}{4}$ donne $\dfrac{1}{4} + c = \dfrac{5}{4}$, soit $c = 1$.

\textbf{Conclusion.} La primitive cherchée est :
\[ F(x) = \frac{\left(x^{2}+1\right)^{4}}{4} + 1. \]
Vérification : $F'(x) = \dfrac{4\times 2x\,(x^2+1)^3}{4} = 2x(x^2+1)^3 = f(x)$ et $F(0) = \dfrac{1}{4}+1 = \dfrac{5}{4}$. \checkmark
\end{exoresolu}

\begin{popfigure}
\begin{center}
\begin{tikzpicture}
\begin{axis}[
  width=0.92\linewidth, height=6.5cm,
  axis lines=middle,
  xlabel={$x$}, ylabel={$y$},
  xmin=-1.6, xmax=1.6, ymin=-0.5, ymax=6.5,
  xtick={-1,0,1}, ytick={1,2,4,6},
  legend style={at={(0.02,0.98)}, anchor=north west, font=\small},
  grid=both, grid style={popline!40}
]
\addplot[popcurve, popblue, domain=-1.45:1.45, samples=200]{2*x*(x^2+1)^3};
\addlegendentry{$f(x)=2x(x^{2}+1)^{3}$}
\addplot[popcurve, poporange, domain=-1.35:1.35, samples=200]{(x^2+1)^4/4};
\addlegendentry{$F(x)=\dfrac{(x^{2}+1)^{4}}{4}$}
\end{axis}
\end{tikzpicture}
\end{center}
\legende{La fonction $f = u'\,u^{3}$ (en bleu) et sa primitive $F = \dfrac{u^{4}}{4}$ (en orange) : en tout point, la pente de la tangente à la courbe de $F$ vaut $f(x)$. On observe que $F$ est décroissante là où $f<0$ et croissante là où $f>0$.}
\end{popfigure}

\begin{savaistu}[Pourquoi cette forme est-elle partout ?]
La forme $u'\,u^{r}$ est le moteur caché de la \emph{substitution} (ou changement de variable), technique que tu rencontreras peut-être à l'université : poser $t = u(x)$ transforme $\int u'(x)\,u(x)^{r}\,\mathrm{d}x$ en $\int t^{r}\,\mathrm{d}t$. Les ingénieurs de CAMTEL qui modélisent l'atténuation d'un signal dans une fibre optique, ou les agronomes qui étudient la croissance d'une plantation de cacao, résolvent quotidiennement des équations dont les solutions s'obtiennent par reconnaissance de telles formes. Maîtriser $u'\,u^{r}$, c'est déjà tenir un outil professionnel.
\end{savaistu}

\begin{exercice}[À toi de jouer]
Déterminer les primitives des fonctions suivantes sur l'intervalle indiqué :
\begin{enumerate}
\item $f(x) = (4x^{3}-2)\left(x^{4}-2x\right)^{7}$ sur $\mathbb{R}$ ;
\item $g(x) = \dfrac{3x^{2}}{(x^{3}+1)^{4}}$ sur $]-1\,;\,+\infty[$ ;
\item $h(x) = \dfrac{7}{(3x+1)^{2}}$ sur $\left]-\dfrac{1}{3}\,;\,+\infty\right[$, puis la primitive $H$ vérifiant $H(0) = 0$.
\end{enumerate}
\end{exercice}

\begin{corrige}[Exercice]
\begin{enumerate}
\item $u = x^{4}-2x$, $u' = 4x^{3}-2$ : forme exacte $u'\,u^{7}$. Donc $F(x) = \dfrac{(x^{4}-2x)^{8}}{8} + c$.
\item $u = x^{3}+1$, $u' = 3x^{2}$ : forme exacte $u'\,u^{-4}$. Donc $G(x) = \dfrac{u^{-3}}{-3} + c = -\dfrac{1}{3(x^{3}+1)^{3}} + c$.
\item $u = 3x+1$, $u' = 3$, donc $h = \dfrac{7}{3}\,u'\,u^{-2}$ et $H(x) = \dfrac{7}{3}\times\dfrac{u^{-1}}{-1} + c = -\dfrac{7}{3(3x+1)} + c$. La condition $H(0) = 0$ donne $-\dfrac{7}{3} + c = 0$, soit $c = \dfrac{7}{3}$. Finalement : $H(x) = \dfrac{7}{3} - \dfrac{7}{3(3x+1)}$.
\end{enumerate}
\end{corrige}

\begin{aretenir}[L'essentiel de la section]
\begin{itemize}
\item Si $u$ est dérivable sur $I$ et $r \in \mathbb{Q}$, $r \neq -1$ (avec $u>0$ si $r$ non entier) : une primitive de $\cle{u'\,u^{r}}$ est $\cle{\dfrac{u^{r+1}}{r+1}}$.
\item En particulier, une primitive de $\dfrac{a}{(cx+d)^{r}}$ est $\dfrac{a}{c(1-r)(cx+d)^{r-1}}$ pour $r \neq 1$ : n'oublie pas de diviser par $c = u'$.
\item La reconnaissance se fait en trois temps : repérer la boîte $u$, calculer $u'$, ajuster une \textbf{constante} multiplicative.
\item Il n'existe pas de formule pour primitiver un produit ou un quotient quelconque : seule la reconnaissance d'une forme fonctionne.
\item Réflexe permanent : \textbf{vérifier le résultat en dérivant}.
\end{itemize}
\end{aretenir}

\coursec{Primitives de $\cos^n x$ et $\sin^n x$}

Dans la section précédente, tu as appris à reconnaître les formes composées $u'\cdot u^r$ et $\dfrac{a}{(cx+d)^r}$. Nous allons maintenant appliquer ces techniques à une famille de fonctions emblématique du programme de Terminale : les puissances de $\cos x$ et $\sin x$. Ces fonctions apparaissent très souvent au Baccalauréat, notamment dans les calculs d'intégrales et d'aires. Bonne nouvelle : il n'y a que \textbf{deux stratégies} à connaître, et le choix entre elles dépend d'un seul critère — la \cle{parité de l'exposant} $n$.

\courssub{Le problème posé : pourquoi $\cos^2 x$ n'est pas une forme $u'\cdot u^r$}

Cherchons une primitive de $f(x)=\cos^2 x$. Un réflexe tentant (mais faux !) serait d'appliquer la formule de $u^r$ : « la primitive de $u^2$ est $\dfrac{u^3}{3}$, donc la primitive de $\cos^2 x$ est $\dfrac{\cos^3 x}{3}$ ». Vérifions par dérivation :
\[
\left(\frac{\cos^3 x}{3}\right)' = \frac{3\cos^2 x \cdot (-\sin x)}{3} = -\sin x \cos^2 x \neq \cos^2 x.
\]
Le résultat est faux. Pourquoi ? Parce que la formule $\displaystyle\int u^r\,u' = \frac{u^{r+1}}{r+1}$ exige la présence du facteur $u'$ \emph{en facteur} dans l'expression. Ici, avec $u=\cos x$, il manque le facteur $u'=-\sin x$.

\begin{attention}[Le piège classique du Bac]
La primitive de $\cos^2 x$ n'est \textbf{pas} $\dfrac{\cos^3 x}{3}$, et celle de $\sin^2 x$ n'est \textbf{pas} $\dfrac{\sin^3 x}{3}$. La forme $u'\cdot u^r$ exige la présence du facteur $u'$ : or $\cos^2 x$ ne contient pas le facteur $-\sin x$. Il faut donc \emph{transformer} l'expression avant de primitiver. Deux cas se présentent selon la parité de $n$.
\end{attention}

\courssub{Cas d'une puissance paire : linéarisation (cas $n=2$)}

\textbf{Intuition.} Les fonctions $\cos^2 x$ et $\sin^2 x$ oscillent entre $0$ et $1$, avec une valeur moyenne égale à $\tfrac12$ (puisque $\cos^2 x + \sin^2 x = 1$). Les formules de duplication vont « casser » le carré et ramener l'expression à une somme de termes que l'on sait primitiver : une constante et un cosinus.

\begin{propriete}[Formules de duplication — à connaître par cœur]
Pour tout réel $x$ :
\[
\cos^2 x = \frac{1+\cos(2x)}{2} \qquad \text{et} \qquad \sin^2 x = \frac{1-\cos(2x)}{2}.
\]
\end{propriete}

\begin{experience}[Démonstration guidée des formules de duplication]
Rappelons la formule d'addition : $\cos(a+b)=\cos a\cos b - \sin a \sin b$.
\begin{enumerate}
\item Pose $a=b=x$ : $\cos(2x)=\cos^2 x - \sin^2 x$.
\item Utilise $\sin^2 x = 1-\cos^2 x$ : alors $\cos(2x) = \cos^2 x - (1-\cos^2 x) = 2\cos^2 x - 1$.
\item Isole $\cos^2 x$ : $\cos^2 x = \dfrac{1+\cos(2x)}{2}$. \quad $\blacksquare$
\item De même, avec $\cos^2 x = 1-\sin^2 x$ : $\cos(2x) = 1 - 2\sin^2 x$, d'où $\sin^2 x = \dfrac{1-\cos(2x)}{2}$. \quad $\blacksquare$
\end{enumerate}
\end{experience}

\begin{methode}[Primitiver une puissance paire : linéariser]
Pour déterminer une primitive de $\cos^2 x$ ou $\sin^2 x$ :
\begin{enumerate}
\item \textbf{Linéariser} : remplacer $\cos^2 x$ par $\dfrac{1+\cos(2x)}{2}$ (resp. $\sin^2 x$ par $\dfrac{1-\cos(2x)}{2}$).
\item \textbf{Primitiver terme à terme} : la constante $\tfrac12$ a pour primitive $\tfrac{x}{2}$, et $\tfrac{\cos(2x)}{2}$ a pour primitive $\tfrac{\sin(2x)}{4}$.
\item \textbf{Vérifier} en dérivant le résultat obtenu.
\end{enumerate}
\end{methode}

\begin{propriete}[Primitives de $\cos^2 x$ et $\sin^2 x$]
Les primitives de $\cos^2 x$ et $\sin^2 x$ sur $\mathbb{R}$ sont :
\[
\int \cos^2 x\,dx = \frac{x}{2} + \frac{\sin(2x)}{4} + c, \qquad
\int \sin^2 x\,dx = \frac{x}{2} - \frac{\sin(2x)}{4} + c, \quad c\in\mathbb{R}.
\]
\end{propriete}

\begin{exemplebox}[Primitive de $\cos^2 x$, avec vérification]
Linéarisons : $\cos^2 x = \dfrac{1}{2} + \dfrac{\cos(2x)}{2}$. En primitivant terme à terme :
\[
F(x) = \frac{x}{2} + \frac{\sin(2x)}{4} + c.
\]
\textbf{Vérification par dérivation :}
\[
F'(x) = \frac{1}{2} + \frac{2\cos(2x)}{4} = \frac{1}{2} + \frac{\cos(2x)}{2} = \frac{1+\cos(2x)}{2} = \cos^2 x. \quad \checkmark
\]
\end{exemplebox}

\begin{popfigure}
\begin{tikzpicture}
\begin{axis}[
  width=\linewidth, height=6.2cm,
  axis lines=middle, xlabel=$x$, ylabel=$y$,
  xmin=-0.4, xmax=6.7, ymin=-0.3, ymax=3.4,
  xtick={1.5708,3.1416,4.7124,6.2832},
  xticklabels={$\frac{\pi}{2}$,$\pi$,$\frac{3\pi}{2}$,$2\pi$},
  legend style={at={(0.98,0.55)},anchor=east,font=\small},
  samples=200, domain=0:6.2832]
\addplot[popcurve,popink,very thick] {cos(deg(x))^2};
\addlegendentry{$f(x)=\cos^2 x$}
\addplot[popcurve,poporange,very thick] {x/2 + sin(deg(2*x))/4};
\addlegendentry{$F(x)=\dfrac{x}{2}+\dfrac{\sin(2x)}{4}$}
\end{axis}
\end{tikzpicture}
\legende{La fonction $f(x)=\cos^2 x$ est toujours positive : sa primitive $F$ est donc croissante sur $\mathbb{R}$. On observe que $F$ « monte » avec de petites ondulations, et que $F$ marque un ralentissement (tangente horizontale) exactement là où $f$ s'annule.}
\end{popfigure}

\begin{savaistu}[La linéarisation, un outil de physicien]
La linéarisation n'est pas qu'une astuce de calcul : elle révèle que la \textbf{valeur moyenne} de $\cos^2 x$ sur une période vaut $\tfrac12$, puisque le terme $\tfrac{\cos(2x)}{2}$ oscille autour de $0$. En physique, la puissance moyenne dissipée par un courant alternatif $i(t)=I\cos(\omega t)$ dans une résistance est proportionnelle à la moyenne de $\cos^2(\omega t)$, donc à $\tfrac{I^2}{2}$ : c'est l'origine de la notion de \emph{valeur efficace} $I_{\text{eff}} = I/\sqrt{2}$ affichée sur les compteurs ENEO. Au Bac, tu réutiliseras la linéarisation pour calculer des intégrales comme $\displaystyle\int_0^{\pi}\cos^2 x\,dx = \frac{\pi}{2}$.
\end{savaistu}

\courssub{Cas d'une puissance impaire : la forme $u'\cdot u^2$ (cas $n=3$)}

\textbf{Intuition.} Quand l'exposant est impair, on peut « détacher » un facteur $\cos x$ (ou $\sin x$) et transformer le carré restant grâce à $\cos^2 x + \sin^2 x = 1$. Le facteur détaché joue alors le rôle de $u'$, et l'on retombe sur la forme $u'\cdot u^2$ étudiée à la section précédente.

\begin{methode}[Primitiver une puissance impaire]
Pour primitiver $\cos^3 x$ (resp. $\sin^3 x$) :
\begin{enumerate}
\item \textbf{Détacher} un facteur : $\cos^3 x = \cos x \cdot \cos^2 x$.
\item \textbf{Transformer} le carré avec la relation fondamentale : $\cos^2 x = 1-\sin^2 x$, d'où
\[
\cos^3 x = \cos x\,(1-\sin^2 x) = \cos x - \cos x \sin^2 x.
\]
\item \textbf{Identifier} la forme $u'\cdot u^2$ : avec $u=\sin x$, on a $u'=\cos x$, donc $\cos x \sin^2 x = u'\cdot u^2$, de primitive $\dfrac{u^3}{3}=\dfrac{\sin^3 x}{3}$.
\item Pour $\sin^3 x$, écrire symétriquement $\sin^3 x = \sin x\,(1-\cos^2 x)$ et poser $u=\cos x$, $u'=-\sin x$.
\end{enumerate}
\end{methode}

\begin{propriete}[Primitives de $\cos^3 x$ et $\sin^3 x$]
Sur $\mathbb{R}$ :
\[
\int \cos^3 x\,dx = \sin x - \frac{\sin^3 x}{3} + c, \qquad
\int \sin^3 x\,dx = -\cos x + \frac{\cos^3 x}{3} + c, \quad c\in\mathbb{R}.
\]
\end{propriete}

\begin{exemplebox}[Primitive de $\sin^3 x$, détaillée]
Écrivons $\sin^3 x = \sin x \cdot \sin^2 x = \sin x\,(1-\cos^2 x) = \sin x - \sin x \cos^2 x$.
\begin{itemize}
\item Une primitive de $\sin x$ est $-\cos x$.
\item Pour $\sin x \cos^2 x$ : posons $u=\cos x$, alors $u'=-\sin x$, donc $\sin x \cos^2 x = -u'\cdot u^2$, de primitive $-\dfrac{\cos^3 x}{3}$.
\end{itemize}
Par conséquent :
\[
F(x) = -\cos x - \left(-\frac{\cos^3 x}{3}\right) + c = -\cos x + \frac{\cos^3 x}{3} + c.
\]
\textbf{Vérification :} $F'(x) = \sin x + \dfrac{3\cos^2 x\,(-\sin x)}{3} = \sin x\,(1-\cos^2 x) = \sin x \sin^2 x = \sin^3 x$. \quad $\checkmark$
\end{exemplebox}

\courssub{Généralisation : la règle de parité}

\begin{aretenir}[Stratégie selon la parité de $n$]
Pour primitiver $\cos^n x$ ou $\sin^n x$ :
\begin{itemize}
\item Si $n$ est \cle{impair} : détacher un facteur $\cos x$ (ou $\sin x$), transformer la puissance paire restante avec $\cos^2+\sin^2=1$, et reconnaître la forme \cle{$u'\cdot u^r$}.
\item Si $n$ est \cle{pair} : \cle{linéariser} à l'aide des formules de duplication (éventuellement itérées : pour $n=4$, on linéarise deux fois), puis primitiver terme à terme.
\end{itemize}
\end{aretenir}

\courssub{Exemples avec coefficients et exercice résolu}

Lorsque l'argument est $ax+b$ au lieu de $x$, la stratégie est identique ; il suffit de prêter attention au coefficient $a$ lors de la primitive (division par $a$).

\begin{exoresolu}[Primitive de $\cos^2(2x)$]
Déterminer les primitives de $f(x)=\cos^2(2x)$ sur $\mathbb{R}$, puis celle qui s'annule en $x=0$.
\tcblower
\textbf{Étape 1 : linéarisation.} La formule $\cos^2 \theta = \dfrac{1+\cos(2\theta)}{2}$ avec $\theta = 2x$ donne :
\[
\cos^2(2x) = \frac{1+\cos(4x)}{2} = \frac12 + \frac{\cos(4x)}{2}.
\]
\textbf{Étape 2 : primitive terme à terme.} Une primitive de $\cos(4x)$ est $\dfrac{\sin(4x)}{4}$, donc :
\[
F(x) = \frac{x}{2} + \frac{\sin(4x)}{8} + c.
\]
\textbf{Étape 3 : condition initiale.} $F(0) = 0 + 0 + c = 0$ impose $c=0$. La primitive cherchée est :
\[
F(x) = \frac{x}{2} + \frac{\sin(4x)}{8}.
\]
\textbf{Vérification :} $F'(x) = \dfrac12 + \dfrac{4\cos(4x)}{8} = \dfrac{1+\cos(4x)}{2} = \cos^2(2x)$. \quad $\checkmark$
\end{exoresolu}

\begin{exoresolu}[Primitive de $\sin^3(3x)$]
Déterminer une primitive de $g(x)=\sin^3(3x)$ sur $\mathbb{R}$.
\tcblower
\textbf{Étape 1 : forme $u'\cdot u^2$.} Écrivons, avec $\theta = 3x$ :
\[
\sin^3(3x) = \sin(3x)\bigl(1-\cos^2(3x)\bigr) = \sin(3x) - \sin(3x)\cos^2(3x).
\]
\textbf{Étape 2 : primitives.}
\begin{itemize}
\item Une primitive de $\sin(3x)$ est $-\dfrac{\cos(3x)}{3}$.
\item Pour $\sin(3x)\cos^2(3x)$ : posons $u=\cos(3x)$, alors $u'=-3\sin(3x)$, donc
\[
\sin(3x)\cos^2(3x) = -\frac13\,u'\cdot u^2, \quad \text{de primitive } -\frac13\cdot\frac{\cos^3(3x)}{3} = -\frac{\cos^3(3x)}{9}.
\]
\end{itemize}
\textbf{Étape 3 : conclusion.}
\[
G(x) = -\frac{\cos(3x)}{3} + \frac{\cos^3(3x)}{9} + c.
\]
\textbf{Vérification :} $G'(x) = \sin(3x) - \dfrac{9\cos^2(3x)\sin(3x)}{9} = \sin(3x)\bigl(1-\cos^2(3x)\bigr) = \sin^3(3x)$. \quad $\checkmark$
\end{exoresolu}

\begin{attention}[Réflexe de sécurité : toujours vérifier par dérivation]
La vérification par dérivation prend moins d'une minute et détecte immédiatement les deux erreurs les plus fréquentes : l'oubli du coefficient $a$ dans la division (par exemple écrire $\tfrac{\sin(4x)}{4}$ au lieu de $\tfrac{\sin(4x)}{8}$ pour la primitive de $\tfrac{\cos(4x)}{2}$) et l'erreur de signe sur la primitive de $\sin$. Au Baccalauréat, cette vérification rédigée en une ligne est toujours valorisée.
\end{attention}

\begin{exercice}[Pour t'entraîner]
\begin{enumerate}
\item Déterminer les primitives de $f(x)=\sin^2 x$ sur $\mathbb{R}$, puis celle vérifiant $F\!\left(\tfrac{\pi}{2}\right)=\tfrac{\pi}{4}$.
\item Déterminer une primitive de $g(x)=\cos^3 x$ sur $\mathbb{R}$ et vérifier le résultat par dérivation.
\item Déterminer une primitive de $h(x)=\cos^2\!\left(\dfrac{x}{2}\right)$. (Indication : linéariser avec $\theta = \tfrac{x}{2}$.)
\end{enumerate}
\end{exercice}

\begin{corrige}[Exercice : Pour t'entraîner]
\textbf{1.} $\sin^2 x = \dfrac{1-\cos(2x)}{2}$, donc $F(x) = \dfrac{x}{2} - \dfrac{\sin(2x)}{4} + c$. La condition $F\!\left(\tfrac{\pi}{2}\right) = \dfrac{\pi}{4} - \dfrac{\sin \pi}{4} + c = \dfrac{\pi}{4} + c = \dfrac{\pi}{4}$ donne $c=0$. D'où $F(x)=\dfrac{x}{2}-\dfrac{\sin(2x)}{4}$.

\textbf{2.} $\cos^3 x = \cos x\,(1-\sin^2 x) = \cos x - \cos x \sin^2 x$. Avec $u=\sin x$, $\cos x \sin^2 x = u'u^2$ a pour primitive $\dfrac{\sin^3 x}{3}$. Donc $G(x) = \sin x - \dfrac{\sin^3 x}{3} + c$. 
\textbf{Vérification :} $G'(x) = \cos x - \dfrac{3\sin^2 x \cos x}{3} = \cos x (1-\sin^2 x) = \cos^3 x$. \quad $\checkmark$

\textbf{3.} $\cos^2\!\left(\dfrac{x}{2}\right) = \dfrac{1+\cos x}{2}$, donc $H(x) = \dfrac{x}{2} + \dfrac{\sin x}{2} + c$. 
\textbf{Vérification :} $H'(x)=\dfrac{1+\cos x}{2}=\cos^2\!\left(\dfrac{x}{2}\right)$. \quad $\checkmark$
\end{corrige}

\begin{aretenir}[L'essentiel de la section]
\begin{itemize}
\item Puissance \textbf{paire} : on \textbf{linéarise} : $\cos^2 x = \dfrac{1+\cos 2x}{2}$, $\sin^2 x = \dfrac{1-\cos 2x}{2}$.
\item Puissance \textbf{impaire} : on \textbf{détache un facteur} pour créer la forme $u'\cdot u^2$ via $\cos^2+\sin^2=1$.
\item $\displaystyle\int\cos^2 x\,dx = \frac{x}{2}+\frac{\sin 2x}{4}+c$ ; \quad $\displaystyle\int\sin^3 x\,dx = -\cos x + \frac{\cos^3 x}{3}+c$.
\item Avec un argument $ax+b$, la méthode est inchangée, mais on divise par $a$ en primitivant.
\item \textbf{Toujours vérifier par dérivation.}
\end{itemize}
\end{aretenir}

\coursec{Compléments utiles au Bac : u′/u et u′eᵘ ; synthèse méthodologique}

Dans la section précédente, tu as appris à primitiver $\cos^n x$ et $\sin^n x$ grâce à la linéarisation : une technique qui consiste à \emph{transformer} l'écriture de la fonction pour la ramener à des formes connues. Cette idée est en réalité la clé de tout le chapitre : \textbf{primitiver, c'est reconnaître}. Reconnaître une fonction du tableau usuel, reconnaître une somme, reconnaître une forme $u'\,u^r$... Dans cette dernière section, nous ajoutons deux formes composées extrêmement fréquentes au Baccalauréat — $\dfrac{u'}{u}$ et $u'\,e^{u}$ — puis nous organisons tout ce que tu sais en une véritable \cle{stratégie de recherche de primitive}, celle qui te fera gagner un temps précieux le jour de l'épreuve.

\courssub{Deux compléments indispensables : les formes $\dfrac{u'}{u}$ et $u'\,e^{u}$}

\textbf{Intuition.} Tu sais dériver $\ln u$ et $e^u$ : les formules de dérivation des fonctions composées donnent $(\ln u)' = \dfrac{u'}{u}$ et $(e^u)' = u'\,e^{u}$. Or primitiver, c'est « remonter » la dérivation. En lisant ces deux formules \emph{à l'envers}, on obtient immédiatement deux nouvelles primitives composées. Le programme officiel ne les liste pas explicitement, mais elles apparaissent dans la quasi-totalité des sujets de Bac C et D : les connaître est un avantage décisif.

\begin{propriete}[Compléments utiles au Bac]
Soit $u$ une fonction dérivable sur un intervalle $I$.
\begin{enumerate}
\item Si $u$ \textbf{ne s'annule pas} sur $I$, alors la fonction $x \mapsto \dfrac{u'(x)}{u(x)}$ admet pour primitives sur $I$ les fonctions :
\[ x \mapsto \ln|u(x)| + c, \qquad c \in \mathbb{R}. \]
\item La fonction $x \mapsto u'(x)\,e^{u(x)}$ admet pour primitives sur $I$ les fonctions :
\[ x \mapsto e^{u(x)} + c, \qquad c \in \mathbb{R}. \]
\end{enumerate}
\end{propriete}

\begin{experience}[Démonstration guidée]
La démonstration est formatrice et peut être demandée sous la forme « vérifier que $F$ est une primitive de $f$ ».
\begin{enumerate}
\item \textbf{Forme $\dfrac{u'}{u}$.} Posons $F(x) = \ln|u(x)|$. Comme $u$ ne s'annule pas sur $I$ et y est continue (car dérivable), $u$ garde un signe constant sur $I$ (théorème des valeurs intermédiaires).
\begin{itemize}
\item Si $u > 0$ sur $I$ : $F(x) = \ln(u(x))$, donc $F'(x) = \dfrac{u'(x)}{u(x)}$ par dérivation de $\ln \circ\, u$.
\item Si $u < 0$ sur $I$ : $F(x) = \ln(-u(x))$, donc $F'(x) = \dfrac{-u'(x)}{-u(x)} = \dfrac{u'(x)}{u(x)}$.
\end{itemize}
Dans les deux cas, $F'(x) = \dfrac{u'(x)}{u(x)}$ : $F$ est bien une primitive de $\dfrac{u'}{u}$.
\item \textbf{Forme $u'\,e^{u}$.} Posons $G(x) = e^{u(x)}$. Par la formule de dérivation des composées, $G'(x) = u'(x)\,e^{u(x)}$. Donc $G$ est bien une primitive de $u'\,e^{u}$. \hfill$\blacksquare$
\end{enumerate}
\end{experience}

\begin{attention}[Le piège classique du dénominateur]
La formule $\displaystyle \int \frac{u'}{u} \rightsquigarrow \ln|u|$ exige que le \textbf{numérateur soit exactement la dérivée du dénominateur} (à une constante multiplicative près que l'on ajuste). Erreur fréquente : écrire qu'une primitive de $\dfrac{1}{x^2+1}$ serait $\ln(x^2+1)$. C'est \textbf{faux} : la dérivée de $x^2+1$ est $2x$, et le numérateur vaut $1$, pas $2x$. La fonction $\dfrac{1}{x^2+1}$ ne relève pas de la forme $\dfrac{u'}{u}$ (sa primitive est la fonction $\arctan$, que tu rencontreras plus tard). Avant d'écrire $\ln|u|$, vérifie toujours : « le numérateur est-il $u'$, à une constante près ? »
\end{attention}

\begin{exemplebox}[Primitive de la fonction tangente]
Déterminons les primitives de $f(x) = \tan x$ sur l'intervalle $I = \left]-\dfrac{\pi}{2}\,;\,\dfrac{\pi}{2}\right[$.
Sur $I$, $\cos x > 0$, donc $\cos$ ne s'annule pas. Écrivons :
\[ \tan x = \frac{\sin x}{\cos x} = -\,\frac{-\sin x}{\cos x} = -\,\frac{u'(x)}{u(x)} \quad \text{avec } u(x) = \cos x. \]
On reconnaît la forme $\dfrac{u'}{u}$ au facteur $-1$ près. Donc :
\[ F(x) = -\ln(\cos x) + c, \qquad c \in \mathbb{R}. \]
(On peut écrire $\ln(\cos x)$ sans valeur absolue puisque $\cos x > 0$ sur $I$.)
\end{exemplebox}

\begin{exemplebox}[Forme $u'\,e^{u}$]
Déterminons les primitives de $g(x) = x\,e^{x^2}$ sur $\mathbb{R}$.
Posons $u(x) = x^2$ : alors $u'(x) = 2x$. Or le facteur devant l'exponentielle est $x = \dfrac{1}{2}\times 2x = \dfrac{1}{2}u'(x)$. On ajuste :
\[ g(x) = \frac{1}{2}\times 2x\,e^{x^2} = \frac{1}{2}\,u'(x)\,e^{u(x)}. \]
Par conséquent :
\[ G(x) = \frac{1}{2}\,e^{x^2} + c, \qquad c \in \mathbb{R}. \]
Vérification : $G'(x) = \dfrac{1}{2}\times 2x\,e^{x^2} = x\,e^{x^2} = g(x)$. \checkmark
\end{exemplebox}

\courssub{Stratégie générale face à une fonction à primitiver}

Face à une fonction $f$ dont on cherche une primitive, le réflexe du bon candidat n'est pas de foncer tête baissée, mais de se poser, \textbf{dans l'ordre}, quatre questions.

\begin{methode}[Les quatre questions de la recherche de primitive]
\begin{enumerate}
\item \textbf{Est-ce une fonction usuelle du tableau ?} (puissances $x^r$, $\dfrac{1}{x}$, $\cos$, $\sin$, $e^x$, formes $a\cos(ax+b)$, $a\sin(ax+b)$, $\dfrac{a}{(cx+d)^r}$...) Si oui, on applique directement la formule.
\item \textbf{Est-ce une somme ou une combinaison linéaire ?} Si oui, on primitive \textbf{terme à terme} (une primitive d'une somme est la somme des primitives).
\item \textbf{Reconnaît-on une forme composée ?} On cherche à faire apparaître, à une constante multiplicative près :
\begin{itemize}
\item $u'\,u^{r}$ \quad $\rightsquigarrow$ \quad $\dfrac{u^{r+1}}{r+1}$ (pour $r \neq -1$) ;
\item $\dfrac{u'}{u}$ \quad $\rightsquigarrow$ \quad $\ln|u|$ (si $u$ ne s'annule pas) ;
\item $u'\,e^{u}$ \quad $\rightsquigarrow$ \quad $e^{u}$.
\end{itemize}
On \textbf{ajuste} la constante, puis on \textbf{vérifie en dérivant}.
\item \textbf{Faut-il transformer l'écriture ?} Pour les puissances de $\cos$ et $\sin$, on \textbf{linéarise} avec les formules d'Euler ou les formules de duplication ; parfois une identité remarquable ou une décomposition suffit.
\end{enumerate}
Enfin, dans tous les cas : on conclut par $+c$ sur l'intervalle considéré, et si une condition initiale est donnée, on détermine $c$.
\end{methode}

\begin{popfigure}
\begin{center}
\begin{tikzpicture}[
  font=\small,
  question/.style={draw=popblue, fill=popblueL, rounded corners=2pt, text width=4.6cm, align=center, minimum height=1cm, inner sep=3pt},
  action/.style={draw=popgreen, fill=popgreenL, rounded corners=2pt, text width=4.2cm, align=center, minimum height=0.9cm, inner sep=3pt},
  lab/.style={font=\footnotesize\bfseries, midway},
  node distance=0.55cm and 1.6cm
]
\node[question] (q1) {Q1 : fonction \textbf{usuelle} du tableau ?};
\node[action, right=of q1] (a1) {Appliquer la formule du tableau};
\node[question, below=of q1] (q2) {Q2 : \textbf{somme} de termes ?};
\node[action, right=of q2] (a2) {Primitiver terme à terme};
\node[question, below=of q2] (q3) {Q3 : forme $u'u^{r}$, $\dfrac{u'}{u}$ ou $u'e^{u}$ ?};
\node[action, right=of q3] (a3) {Ajuster la constante, appliquer la formule};
\node[question, below=of q3] (q4) {Q4 : puissance de $\cos$ ou $\sin$ ?};
\node[action, right=of q4] (a4) {\textbf{Linéariser} puis primitiver};
\node[action, below=of q4, fill=popgoldL, draw=popgold, align=center] (fin) {Conclure : $F(x) = \ldots + c$ \\ puis \textbf{vérifier en dérivant}};

\draw[-{Stealth}, thick, popdark] (q1) -- node[lab, above]{oui} (a1);
\draw[-{Stealth}, thick, popdark] (q2) -- node[lab, above]{oui} (a2);
\draw[-{Stealth}, thick, popdark] (q3) -- node[lab, above]{oui} (a3);
\draw[-{Stealth}, thick, popdark] (q4) -- node[lab, above]{oui} (a4);
\draw[-{Stealth}, thick, popdark] (q1) -- node[lab, right]{non} (q2);
\draw[-{Stealth}, thick, popdark] (q2) -- node[lab, right]{non} (q3);
\draw[-{Stealth}, thick, popdark] (q3) -- node[lab, right]{non} (q4);
\draw[-{Stealth}, thick, popdark] (a1.east) -- ++(0.7,0) |- (fin.east);
\draw[-{Stealth}, thick, popdark] (a2.east) -- ++(0.7,0) |- (fin.east);
\draw[-{Stealth}, thick, popdark] (a3.east) -- ++(0.7,0) |- (fin.east);
\draw[-{Stealth}, thick, popdark] (a4) -- (fin);
\end{tikzpicture}
\end{center}
\legende{Organigramme de décision : la stratégie complète de recherche d'une primitive.}
\end{popfigure}

\begin{attention}[Rédaction type Bac : les cinq réflexes du candidat sérieux]
Une question de recherche de primitive se rédige en cinq temps :
\begin{enumerate}
\item \textbf{Annoncer l'intervalle} $I$ sur lequel on travaille et justifier la continuité de $f$ sur $I$ (donc l'existence de primitives).
\item \textbf{Identifier la forme} : « on reconnaît la forme $u'\,u^{r}$ avec $u(x) = \ldots$ ».
\item \textbf{Ajuster} la constante multiplicative en écrivant explicitement $f(x) = k \times u'(x)\,u(x)^{r}$.
\item \textbf{Conclure} avec la famille de primitives $F(x) = \ldots + c$, $c \in \mathbb{R}$, et déterminer $c$ si une condition initiale est donnée.
\item \textbf{Vérifier} mentalement (ou au brouillon) en dérivant $F$ : on doit retrouver $f$.
\end{enumerate}
Oublier le $+c$ ou l'intervalle coûte des points, même quand le calcul est juste.
\end{attention}

\courssub{Lecture graphique : de $f$ à $F$, le lien des pentes}

Il est précieux de \emph{voir} ce que signifie « $F$ est une primitive de $f$ » : cela veut dire qu'\textbf{en chaque point, la pente de la tangente à la courbe de $F$ est égale à la valeur de $f$}. Prenons $f(x) = \dfrac{1}{x}$ sur $]0\,;\,+\infty[$ et sa primitive $F(x) = \ln x$ : là où $f$ est grande (près de $0$), la courbe de $\ln$ monte très raide ; là où $f$ devient petite (pour $x$ grand), la courbe de $\ln$ s'aplatit. Et au point d'abscisse $1$, $f(1) = 1$ : la tangente à la courbe de $\ln$ en $x = 1$ a pour coefficient directeur exactement $1$.

\begin{popfigure}
\begin{center}
\begin{tikzpicture}
\begin{axis}[
  width=0.95\linewidth, height=7.2cm,
  axis lines=middle,
  xlabel={$x$}, ylabel={$y$},
  xmin=0, xmax=5.2, ymin=-2.6, ymax=3.2,
  xtick={1,2,3,4,5}, ytick={-2,-1,1,2,3},
  grid=both, grid style={popline!40},
  legend style={at={(0.97,0.03)}, anchor=south east, font=\small, draw=popline}
]
\addplot[popcurve, popink, domain=0.32:5, samples=200]{1/x};
\addlegendentry{$f(x)=\dfrac{1}{x}$}
\addplot[popcurve, popblue, domain=0.08:5, samples=200]{ln(x)};
\addlegendentry{$F(x)=\ln x$}
\addplot[poporange, thick, domain=0.3:1.9, samples=2]{x-1};
\addlegendentry{tangente à $F$ en $x=1$ (pente $f(1)=1$)}
\addplot[only marks, mark=*, poporange] coordinates {(1,0)};
\end{axis}
\end{tikzpicture}
\end{center}
\legende{$F' = f$ se lit graphiquement : la pente de la courbe de $F = \ln$ en tout point d'abscisse $x$ vaut $f(x) = \dfrac{1}{x}$. En $x = 1$, la tangente (orange) a pour pente $f(1) = 1$.}
\end{popfigure}

\begin{aretenir}[Les deux formes à retenir de cette section]
Pour $u$ dérivable sur $I$ :
\[ \frac{u'}{u} \;\xrightarrow{\ \text{primitive}\ }\; \ln|u| + c \quad (u \neq 0 \text{ sur } I) \qquad\qquad u'\,e^{u} \;\xrightarrow{\ \text{primitive}\ }\; e^{u} + c. \]
Et le réflexe permanent : \textbf{après avoir primitivé, dérive pour vérifier}.
\end{aretenir}

\courssub{Exercice de synthèse guidé}

\begin{exoresolu}[Primitive avec condition initiale — sujet type Bac]
On considère la fonction $f$ définie sur $J = \left]\dfrac{-3+\sqrt{5}}{2}\,;\,+\infty\right[$ par :
\[ f(x) = \frac{2x+3}{x^2+3x+1}. \]
Déterminer la primitive $F$ de $f$ sur $J$ telle que $F(0) = 1$. (Remarque : $J$ contient $0$, ce qui permet d'imposer une condition en $0$.)
\tcblower
\textbf{Solution détaillée.}
\begin{enumerate}
\item \textbf{Intervalle et continuité.} Le trinôme $x^2+3x+1$ a pour discriminant $\Delta = 9-4 = 5$ et pour racines $\dfrac{-3\pm\sqrt{5}}{2}$. Sur $J = \left]\dfrac{-3+\sqrt{5}}{2}\,;\,+\infty\right[$ (qui contient $0$), on a $x^2+3x+1 > 0$ : le dénominateur ne s'annule pas, donc $f$ est continue sur $J$ et y admet des primitives.
\item \textbf{Identification de la forme.} Posons $u(x) = x^2+3x+1$. Alors $u'(x) = 2x+3$ : c'est \emph{exactement} le numérateur ! On reconnaît la forme $\dfrac{u'}{u}$, sans aucun ajustement :
\[ f(x) = \frac{u'(x)}{u(x)}. \]
\item \textbf{Conclusion générale.} Comme $u > 0$ sur $J$, les primitives de $f$ sur $J$ sont :
\[ F(x) = \ln\left(x^2+3x+1\right) + c, \qquad c \in \mathbb{R}. \]
\item \textbf{Condition initiale.} $F(0) = \ln(1) + c = c$. La condition $F(0) = 1$ donne $c = 1$.
\item \textbf{Vérification.} $F'(x) = \dfrac{2x+3}{x^2+3x+1} = f(x)$ \checkmark \ et $F(0) = \ln 1 + 1 = 1$ \checkmark.
\end{enumerate}
\textbf{Réponse :} $\displaystyle F(x) = \ln\left(x^2+3x+1\right) + 1$.
\end{exoresolu}

\begin{exercice}[Pour t'entraîner]
\begin{enumerate}
\item Déterminer les primitives de $f(x) = \dfrac{4x}{x^2+1}$ sur $\mathbb{R}$.
\item Déterminer la primitive de $g(x) = (2x+1)\,e^{x^2+x}$ sur $\mathbb{R}$ qui s'annule en $0$.
\item Déterminer les primitives de $h(x) = \dfrac{1}{x\,\ln x}$ sur $]1\,;\,+\infty[$ (indice : poser $u = \ln x$).
\end{enumerate}
\end{exercice}

\begin{corrige}[Exercice : Pour t'entraîner]
\begin{enumerate}
\item Avec $u(x) = x^2+1$, $u'(x) = 2x$, donc $f(x) = 2\,\dfrac{u'(x)}{u(x)}$ et $u > 0$ sur $\mathbb{R}$ : $F(x) = 2\ln(x^2+1) + c$.
\item Avec $u(x) = x^2+x$, $u'(x) = 2x+1$ : forme $u'e^u$ exacte, donc $G(x) = e^{x^2+x} + c$. La condition $G(0) = 0$ donne $1 + c = 0$, soit $c = -1$ : $G(x) = e^{x^2+x} - 1$.
\item Avec $u(x) = \ln x$, $u'(x) = \dfrac{1}{x}$, donc $h(x) = \dfrac{u'(x)}{u(x)}$ et $\ln x > 0$ sur $]1\,;\,+\infty[$ : $H(x) = \ln(\ln x) + c$.
\end{enumerate}
\end{corrige}

\courssub{Vers la leçon suivante : le calcul intégral}

\begin{savaistu}[À quoi servent vraiment les primitives au Bac ?]
Au Baccalauréat C et D, la recherche de primitives n'est presque jamais une question isolée : c'est un \textbf{outil} au service de problèmes plus riches. Elle intervient dans le \textbf{calcul d'intégrales} et d'\textbf{aires} sous une courbe (par exemple l'aire du domaine délimité par la courbe d'une fonction de coût marginal, très utilisée en économie), dans le calcul de \textbf{volumes de solides de révolution}, et dans la résolution des \textbf{équations différentielles} $y' = ay + b$ qui modélisent des phénomènes concrets : la désintégration radioactive, le refroidissement d'un liquide, la charge d'un condensateur, ou encore la croissance d'une population. C'est précisément l'objet de la leçon suivante : nous y définirons l'intégrale d'une fonction continue sur un segment et nous montrerons que l'aire sous la courbe de $f$ entre $a$ et $b$ se calcule par la formule fondamentale $F(b) - F(a)$, où $F$ est une primitive quelconque de $f$. Tout le travail accompli dans cette leçon deviendra alors un véritable calculateur d'aires.
\end{savaistu}

Tu disposes désormais de la panoplie complète : tableau des primitives usuelles, linéarité, formes $u'\,u^{r}$, $\dfrac{u'}{u}$, $u'\,e^{u}$, et linéarisation trigonométrique. Avec l'organigramme de décision en tête et le réflexe de vérification par dérivation, aucune primitive du programme ne te résistera. Rendez-vous au calcul intégral !

\newpage

\coursec{Activité d'intégration}

\courssub{Objectifs de l'activité}

À l'issue de cette activité d'intégration, tu devras être capable de :
\begin{itemize}
\item modéliser une situation concrète (remplissage d'une cuve) par un problème de recherche de primitive ;
\item justifier qu'une fonction inconnue est une \cle{primitive} d'une fonction donnée, en t'appuyant sur la définition ;
\item déterminer l'\cle{ensemble des primitives} d'une fonction somme, en reconnaissant en particulier la forme $u'\cdot u^{r}$ ;
\item utiliser une \cle{condition initiale} pour sélectionner l'unique primitive adaptée au problème ;
\item exploiter la fonction obtenue pour calculer des valeurs et résoudre (par tâtonnement) une équation issue du contexte ;
\item \cle{vérifier} un résultat en dérivant la fonction trouvée, démarche de contrôle indispensable au Baccalauréat.
\end{itemize}

\courssub{Situation}

\textbf{Le réservoir de la SONAF à Bafoussam.}

Dans le cadre de la distribution d'eau minérale dans la ville de Bafoussam (région de l'Ouest), une station de remplissage de la SONAF utilise une cuve de stockage. À l'ouverture de la station le matin (instant $t = 0$), la cuve contient déjà $200$ litres d'eau. Une pompe alimente alors la cuve avec un débit qui varie au cours du temps : le débit est fort au démarrage, puis se stabilise progressivement.

On modélise le débit d'alimentation (en litres par heure) à l'instant $t$ (exprimé en heures, $t \geq 0$) par la fonction :
\[ d(t) = 6t + 4 + \frac{12}{(t+1)^{2}}. \]
On note $Q(t)$ la quantité d'eau (en litres) contenue dans la cuve à l'instant $t$. On admet que le débit est la dérivée de la quantité d'eau : $Q'(t) = d(t)$ pour tout $t \geq 0$.

\textbf{Problème posé par le gestionnaire de la station :} quelle quantité d'eau la cuve contiendra-t-elle au bout de $3$ heures de fonctionnement, et à quel moment la cuve atteindra-t-elle le niveau de $300$ litres, seuil à partir duquel une vanne de sécurité doit être ouverte ?

\courssub{Consignes}

Travail en groupes de 3 à 4 élèves. Chaque groupe rédige sa solution sur une feuille, puis un rapporteur présente les résultats au tableau pour une correction collective portant aussi sur la \emph{qualité de la rédaction}.

\begin{enumerate}
\item \textbf{Modélisation.} Expliquer, à l'aide de la définition d'une primitive, pourquoi la fonction $Q$ est une primitive de la fonction $d$ sur $[0\,; +\infty[$. Vérifier au préalable que $d$ est continue sur cet intervalle.
\item \textbf{Ensemble des primitives.} Déterminer l'ensemble des primitives de la fonction $d$ sur $[0\,; +\infty[$. \emph{Indication :} écrire $\dfrac{12}{(t+1)^{2}} = 12(t+1)^{-2}$ et reconnaître la forme $u'\cdot u^{r}$ avec $u(t) = t+1$.
\item \textbf{Condition initiale.} Déterminer l'unique primitive $Q$ de $d$ qui vérifie $Q(0) = 200$.
\item \textbf{Exploitation.}
\begin{enumerate}
\item Calculer la quantité d'eau contenue dans la cuve au bout de $3$ heures.
\item Déterminer, par tâtonnement à la calculatrice (tableau de valeurs), un encadrement à $0{,}1$ près de l'instant $t_{0}$ où la cuve atteint $300$ litres.
\end{enumerate}
\item \textbf{Vérification.} Dériver la fonction $Q$ obtenue à la question 3 et vérifier que l'on retrouve bien $d(t)$, ainsi que $Q(0) = 200$.
\item \textbf{Interprétation.} Rédiger une phrase de conclusion à l'attention du gestionnaire de la station, répondant à son problème.
\end{enumerate}

\courssub{Ressources à mobiliser}

\begin{aretenir}[Boîte à outils de la leçon]
\begin{itemize}
\item \textbf{Définition :} $F$ est une primitive de $f$ sur un intervalle $I$ si $F$ est dérivable sur $I$ et si $F'(x) = f(x)$ pour tout $x \in I$. Toute fonction \cle{continue} sur un intervalle y admet des primitives.
\item \textbf{Ensemble des primitives :} si $F$ est une primitive de $f$ sur $I$, alors les primitives de $f$ sur $I$ sont exactement les fonctions $F + k$, $k \in \mathbb{R}$.
\item \textbf{Condition initiale :} pour tout $x_{0} \in I$ et tout $y_{0} \in \mathbb{R}$, il existe une \emph{unique} primitive $F$ de $f$ vérifiant $F(x_{0}) = y_{0}$.
\item \textbf{Primitives usuelles :} une primitive de $t^{n}$ ($n \in \mathbb{N}$) est $\dfrac{t^{n+1}}{n+1}$ ; une primitive de la constante $a$ est $at$.
\item \textbf{Forme composée :} une primitive de $u'\cdot u^{r}$ ($r$ rationnel, $r \neq -1$, avec $u$ qui ne s'annule pas si $r < 0$) est $\dfrac{u^{r+1}}{r+1}$. En particulier, une primitive de $\dfrac{u'}{u^{2}}$ est $-\dfrac{1}{u}$.
\item \textbf{Linéarité :} une primitive d'une somme est la somme des primitives ; une primitive de $\lambda f$ est $\lambda F$.
\end{itemize}
\end{aretenir}

\courssub{Démarche et corrigé commenté}

\begin{exoresolu}[Le réservoir de la SONAF — corrigé complet]
Énoncé : $d(t) = 6t + 4 + \dfrac{12}{(t+1)^{2}}$ (débit en L/h), $Q'(t) = d(t)$, $Q(0) = 200$. Déterminer $Q$, calculer $Q(3)$, trouver l'instant où $Q(t) = 300$.

\tcblower

\textbf{Question 1 — $Q$ est une primitive de $d$.}

La fonction $t \mapsto 6t+4$ est polynomiale, donc continue sur $\mathbb{R}$. La fonction $t \mapsto \dfrac{12}{(t+1)^{2}}$ est un quotient de fonctions continues dont le dénominateur $(t+1)^{2}$ ne s'annule pas sur $[0\,; +\infty[$ (car pour $t \geq 0$, on a $t+1 \geq 1 > 0$) : elle y est donc continue. Par somme, $d$ est \cle{continue} sur $[0\,; +\infty[$.

Or, par hypothèse de la modélisation, $Q$ est dérivable sur $[0\,; +\infty[$ et $Q'(t) = d(t)$ pour tout $t \geq 0$. Par \textbf{définition}, $Q$ est donc une \cle{primitive} de $d$ sur $[0\,; +\infty[$. L'existence d'une telle fonction $Q$ est garantie par le théorème : toute fonction continue sur un intervalle y admet des primitives.

\medskip
\textbf{Question 2 — Ensemble des primitives de $d$.}

On cherche une primitive de chacun des trois termes.
\begin{itemize}
\item Une primitive de $6t$ est $6 \times \dfrac{t^{2}}{2} = 3t^{2}$.
\item Une primitive de la constante $4$ est $4t$.
\item Pour $\dfrac{12}{(t+1)^{2}}$, on pose $u(t) = t+1$, donc $u'(t) = 1$. Alors
\[ \frac{12}{(t+1)^{2}} = 12\,\frac{u'(t)}{[u(t)]^{2}} = 12\,u'(t)\,[u(t)]^{-2}. \]
C'est la forme $u'\cdot u^{r}$ avec $r = -2 \neq -1$, et $u(t) = t+1 \geq 1$ ne s'annule pas sur $[0\,; +\infty[$. Une primitive est
\[ 12 \times \frac{[u(t)]^{-2+1}}{-2+1} = 12 \times \frac{[u(t)]^{-1}}{-1} = -\frac{12}{u(t)} = -\frac{12}{t+1}. \]
\end{itemize}
Par linéarité, l'ensemble des primitives de $d$ sur $[0\,; +\infty[$ est l'ensemble des fonctions :
\[ F_{k}(t) = 3t^{2} + 4t - \frac{12}{t+1} + k, \qquad k \in \mathbb{R}. \]

\medskip
\textbf{Question 3 — Condition initiale $Q(0) = 200$.}

On calcule :
\[ F_{k}(0) = 3(0)^{2} + 4(0) - \frac{12}{0+1} + k = -12 + k. \]
La condition $Q(0) = 200$ s'écrit $-12 + k = 200$, d'où $k = 212$. Il existe une \emph{unique} primitive convenable (théorème de la condition initiale) :
\[ Q(t) = 3t^{2} + 4t - \frac{12}{t+1} + 212. \]

\medskip
\textbf{Question 4 — Exploitation.}

\textbf{a)} Quantité d'eau au bout de $3$ heures :
\begin{align*}
Q(3) &= 3(3)^{2} + 4(3) - \frac{12}{3+1} + 212 \\
&= 27 + 12 - 3 + 212 = 248.
\end{align*}
Au bout de $3$ heures, la cuve contient $\mathbf{248}$ \textbf{litres}.

\textbf{b)} Instant où la cuve atteint $300$ litres.

On résout l'équation $Q(t) = 300$ par tâtonnement, à l'aide du tableau de valeurs de la calculatrice. On calcule d'abord quelques valeurs :
\[ Q(4) = 3(16) + 16 - \frac{12}{5} + 212 = 48 + 16 - 2{,}4 + 212 = 273{,}6 \ ; \]
\[ Q(5) = 3(25) + 20 - \frac{12}{6} + 212 = 75 + 20 - 2 + 212 = 305. \]
Comme $Q(4) = 273{,}6 < 300 < 305 = Q(5)$ et que $Q$ est strictement croissante sur $[0\,; +\infty[$ (sa dérivée $d$ est strictement positive, comme somme de termes strictement positifs), l'équation $Q(t) = 300$ admet une unique solution $t_{0}$ et $t_{0} \in [4\,; 5]$. On affine au dixième :
\begin{align*}
Q(4{,}8) &= 3(4{,}8)^{2} + 4(4{,}8) - \frac{12}{5{,}8} + 212 \\
&= 69{,}12 + 19{,}2 - 2{,}069 + 212 \approx 298{,}25, \\[4pt]
Q(4{,}9) &= 3(4{,}9)^{2} + 4(4{,}9) - \frac{12}{5{,}9} + 212 \\
&= 72{,}03 + 19{,}6 - 2{,}034 + 212 \approx 301{,}60.
\end{align*}
Comme $Q(4{,}8) < 300 < Q(4{,}9)$, on obtient l'encadrement $4{,}8 < t_{0} < 4{,}9$. La cuve atteint $300$ litres au bout d'environ $\mathbf{4{,}8}$ \textbf{heures}, soit environ $4$ h $48$ min.

\medskip
\textbf{Question 5 — Vérification par dérivation.}

$Q$ est dérivable sur $[0\,; +\infty[$ comme somme de fonctions dérivables, et :
\begin{align*}
Q'(t) &= 3 \times 2t + 4 - 12 \times \left(-\frac{1}{(t+1)^{2}}\right) + 0 \\
&= 6t + 4 + \frac{12}{(t+1)^{2}} = d(t). \quad \checkmark
\end{align*}
De plus $Q(0) = 0 + 0 - 12 + 212 = 200$. \quad $\checkmark$

\medskip
\textbf{Question 6 — Conclusion pour le gestionnaire.}

Au bout de $3$ heures de fonctionnement, la cuve contient $248$ litres d'eau minérale. Le niveau critique de $300$ litres est atteint au bout d'environ $4$ h $48$ min : le gestionnaire doit donc prévoir l'ouverture de la vanne de sécurité avant cet instant.
\end{exoresolu}

\begin{attention}[Pièges relevés lors de la correction collective]
\begin{itemize}
\item \textbf{Oubli de la constante $k$ :} écrire directement $Q(t) = 3t^{2} + 4t - \dfrac{12}{t+1}$ sans constante, c'est choisir \emph{une} primitive parmi une infinité, et elle ne vérifie pas $Q(0) = 200$ (elle donne $Q(0) = -12$ !). La condition initiale est indispensable pour déterminer $k$.
\item \textbf{Erreur de signe :} une primitive de $\dfrac{12}{(t+1)^{2}}$ est $-\dfrac{12}{t+1}$ et non $+\dfrac{12}{t+1}$. En cas de doute, dérive ta réponse : c'est le réflexe de vérification.
\item \textbf{Confusion entre dérivation et primitivation :} pour une primitive de $t^{n}$, on \emph{augmente} l'exposant de $1$ et on \emph{divise} par le nouvel exposant ; c'est l'inverse de la dérivation.
\item \textbf{Condition d'application :} la formule de $u'\cdot u^{r}$ exige $r \neq -1$ et, si $r < 0$, que $u$ ne s'annule pas sur l'intervalle. Ici $u(t) = t+1 \geq 1$ : la condition est bien remplie, et il faut le dire.
\item \textbf{Rédaction :} annoncer la forme $u'\cdot u^{r}$ en nommant explicitement $u$ et $r$ est une justification attendue au Baccalauréat.
\end{itemize}
\end{attention}

\courssub{Bilan de l'activité}

Cette activité a permis de mobiliser, dans un cadre concret, l'ensemble des savoir-faire de la leçon :
\begin{itemize}
\item le lien entre une grandeur et son taux de variation s'exprime naturellement par la notion de \cle{primitive} : la quantité d'eau est une primitive du débit ;
\item la recherche de primitives exige de \emph{décomposer} la fonction et de \emph{reconnaître des formes} : ici, un polynôme et la forme $u'\cdot u^{-2}$ ;
\item une condition initiale ($Q(0) = 200$) permet de sélectionner, parmi l'infinité de primitives, l'\emph{unique} fonction adaptée au problème physique ;
\item la fonction obtenue s'exploite ensuite pour calculer des valeurs et résoudre des équations par tâtonnement, en justifiant l'existence et l'unicité de la solution grâce à la stricte monotonie ;
\item enfin, la \cle{vérification par dérivation} est apparue comme un outil de contrôle simple et systématique : dériver la réponse doit devenir un réflexe.
\end{itemize}

\begin{savaistu}[Pour aller plus loin]
La démarche utilisée ici — connaître un débit (une vitesse de variation) et en déduire la quantité accumulée — est au cœur de nombreux métiers : gestion des ressources en eau à la SONAF ou à la CAMWATER, suivi de la consommation électrique par ENEO, estimation des stocks dans les silos à céréales de la SODECOTON. En classe de Terminale, elle trouvera son prolongement naturel dans le chapitre sur le \emph{calcul intégral}, où l'on apprendra que la quantité totale écoulée entre deux instants s'obtient par une \emph{intégrale} du débit.
\end{savaistu}

\newpage

\coursec{Méthodes et exercices résolus}

\coursec{Primitives d'une fonction continue}

\courssub{Notion de primitive et existence}

\begin{definition}[Primitive d'une fonction]
Soit $I$ un intervalle de $\mathbb{R}$ et $f$ une fonction définie sur $I$. On appelle \cle{primitive} de $f$ sur $I$ toute fonction $F$ dérivable sur $I$ telle que $F'(x) = f(x)$ pour tout $x \in I$.
\end{definition}

\begin{propriete}[Existence des primitives]
Toute fonction \cle{continue} sur un intervalle $I$ admet au moins une primitive sur $I$.
\end{propriete}
\begin{attention}
La continuité sur un \emph{intervalle} est une condition suffisante (et nécessaire pour l'existence sur tout l'intervalle). Une fonction continue par morceaux admet des primitives sur chaque intervalle de continuité, mais pas nécessairement sur la réunion.
\end{attention}

\begin{propriete}[Ensemble des primitives]
Si $F_0$ est une primitive de $f$ sur $I$, alors l'ensemble des primitives de $f$ sur $I$ est :
\[ \{ F_0 + C \mid C \in \mathbb{R} \}. \]
Deux primitives d'une même fonction sur un intervalle ne diffèrent que d'une constante.
\end{propriete}

\begin{aretenir}[Primitive et condition initiale]
Soit $f$ continue sur $I$, $x_0 \in I$ et $y_0 \in \mathbb{R}$. Il existe une \cle{unique} primitive $F$ de $f$ sur $I$ telle que $F(x_0) = y_0$. Elle s'obtient en fixant la constante $C = y_0 - F_0(x_0)$ où $F_0$ est une primitive quelconque de $f$.
\end{aretenir}

\courssub{Ensemble des primitives ; condition initiale}

\begin{methode}[Déterminer la primitive vérifiant une condition initiale]
Pour déterminer \textbf{la} primitive $F$ de $f$ sur $I$ telle que $F(x_0) = y_0$ :
\begin{enumerate}
\item \textbf{Justifier l'existence} : $f$ continue sur $I$ et $x_0 \in I$ $\implies$ existence et unicité.
\item \textbf{Calculer l'ensemble des primitives} : $F(x) = F_0(x) + C,\quad C \in \mathbb{R}$.
\item \textbf{Traduire la condition} : $F(x_0) = y_0 \iff F_0(x_0) + C = y_0$.
\item \textbf{Résoudre} pour $C$ : $C = y_0 - F_0(x_0)$.
\item \textbf{Conclure} en donnant $F(x) = F_0(x) + C$ (sans $+C$ final).
\end{enumerate}
\end{methode}

\begin{attention}[Piège de rédaction]
La condition initiale ne se vérifie qu'\emph{après} avoir écrit l'ensemble des primitives $F_0(x)+C$. Écrire directement une primitive particulière sans constante, puis « ajuster » est une démarche mal rédigée et source d'erreurs : passe toujours par l'étape $F_0(x) + C$.
\end{attention}

\begin{exoresolu}[Primitive avec condition initiale (type Bac)]
\textbf{Énoncé.} Soit $f$ la fonction définie sur $\mathbb{R}$ par $f(x) = 2\cos(2x) - 3\sin(3x)$. Déterminer la primitive $F$ de $f$ sur $\mathbb{R}$ telle que $F(0) = 1$.
\tcblower
\textbf{Solution détaillée.}

\textbf{Étape 1 — Existence.} Les fonctions $x \mapsto \cos(2x)$ et $x \mapsto \sin(3x)$ sont continues sur $\mathbb{R}$, donc $f$ est continue sur $\mathbb{R}$. Elle admet des primitives sur $\mathbb{R}$, et une unique primitive vérifiant $F(0) = 1$.

\textbf{Étape 2 — Ensemble des primitives.} On utilise les formules :
\[ \int \cos(ax)\,dx = \frac{1}{a}\sin(ax), \qquad \int \sin(ax)\,dx = -\frac{1}{a}\cos(ax). \]
\begin{align*}
\int 2\cos(2x)\,dx &= 2\times\frac{1}{2}\sin(2x) = \sin(2x), \\
\int \big(-3\sin(3x)\big)\,dx &= -3\times\left(-\frac{1}{3}\cos(3x)\right) = \cos(3x).
\end{align*}
Donc l'ensemble des primitives de $f$ est :
\[ F(x) = \sin(2x) + \cos(3x) + C, \quad C \in \mathbb{R}. \]

\textbf{Étape 3 — Condition initiale.} On impose $F(0) = 1$ :
\[ F(0) = \sin(0) + \cos(0) + C = 0 + 1 + C = 1 \iff C = 0. \]

\textbf{Conclusion.} L'unique primitive de $f$ sur $\mathbb{R}$ vérifiant $F(0) = 1$ est :
\[ \boxed{F(x) = \sin(2x) + \cos(3x).} \]
\emph{Vérification :} $F'(x) = 2\cos(2x) - 3\sin(3x) = f(x)$ et $F(0) = 0 + 1 = 1$. \checkmark
\end{exoresolu}

\courssub{Primitives des fonctions usuelles}

\begin{propriete}[Tableau des primitives usuelles]
Soit $I$ un intervalle où les expressions sont définies. Pour toute constante $a \in \mathbb{R}^*$ et $r \in \mathbb{Q}$ :
\begin{center}
\begin{tabularx}{\linewidth}{|l|X|l|}
\hline
\rowcolor{popblueL} \textbf{Fonction $f(x)$} & \textbf{Primitive $F(x)$ sur $I$} & \textbf{Condition} \\
\hline
$x^r$ & $\dfrac{x^{r+1}}{r+1}$ & $r \neq -1$ \\
\hline
$\dfrac{1}{x}$ & $\ln|x|$ & $0 \notin I$ \\
\hline
$\cos(ax+b)$ & $\dfrac{1}{a}\sin(ax+b)$ & $a \neq 0$ \\
\hline
$\sin(ax+b)$ & $-\dfrac{1}{a}\cos(ax+b)$ & $a \neq 0$ \\
\hline
$e^{ax+b}$ & $\dfrac{1}{a}e^{ax+b}$ & $a \neq 0$ \\
\hline
\end{tabularx}
\end{center}
\emph{Règle de linéarité :} Une primitive de $\lambda f + \mu g$ est $\lambda F + \mu G$ (somme des primitives).
\end{propriete}

\begin{methode}[Trouver toutes les primitives d'une fonction usuelle]
Pour déterminer l'ensemble des primitives d'une fonction $f$ continue sur un intervalle $I$ :
\begin{enumerate}
\item \textbf{Identifier} la forme de $f$ : polynôme, puissance $x^r$, $\dfrac{a}{(cx+d)^r}$, $\cos(ax+b)$, $\sin(ax+b)$, etc.
\item \textbf{Choisir la formule adaptée} dans le tableau ci-dessus.
\item \textbf{Ajuster les coefficients} : si la fonction est de la forme $f(ax+b)$, penser à diviser par $a$ (car la dérivation de $F(ax+b)$ fait apparaître le facteur $a$).
\item \textbf{Écrire l'ensemble des primitives} : $F(x) = F_0(x) + C$, où $C \in \mathbb{R}$ est une constante arbitraire. \textbf{Ne jamais oublier le $+C$}.
\item \textbf{Vérifier mentalement} en dérivant le résultat : on doit retrouver $f$.
\end{enumerate}
\end{methode}

\begin{exoresolu}[Primitives d'une fonction polynôme et d'une puissance rationnelle]
\textbf{Énoncé.} Déterminer l'ensemble des primitives sur l'intervalle indiqué de chacune des fonctions suivantes :
\begin{enumerate}
\item $f(x) = 3x^2 - 4x + 5$ sur $\mathbb{R}$ ;
\item $g(x) = \dfrac{2}{x^3} + \sqrt{x}$ sur $]0\,; +\infty[$.
\end{enumerate}
\tcblower
\textbf{Solution détaillée.}

\textbf{1.} La fonction $f$ est une fonction polynôme, donc continue sur $\mathbb{R}$ : elle admet des primitives sur $\mathbb{R}$. On primitive terme à terme avec la formule $\displaystyle\int x^n\,dx = \frac{x^{n+1}}{n+1}$ :
\begin{align*}
\int 3x^2\,dx &= 3\times\frac{x^3}{3} = x^3, &
\int (-4x)\,dx &= -4\times\frac{x^2}{2} = -2x^2, &
\int 5\,dx &= 5x.
\end{align*}
L'ensemble des primitives de $f$ sur $\mathbb{R}$ est donc :
\[ F(x) = x^3 - 2x^2 + 5x + C, \quad C \in \mathbb{R}. \]
\emph{Vérification :} $F'(x) = 3x^2 - 4x + 5 = f(x)$. \checkmark

\textbf{2.} Sur $]0\,;+\infty[$, la fonction $g$ est continue (somme de fonctions continues). On réécrit $g$ sous forme de puissances pour appliquer la formule $\displaystyle\int x^r\,dx = \frac{x^{r+1}}{r+1}$ :
\[ g(x) = 2x^{-3} + x^{1/2}. \]
\begin{itemize}
\item Pour $2x^{-3}$ : $r = -3 \neq -1$, donc $\displaystyle\int 2x^{-3}\,dx = 2\times\frac{x^{-2}}{-2} = -\frac{1}{x^2}$.
\item Pour $x^{1/2}$ : $r = \tfrac12$, donc $\displaystyle\int x^{1/2}\,dx = \frac{x^{3/2}}{3/2} = \frac{2}{3}x^{3/2} = \frac{2}{3}x\sqrt{x}$.
\end{itemize}
L'ensemble des primitives de $g$ sur $]0\,;+\infty[$ est :
\[ G(x) = -\frac{1}{x^2} + \frac{2}{3}x\sqrt{x} + C, \quad C \in \mathbb{R}. \]
\emph{Vérification :} $G'(x) = \dfrac{2}{x^3} + \dfrac{2}{3}\times\dfrac{3}{2}x^{1/2} = \dfrac{2}{x^3} + \sqrt{x} = g(x)$. \checkmark

\textbf{Conclusion.} Les primitives de $f$ sont les fonctions $x \mapsto x^3 - 2x^2 + 5x + C$, et celles de $g$ sont les fonctions $x \mapsto -\dfrac{1}{x^2} + \dfrac{2}{3}x\sqrt{x} + C$, avec $C \in \mathbb{R}$.
\end{exoresolu}

\courssub{Primitives des formes composées : $u'\cdot u^r$ et $\dfrac{a}{(cx+d)^r}$}

\begin{methode}[Primitiver une puissance d'une fonction affine $\dfrac{a}{(cx+d)^r}$]
Pour primitiver $f(x) = \dfrac{a}{(cx+d)^r} = a(cx+d)^{-r}$ avec $c \neq 0$ et $r \in \mathbb{Q}$, $r \neq 1$ (cas $r=1$ : voir forme $\frac{u'}{u}$) :
\begin{enumerate}
\item \textbf{Reconnaître la forme} $u^{-r}$ avec $u(x) = cx+d$ et $u'(x) = c$.
\item \textbf{Faire apparaître $u'$} : écrire $f(x) = \dfrac{a}{c}\times c\,(cx+d)^{-r} = \dfrac{a}{c}\,u'(x)\,u(x)^{-r}$.
\item \textbf{Appliquer la formule} $\displaystyle\int u'\,u^{n}\,dx = \frac{u^{n+1}}{n+1}$ avec $n = -r$ :
\[ \int \frac{a}{(cx+d)^r}\,dx = \frac{a}{c}\times\frac{(cx+d)^{-r+1}}{-r+1} = \frac{a}{c(1-r)\,(cx+d)^{r-1}}. \]
\item \textbf{Préciser l'intervalle} : la fonction n'est définie que si $cx+d \neq 0$ ; on travaille sur un intervalle ne contenant pas $-\dfrac{d}{c}$.
\item \textbf{Ajouter $+C$} et vérifier par dérivation.
\end{enumerate}
\end{methode}

\begin{exoresolu}[Primitives de $\dfrac{a}{(cx+d)^r}$]
\textbf{Énoncé.} 
\begin{enumerate}
\item Déterminer l'ensemble des primitives de $f(x) = \dfrac{6}{(2x+1)^2}$ sur $\left]-\dfrac12\,; +\infty\right[$.
\item En déduire la primitive $F$ de $f$ telle que $F(0) = 0$.
\end{enumerate}
\tcblower
\textbf{Solution détaillée.}

\textbf{1.} Sur $\left]-\tfrac12\,;+\infty\right[$, on a $2x+1 > 0$, donc $f$ est bien définie et continue sur cet intervalle.

On reconnaît la forme $u^{-2}$ avec $u(x) = 2x+1$, donc $u'(x) = 2$. On fait apparaître $u'$ :
\[ f(x) = \frac{6}{(2x+1)^2} = 3\times\frac{2}{(2x+1)^2} = 3\,u'(x)\,u(x)^{-2}. \]
On applique la formule $\displaystyle\int u'u^{-2}\,dx = \frac{u^{-1}}{-1} = -\frac{1}{u}$ :
\[ \int f(x)\,dx = 3\times\left(-\frac{1}{2x+1}\right) + C = -\frac{3}{2x+1} + C, \quad C \in \mathbb{R}. \]
\emph{Vérification :} $\left(-\dfrac{3}{2x+1}\right)' = -3\times\left(-\dfrac{2}{(2x+1)^2}\right) = \dfrac{6}{(2x+1)^2} = f(x)$. \checkmark

\textbf{2.} On impose $F(0) = 0$ :
\[ F(0) = -\frac{3}{2\times 0 + 1} + C = -3 + C = 0 \iff C = 3. \]

\textbf{Conclusion.} L'ensemble des primitives de $f$ est $x \mapsto -\dfrac{3}{2x+1} + C$, et la primitive cherchée est :
\[ \boxed{F(x) = -\frac{3}{2x+1} + 3 = \frac{6x}{2x+1}.} \]
\end{exoresolu}

\begin{methode}[Primitiver une forme composée $u' \cdot u^r$]
C'est la compétence-clé du Bac. Pour primitiver $f$ :
\begin{enumerate}
\item \textbf{Repérer une « fonction intérieure »} $u$ : souvent sous une puissance, une racine, un cosinus... Poser $u(x)$ et calculer $u'(x)$.
\item \textbf{Comparer} : l'expression contient-elle $u'(x)$ en facteur, à une constante multiplicative près ?
\item \textbf{Ajuster la constante} : si $f = k \cdot u' \cdot u^r$, écrire $f = k\,u'u^r$ et appliquer
\[ \int u'\,u^r\,dx = \frac{u^{r+1}}{r+1} \quad (r \neq -1). \]
\item \textbf{Conclure} avec $+C$, sur un intervalle où $u$ ne s'annule pas si $r < 0$.
\item \textbf{Vérifier} en dérivant : la dérivée de $\dfrac{u^{r+1}}{r+1}$ est $u'u^r$ (dérivation des composées).
\end{enumerate}
\textbf{Réflexe.} Quand tu vois un produit « dérivée $\times$ puissance », pense immédiatement à $u'u^r$.
\end{methode}

\begin{exoresolu}[Reconnaissance de la forme $u'u^r$ (type Bac)]
\textbf{Énoncé.} Déterminer l'ensemble des primitives des fonctions suivantes sur l'intervalle indiqué :
\begin{enumerate}
\item $f(x) = 2x\left(x^2+3\right)^4$ sur $\mathbb{R}$ ;
\item $g(x) = \dfrac{3x^2}{\sqrt{x^3+1}}$ sur $]-1\,;+\infty[$.
\end{enumerate}
\tcblower
\textbf{Solution détaillée.}

\textbf{1.} Posons $u(x) = x^2+3$ ; alors $u'(x) = 2x$. La fonction $f$ s'écrit exactement :
\[ f(x) = u'(x)\,\big(u(x)\big)^4. \]
C'est la forme $u'u^r$ avec $r = 4 \neq -1$. Donc :
\[ \int f(x)\,dx = \frac{u(x)^5}{5} + C = \frac{\left(x^2+3\right)^5}{5} + C, \quad C \in \mathbb{R}. \]
\emph{Vérification :} $\left(\dfrac{(x^2+3)^5}{5}\right)' = \dfrac{5(x^2+3)^4 \times 2x}{5} = 2x(x^2+3)^4 = f(x)$. \checkmark

\textbf{2.} Sur $]-1\,;+\infty[$, $x^3+1 > 0$, donc $g$ est bien définie et continue. Réécrivons :
\[ g(x) = 3x^2\left(x^3+1\right)^{-1/2}. \]
Posons $u(x) = x^3+1$ ; alors $u'(x) = 3x^2$, qui apparaît exactement au numérateur. Donc $g(x) = u'(x)\,u(x)^{-1/2}$, forme $u'u^r$ avec $r = -\tfrac12 \neq -1$ :
\[ \int g(x)\,dx = \frac{u(x)^{1/2}}{1/2} + C = 2\sqrt{x^3+1} + C, \quad C \in \mathbb{R}. \]
\emph{Vérification :} $\left(2\sqrt{x^3+1}\right)' = 2\times\dfrac{3x^2}{2\sqrt{x^3+1}} = \dfrac{3x^2}{\sqrt{x^3+1}} = g(x)$. \checkmark

\textbf{Conclusion.} Les primitives de $f$ sont les fonctions $x \mapsto \dfrac{(x^2+3)^5}{5} + C$ et celles de $g$ sont les fonctions $x \mapsto 2\sqrt{x^3+1} + C$, avec $C \in \mathbb{R}$.
\end{exoresolu}

\courssub{Primitives de $\cos^n x$ et $\sin^n x$}

\begin{methode}[Primitiver $\cos^n x$ et $\sin^n x$]
Deux cas à distinguer :
\begin{enumerate}
\item \textbf{Si $n$ est impair} ($n = 2k+1$, $k \ge 1$) : on utilise $\cos^2 x + \sin^2 x = 1$ pour faire apparaître une forme $u'u^r$.
\begin{itemize}
\item Pour $\sin^{2k+1} x = \sin x \,(1-\cos^2 x)^k$ : poser $u = \cos x$, $u' = -\sin x$.
\item Pour $\cos^{2k+1} x = \cos x\,(1-\sin^2 x)^k$ : poser $u = \sin x$, $u' = \cos x$.
\end{itemize}
\item \textbf{Si $n$ est pair} ($n = 2k$) : on \textbf{linéarise} grâce aux formules d'Euler / Moivre :
\[ \cos^2 x = \frac{1+\cos(2x)}{2}, \qquad \sin^2 x = \frac{1-\cos(2x)}{2}, \]
et on réduit les puissances supérieures par récurrence ou formules de linéarisation ($\cos^4 x$, $\sin^4 x$, etc.), puis on primitive les fonctions affines de $\cos(2mx)$.
\end{enumerate}
\textbf{Réflexe.} Puissance impaire $\to$ forme $u'u^r$ ; puissance paire $\to$ linéarisation.
\end{methode}

\begin{exoresolu}[Primitives de $\sin^3 x$ et $\cos^2 x$]
\textbf{Énoncé.} Déterminer l'ensemble des primitives sur $\mathbb{R}$ de :
\begin{enumerate}
\item $f(x) = \sin^3 x$ ;
\item $g(x) = \cos^2 x$.
\end{enumerate}
\tcblower
\textbf{Solution détaillée.}

\textbf{1.} La puissance est impaire ($n=3$). On écrit, grâce à $\sin^2 x = 1 - \cos^2 x$ :
\[ f(x) = \sin x \cdot \sin^2 x = \sin x\left(1 - \cos^2 x\right) = \sin x - \sin x \cos^2 x. \]
\begin{itemize}
\item Une primitive de $\sin x$ est $-\cos x$.
\item Pour $\sin x \cos^2 x$ : posons $u(x) = \cos x$, alors $u'(x) = -\sin x$, donc $\sin x \cos^2 x = -u'(x)\,u(x)^2$, forme $u'u^2$ de primitive $-\dfrac{u^3}{3} = -\dfrac{\cos^3 x}{3}$.
\end{itemize}
Par conséquent :
\[ \int \sin^3 x\,dx = -\cos x - \left(-\frac{\cos^3 x}{3}\right) + C = -\cos x + \frac{\cos^3 x}{3} + C, \quad C \in \mathbb{R}. \]
\emph{Vérification :} $\left(-\cos x + \dfrac{\cos^3 x}{3}\right)' = \sin x + \dfrac{3\cos^2 x \times (-\sin x)}{3} = \sin x (1-\cos^2 x) = \sin^3 x$. \checkmark

\textbf{2.} La puissance est paire ($n=2$) : on linéarise avec $\cos^2 x = \dfrac{1+\cos(2x)}{2}$ :
\[ g(x) = \frac{1}{2} + \frac{1}{2}\cos(2x). \]
On primitive terme à terme, avec $\displaystyle\int \cos(2x)\,dx = \frac{1}{2}\sin(2x)$ :
\[ \int \cos^2 x\,dx = \frac{x}{2} + \frac{1}{2}\times\frac{1}{2}\sin(2x) + C = \frac{x}{2} + \frac{\sin(2x)}{4} + C, \quad C \in \mathbb{R}. \]
\emph{Vérification :} $\left(\dfrac{x}{2} + \dfrac{\sin(2x)}{4}\right)' = \dfrac12 + \dfrac{2\cos(2x)}{4} = \dfrac{1+\cos(2x)}{2} = \cos^2 x$. \checkmark

\textbf{Conclusion.} Les primitives de $\sin^3 x$ sont $x \mapsto -\cos x + \dfrac{\cos^3 x}{3} + C$, et celles de $\cos^2 x$ sont $x \mapsto \dfrac{x}{2} + \dfrac{\sin(2x)}{4} + C$, avec $C \in \mathbb{R}$.
\end{exoresolu}

\courssub{Compléments utiles au Bac : $\dfrac{u'}{u}$ et $u'e^u$ ; synthèse méthodologique}

\begin{methode}[Primitiver $\dfrac{u'}{u}$ et $u'e^u$]
Ces deux formes complètent efficacement la boîte à outils pour le Bac :
\begin{enumerate}
\item \textbf{Forme $\dfrac{u'}{u}$} : si $u$ est dérivable et \textbf{ne s'annule pas} sur $I$, alors
\[ \int \frac{u'(x)}{u(x)}\,dx = \ln|u(x)| + C. \]
Sur un intervalle où $u > 0$, on écrit simplement $\ln\big(u(x)\big) + C$. Cas particulier fréquent : $\displaystyle\int \frac{dx}{cx+d} = \frac{1}{c}\ln|cx+d| + C$.
\item \textbf{Forme $u'e^u$} : si $u$ est dérivable sur $I$,
\[ \int u'(x)\,e^{u(x)}\,dx = e^{u(x)} + C. \]
\item \textbf{Démarche} : identifier $u$, calculer $u'$, ajuster la constante multiplicative pour faire apparaître exactement $u'$, appliquer la formule, vérifier par dérivation.
\end{enumerate}
\end{methode}

\begin{exoresolu}[Formes $\dfrac{u'}{u}$ et $u'e^u$]
\textbf{Énoncé.} Déterminer l'ensemble des primitives de :
\begin{enumerate}
\item $f(x) = \dfrac{4x}{x^2+1}$ sur $\mathbb{R}$ ;
\item $g(x) = x\,e^{x^2}$ sur $\mathbb{R}$.
\end{enumerate}
\tcblower
\textbf{Solution détaillée.}

\textbf{1.} Pour tout $x \in \mathbb{R}$, $x^2+1 \geq 1 > 0$ : $f$ est définie et continue sur $\mathbb{R}$. Posons $u(x) = x^2+1$ ; alors $u'(x) = 2x$. Le numérateur est $4x = 2\times 2x = 2u'(x)$. Donc :
\[ f(x) = 2\,\frac{u'(x)}{u(x)}. \]
Comme $u(x) > 0$ sur $\mathbb{R}$, la formule $\displaystyle\int \frac{u'}{u}\,dx = \ln(u) + C$ donne :
\[ \int f(x)\,dx = 2\ln\left(x^2+1\right) + C, \quad C \in \mathbb{R}. \]
\emph{Vérification :} $\left(2\ln(x^2+1)\right)' = 2\times\dfrac{2x}{x^2+1} = \dfrac{4x}{x^2+1} = f(x)$. \checkmark

\textbf{2.} La fonction $g$ est continue sur $\mathbb{R}$. Posons $u(x) = x^2$ ; alors $u'(x) = 2x$. On fait apparaître $u'$ :
\[ g(x) = x\,e^{x^2} = \frac{1}{2}\times 2x\,e^{x^2} = \frac{1}{2}\,u'(x)\,e^{u(x)}. \]
La formule $\displaystyle\int u'e^u\,dx = e^u + C$ donne :
\[ \int g(x)\,dx = \frac{1}{2}e^{x^2} + C, \quad C \in \mathbb{R}. \]
\emph{Vérification :} $\left(\dfrac12 e^{x^2}\right)' = \dfrac12 \times 2x\,e^{x^2} = x e^{x^2} = g(x)$. \checkmark

\textbf{Conclusion.} Les primitives de $f$ sont $x \mapsto 2\ln(x^2+1) + C$ et celles de $g$ sont $x \mapsto \dfrac12 e^{x^2} + C$, avec $C \in \mathbb{R}$.
\end{exoresolu}

\begin{methode}[Arbre de décision pour chercher une primitive]
Face à une fonction $f$ dont on cherche une primitive, procéder dans l'ordre :
\begin{enumerate}
\item \textbf{Est-ce une fonction usuelle} (polynôme, $x^r$, $\cos(ax+b)$, $\sin(ax+b)$, $e^{ax+b}$) ? $\to$ appliquer directement le tableau des primitives usuelles.
\item \textbf{Est-ce une somme} ? $\to$ primitiver terme à terme (une primitive d'une somme est la somme des primitives).
\item \textbf{Y a-t-il une « fonction intérieure » $u$ dont la dérivée apparaît en facteur} (à une constante près) ? $\to$ identifier la forme :
\begin{center}
\begin{tabularx}{\linewidth}{|l|X|}
\hline
\rowcolor{popblueL} \textbf{Forme reconnue} & \textbf{Primitive} \\
\hline
$u'\,u^r$ ($r \neq -1$) & $\dfrac{u^{r+1}}{r+1}$ \\
\hline
$\dfrac{u'}{u}$ & $\ln|u|$ \\
\hline
$u'\,e^u$ & $e^u$ \\
\hline
$u'\cos(u)$ ou $u'\sin(u)$ & $\sin(u)$ ou $-\cos(u)$ \\
\hline
\end{tabularx}
\end{center}
\item \textbf{Cas $\cos^n x$, $\sin^n x$} : puissance impaire $\to$ forme $u'u^r$ ; puissance paire $\to$ linéarisation.
\item \textbf{Toujours} : préciser l'intervalle, ajouter $+C$, et vérifier en dérivant.
\end{enumerate}
\end{methode}

\begin{exoresolu}[Exercice de synthèse type Bac]
\textbf{Énoncé.} On considère la fonction $f$ définie sur $\mathbb{R}$ par :
\[ f(x) = (2x+1)\,e^{x^2+x} + \frac{1}{(x+2)^2} \quad \text{restreinte à } I = \left]-2\,;+\infty\right[. \]
\begin{enumerate}
\item Justifier que $f$ admet des primitives sur $I$.
\item Déterminer l'ensemble des primitives de $f$ sur $I$.
\item Déterminer la primitive $F$ de $f$ sur $I$ telle que $F(0) = 2$.
\end{enumerate}
\tcblower
\textbf{Solution détaillée.}

\textbf{1.} Sur $I = ]-2\,;+\infty[$ :
\begin{itemize}
\item la fonction $x \mapsto (2x+1)e^{x^2+x}$ est continue sur $\mathbb{R}$ (produit et composée de fonctions continues), donc sur $I$ ;
\item la fonction $x \mapsto \dfrac{1}{(x+2)^2}$ est continue sur $I$ car $x+2 \neq 0$ sur $I$.
\end{itemize}
La somme $f$ est donc continue sur $I$ : elle y admet des primitives.

\textbf{2.} On primitive chaque terme séparément.

\emph{Premier terme :} $(2x+1)e^{x^2+x}$. Posons $u(x) = x^2+x$ ; alors $u'(x) = 2x+1$, qui est exactement le facteur devant l'exponentielle. C'est la forme $u'e^u$ :
\[ \int (2x+1)e^{x^2+x}\,dx = e^{x^2+x}. \]

\emph{Second terme :} $\dfrac{1}{(x+2)^2} = (x+2)^{-2}$. C'est la forme $u'u^{-2}$ avec $u(x) = x+2$, $u'(x) = 1$ (déjà présent) :
\[ \int \frac{dx}{(x+2)^2} = \frac{(x+2)^{-1}}{-1} = -\frac{1}{x+2}. \]

L'ensemble des primitives de $f$ sur $I$ est donc :
\[ F(x) = e^{x^2+x} - \frac{1}{x+2} + C, \quad C \in \mathbb{R}. \]
\emph{Vérification :} $F'(x) = (2x+1)e^{x^2+x} + \dfrac{1}{(x+2)^2} = f(x)$. \checkmark

\textbf{3.} On impose $F(0) = 2$ :
\[ F(0) = e^{0} - \frac{1}{0+2} + C = 1 - \frac{1}{2} + C = \frac{1}{2} + C = 2 \iff C = \frac{3}{2}. \]

\textbf{Conclusion.} La primitive de $f$ sur $I$ vérifiant $F(0) = 2$ est :
\[ \boxed{F(x) = e^{x^2+x} - \frac{1}{x+2} + \frac{3}{2}.} \]
\end{exoresolu}

\begin{attention}[Les 5 erreurs qui coûtent des points au Bac]
\begin{enumerate}
\item \textbf{Oublier le $+C$.} « Déterminer l'ensemble des primitives » exige $F_0(x) + C$, $C \in \mathbb{R}$. Sans le $+C$, tu ne donnes qu'\emph{une} primitive : la réponse est incomplète. À l'inverse, quand une condition initiale est donnée, la réponse finale ne doit \emph{plus} contenir de $C$.
\item \textbf{Oublier de diviser par $a$ dans $f(ax+b)$.} Une primitive de $\cos(3x)$ est $\dfrac{1}{3}\sin(3x)$ et non $\sin(3x)$ ; une primitive de $\dfrac{1}{(2x+1)^2}$ est $-\dfrac{1}{2(2x+1)}$ et non $-\dfrac{1}{2x+1}$. Réflexe : \emph{toujours vérifier en dérivant}.
\item \textbf{Mal ajuster la constante dans $u'u^r$.} Pour $x(x^2+1)^3$, on a $u' = 2x$, donc il faut écrire $x(x^2+1)^3 = \dfrac12 \times 2x(x^2+1)^3$ avant d'appliquer la formule. Ne jamais « inventer » un facteur sans compenser.
\item \textbf{Appliquer $\int u'u^r = \frac{u^{r+1}}{r+1}$ avec $r = -1$.} Le cas $r=-1$ est la forme $\dfrac{u'}{u}$, de primitive $\ln|u|$, et non une puissance. De même, $\displaystyle\int \frac{dx}{x^2}$ se traite comme $x^{-2}$ (primitive $-\frac1x$), pas avec un logarithme.
\item \textbf{Négliger l'intervalle de validité.} $\sqrt{x}$ et $\dfrac{1}{x^r}$ n'ont de primitives que sur un intervalle où elles sont définies et continues ; $\dfrac{1}{(cx+d)^r}$ impose $cx+d \neq 0$. Annoncer des primitives « sur $\mathbb{R}$ » pour $\dfrac{1}{x}$ est une erreur de rigueur sanctionnée.
\end{enumerate}
\end{attention}

\newpage

\coursec{Exercices}

\coursec{Primitives d'une fonction continue}

\courssub{Notion de primitive et existence}

\begin{definition}[Primitive d'une fonction]
Soit $I$ un intervalle de $\mathbb{R}$ et $f$ une fonction définie sur $I$. On appelle \textbf{primitive de $f$ sur $I$} toute fonction $F$ dérivable sur $I$ telle que $F'(x) = f(x)$ pour tout $x \in I$.
\end{definition}

\begin{propriete}[Existence des primitives]
\emph{Théorème fondamental :} Toute fonction continue sur un intervalle $I$ admet au moins une primitive sur $I$.
\end{propriete}

\begin{attention}[Condition d'intervalle]
L'hypothèse « $I$ intervalle » est \textbf{essentielle}. Par exemple, $f(x) = \frac{1}{x}$ est continue sur $\mathbb{R}^* = ]-\infty,0[ \cup ]0,+\infty[$ (qui n'est \textbf{pas} un intervalle), mais elle n'admet pas de primitive \emph{unique} sur $\mathbb{R}^*$ au sens global (voir exercice Vrai/Faux n°2).
\end{attention}

\begin{exemplebox}[Exemples de primitives]
\begin{itemize}
\item $F(x) = \frac{x^3}{3}$ est une primitive de $f(x)=x^2$ sur $\mathbb{R}$.
\item $F(x) = -\frac{1}{x}$ est une primitive de $f(x)=\frac{1}{x^2}$ sur $]0,+\infty[$.
\item $F(x) = \sin x$ est une primitive de $f(x)=\cos x$ sur $\mathbb{R}$.
\end{itemize}
\end{exemplebox}

\courssub{Ensemble des primitives ; condition initiale}

\begin{propriete}[Ensemble des primitives]
Si $F$ est une primitive de $f$ sur un intervalle $I$, alors l'ensemble des primitives de $f$ sur $I$ est exactement $\{ F + C \mid C \in \mathbb{R} \}$. On note souvent $\int f(x) \, dx = F(x) + C$.
\end{propriete}

\begin{propriete}[Primitive avec condition initiale]
Soit $f$ continue sur un intervalle $I$, $a \in I$ et $y_0 \in \mathbb{R}$. Il existe \textbf{une et une seule} primitive $F$ de $f$ sur $I$ telle que $F(a) = y_0$. Elle s'obtient en fixant la constante $C$ dans $F(x) = F_0(x) + C$ où $F_0$ est une primitive quelconque.
\end{propriete}

\begin{methode}[Déterminer la primitive avec condition initiale]
\begin{enumerate}
\item Calculer une primitive $F_0$ de $f$ (sans constante).
\item Écrire $F(x) = F_0(x) + C$.
\item Utiliser la condition $F(a) = y_0$ pour trouver $C$.
\item Écrire l'expression finale de $F$.
\end{enumerate}
\end{methode}

\courssub{Primitives des fonctions usuelles}

\begin{aretenir}[Tableau des primitives usuelles (constante $C \in \mathbb{R}$)]
\begin{center}
\begin{tabularx}{\linewidth}{|X|X|X|}
\hline
\rowcolor{popblueL}
\textbf{Fonction $f(x)$} & \textbf{Une primitive $F(x)$} & \textbf{Intervalle de validité} \\
\hline
$x^r$ ($r \in \mathbb{Q}$, $r \neq -1$) & $\dfrac{x^{r+1}}{r+1}$ & $\mathbb{R}$ si $r \in \mathbb{N}$ ; $]0,+\infty[$ si $r \notin \mathbb{Z}$ \\
\hline
$\dfrac{1}{x}$ & $\ln x$ & $]0,+\infty[$ \\
\hline
$\cos(ax+b)$ ($a \neq 0$) & $\dfrac{1}{a}\sin(ax+b)$ & $\mathbb{R}$ \\
\hline
$\sin(ax+b)$ ($a \neq 0$) & $-\dfrac{1}{a}\cos(ax+b)$ & $\mathbb{R}$ \\
\hline
$e^{ax+b}$ ($a \neq 0$) & $\dfrac{1}{a}e^{ax+b}$ & $\mathbb{R}$ \\
\hline
\end{tabularx}
\end{center}
\end{aretenir}

\courssub{Primitives des formes composées : $u' \cdot u^r$ et $\frac{a}{(cx+d)^r}$}

\begin{propriete}[Forme $u' \cdot u^r$ ($r \in \mathbb{Q}$, $r \neq -1$)]
Soit $u$ une fonction dérivable et strictement positive (ou $r \in \mathbb{Z}$) sur un intervalle $I$. Une primitive de $u'(x) \cdot [u(x)]^r$ sur $I$ est $\dfrac{[u(x)]^{r+1}}{r+1}$.
\end{propriete}

\begin{propriete}[Forme $\frac{a}{(cx+d)^r}$ ($r \in \mathbb{Q}$, $r \neq -1$, $c \neq 0$)]
Une primitive sur tout intervalle ne contenant pas $-\frac{d}{c}$ est $\dfrac{a}{c(1-r)}(cx+d)^{1-r}$.
\emph{Preuve :} Poser $u(x)=cx+d$, $u'(x)=c$, alors $\frac{a}{(cx+d)^r} = \frac{a}{c} u'(x) [u(x)]^{-r}$.
\end{propriete}

\begin{methode}[Reconnaître et traiter la forme $u' \cdot u^r$]
\begin{enumerate}
\item Identifier une fonction $u(x)$ « intérieure » (souvent ce qui est entre parenthèses, au dénominateur, ou sous racine).
\item Calculer $u'(x)$.
\item Vérifier si l'expression de $f(x)$ est un multiple constant de $u'(x) \cdot [u(x)]^r$.
\item Si oui, la primitive est $\frac{k}{r+1} [u(x)]^{r+1}$ où $k$ est le coefficient ajusté.
\end{enumerate}
\end{methode}

\begin{exemplebox}[Application : $f(x) = \frac{x^2}{\sqrt{x^3+2}}$]
$u(x) = x^3+2 \Rightarrow u'(x) = 3x^2$. On a $f(x) = \frac{1}{3} \cdot 3x^2 \cdot (x^3+2)^{-1/2} = \frac{1}{3} u'(x) [u(x)]^{-1/2}$.
Primitive : $\frac{1}{3} \cdot \frac{[u(x)]^{1/2}}{1/2} = \frac{2}{3}\sqrt{x^3+2} + C$. Valide sur $]-\sqrt[3]{2}, +\infty[$.
\end{exemplebox}

\courssub{Primitives de $\cos^n x$ et $\sin^n x$}

\begin{methode}[Stratégies pour $\cos^n x$ et $\sin^n x$]
\begin{enumerate}
\item \textbf{Pair ($n=2k$) :} Linéariser avec $\cos^2 \theta = \frac{1+\cos 2\theta}{2}$, $\sin^2 \theta = \frac{1-\cos 2\theta}{2}$.
\item \textbf{Impair ($n=2k+1$) :} Factoriser $\cos x$ ou $\sin x$ et utiliser $\cos^2 = 1-\sin^2$ ou $\sin^2 = 1-\cos^2$ pour obtenir une forme $u' \cdot u^r$ avec $u = \sin x$ ou $u = \cos x$.
\end{enumerate}
\end{methode}

\begin{exemplebox}[$\sin^3 x$]
$\sin^3 x = \sin x (1-\cos^2 x) = -\cos' x + \cos' x \cdot \cos^2 x$.
Primitive : $-\cos x + \frac{\cos^3 x}{3} + C$.
\end{exemplebox}

\courssub{Compléments Bac : $u'/u$, $u'e^u$ et synthèse}

\begin{propriete}[Formes $u'/u$ et $u'e^u$]
Soit $u$ dérivable sur un intervalle $I$.
\begin{itemize}
\item Si $u(x) > 0$ sur $I$, une primitive de $\frac{u'(x)}{u(x)}$ est $\ln(u(x))$.
\item Une primitive de $u'(x) e^{u(x)}$ est $e^{u(x)}$.
\end{itemize}
\end{propriete}

\begin{aretenir}[Tableau récapitulatif des formes composées (Bac)]
\begin{center}
\begin{tabularx}{\linewidth}{|l|X|X|}
\hline
\rowcolor{popblueL}
\textbf{Forme} & \textbf{Condition} & \textbf{Primitive} \\
\hline
$u'(x) [u(x)]^r$ & $r \neq -1$ & $\dfrac{[u(x)]^{r+1}}{r+1}$ \\
\hline
$\dfrac{u'(x)}{u(x)}$ & $u(x) > 0$ & $\ln(u(x))$ \\
\hline
$u'(x) e^{u(x)}$ & -- & $e^{u(x)}$ \\
\hline
\end{tabularx}
\end{center}
\end{aretenir}

\begin{savaistu}[Histoire : Le symbole $\int$]
Le symbole d'intégrale $\int$ (un « S » allongé) a été introduit par Leibniz en 1675 pour « Summa » (somme). Il voyait l'intégrale comme une somme infinie de quantités infinitésimales. Newton, de son côté, utilisait la notation $\dot{x}$ pour les dérivées (fluxions). La notation $F(x)+C$ pour les primitives rappelle que la dérivation « efface » les constantes : c'est le cœur du \emph{Théorème fondamental de l'analyse}.
\end{savaistu}

% ============================================================
% EXERCICES
% ============================================================

\courssub{Je vérifie mes connaissances}

\begin{exercice}[QCM : une seule bonne réponse]
Pour chaque question, entoure la bonne réponse. Aucune justification n'est demandée.
\begin{enumerate}
\item Une primitive de $f(x) = \dfrac{1}{x^2}$ sur $]0\,;+\infty[$ est :
\begin{enumerate}
\item[a)] $F(x) = \ln x$ \qquad b) $F(x) = -\dfrac{1}{x}$ \qquad c) $F(x) = \dfrac{2}{x^3}$ \qquad d) $F(x) = -\dfrac{2}{x^3}$
\end{enumerate}
\item L'ensemble des primitives de $f(x) = \cos(2x)$ sur $\mathbb{R}$ est :
\begin{enumerate}
\item[a)] $F(x) = -\sin(2x) + C$ \qquad b) $F(x) = 2\sin(2x) + C$ \qquad c) $F(x) = \dfrac{1}{2}\sin(2x) + C$ \qquad d) $F(x) = -\dfrac{1}{2}\sin(2x) + C$
\end{enumerate}
\item La fonction $F(x) = (x^2 + 1)^5$ est une primitive sur $\mathbb{R}$ de :
\begin{enumerate}
\item[a)] $f(x) = 5(x^2+1)^4$ \qquad b) $f(x) = 10x(x^2+1)^4$ \qquad c) $f(x) = 2x(x^2+1)^5$ \qquad d) $f(x) = 10(x^2+1)^4$
\end{enumerate}
\item La primitive de $f(x) = 3x^2 - 4x + 1$ sur $\mathbb{R}$ qui s'annule en $x = 1$ est :
\begin{enumerate}
\item[a)] $F(x) = x^3 - 2x^2 + x$ \qquad b) $F(x) = x^3 - 2x^2 + x + 1$ \qquad c) $F(x) = 6x - 4$ \qquad d) $F(x) = x^3 - 2x^2 + x - 2$
\end{enumerate}
\end{enumerate}
\end{exercice}

\begin{exercice}[Vrai ou faux, en justifiant]
Pour chacune des affirmations suivantes, dire si elle est vraie ou fausse et \textbf{justifier rigoureusement} la réponse.
\begin{enumerate}
\item « $F(x) = (x^2+1)^5$ est une primitive de $f(x) = 10x(x^2+1)^4$ sur $\mathbb{R}$. »
\item « Deux primitives quelconques de la fonction $x \mapsto \dfrac{1}{x}$ sur $\mathbb{R}^*$ diffèrent toujours d'une constante. »
\item « Si $F$ est une primitive de $f$ sur un intervalle $I$, alors $F + 2024$ est aussi une primitive de $f$ sur $I$. »
\item « La fonction $x \mapsto \sqrt{x}$ admet des primitives sur $\mathbb{R}$. »
\item « Une fonction qui admet une primitive sur un intervalle $I$ possède exactement une primitive sur $I$. »
\end{enumerate}
\end{exercice}

\begin{exercice}[Définitions et cours]
\begin{enumerate}
\item Donner la définition d'une primitive d'une fonction $f$ sur un intervalle $I$.
\item Énoncer le théorème qui garantit l'existence de primitives pour une fonction continue sur un intervalle.
\item Soit $f$ une fonction continue sur un intervalle $I$, $a \in I$ et $y_0 \in \mathbb{R}$. Que peut-on dire de la primitive de $f$ prenant la valeur $y_0$ en $a$ ? (Existence ? Unicité ?)
\item Compléter le tableau des primitives usuelles suivant (constante $C \in \mathbb{R}$) :
\begin{center}
\begin{tabularx}{\linewidth}{|X|X|X|}
\hline
\rowcolor{popblueL}
\textbf{Fonction $f(x)$} & \textbf{Une primitive $F(x)$} & \textbf{Intervalle} \\
\hline
$x^r$ ($r \in \mathbb{Q}$, $r \neq -1$) & & \\
\hline
$\cos(ax+b)$ ($a \neq 0$) & & \\
\hline
$\sin(ax+b)$ ($a \neq 0$) & & \\
\hline
$u'(x)\,[u(x)]^r$ ($r \neq -1$) & & \\
\hline
\end{tabularx}
\end{center}
\end{enumerate}
\end{exercice}

\begin{exercice}[Calculs directs]
Déterminer l'ensemble des primitives de chacune des fonctions suivantes sur l'intervalle indiqué.
\begin{enumerate}
\item $f(x) = 5x^4 - 3x^2 + 7$ sur $\mathbb{R}$ ;
\item $g(x) = \dfrac{2}{x^3}$ sur $]0\,;+\infty[$ ;
\item $h(x) = 4\cos\left(2x - \dfrac{\pi}{3}\right)$ sur $\mathbb{R}$ ;
\item $k(x) = -6\sin(3x) + \dfrac{1}{2\sqrt{x}}$ sur $]0\,;+\infty[$.
\end{enumerate}
\end{exercice}

\courssub{Je m'entraîne}

\begin{exercice}[Reconnaître les formes]
En reconnaissant à chaque fois une forme du type $u' \cdot u^r$ ou $\dfrac{a}{(cx+d)^r}$, déterminer l'ensemble des primitives des fonctions suivantes sur un intervalle à préciser.
\begin{enumerate}
\item[a)] $f(x) = (4x-1)(2x^2 - x + 3)^3$ ;
\item[b)] $g(x) = \dfrac{x^2}{\sqrt{x^3 + 2}}$ ;
\item[c)] $h(x) = \dfrac{3}{(5x-2)^4}$ ;
\item[d)] $k(x) = \cos x \cdot \sin^4 x$.
\end{enumerate}
\emph{Indication : pour chaque fonction, identifier d'abord $u$, calculer $u'$, puis ajuster le coefficient.}
\end{exercice}

\begin{exercice}[Primitives avec condition initiale]
\begin{enumerate}
\item Déterminer la primitive $F$ de $f(x) = 6x^2 - \dfrac{4}{x^2}$ sur $]0\,;+\infty[$ vérifiant $F(1) = 0$.
\item Déterminer la primitive $G$ de $g(x) = \sin\left(3x + \dfrac{\pi}{6}\right)$ sur $\mathbb{R}$ vérifiant $G(0) = \dfrac{1}{2}$.
\item Déterminer la primitive $H$ de $h(x) = \dfrac{2x}{(x^2+1)^2}$ sur $\mathbb{R}$ vérifiant $H(0) = 3$.
\end{enumerate}
\end{exercice}

\begin{exercice}[Retrouver une fonction à partir de sa primitive]
Une fonction $f$, définie sur $\mathbb{R}$, admet pour primitive la fonction
\[ F(x) = \dfrac{x\,e^{2x}}{2} - \dfrac{e^{2x}}{4}. \]
\begin{enumerate}
\item En dérivant $F$, retrouver l'expression de $f(x)$.
\item Vérifier que $f$ est de la forme $f(x) = x\,e^{2x}$.
\item Déterminer alors l'ensemble des primitives de $f$ sur $\mathbb{R}$.
\item En déduire la primitive de $f$ qui s'annule en $0$.
\end{enumerate}
\end{exercice}

\begin{exercice}[Forme $u'/u$ — complément Bac]
On considère la fonction $f$ définie sur $\mathbb{R}$ par
\[ f(x) = \dfrac{6x + 2}{3x^2 + 2x + 1}. \]
\begin{enumerate}
\item Justifier que $f$ est définie et continue sur $\mathbb{R}$ (on étudiera le signe du discriminant du dénominateur).
\item Montrer que $f$ est de la forme $\dfrac{u'}{u}$ avec une fonction $u$ strictement positive sur $\mathbb{R}$ que l'on précisera.
\item En déduire l'ensemble des primitives de $f$ sur $\mathbb{R}$.
\item Déterminer la primitive $F$ de $f$ vérifiant $F(0) = 2$.
\end{enumerate}
\end{exercice}

\courssub{Je me prépare au Bac}

\begin{exercice}[Type Bac — Primitives et condition initiale]
On considère la fonction $f$ définie sur $]0\,;+\infty[$ par
\[ f(x) = 4x^3 - \dfrac{6}{x^2} + 2\sin\left(3x - \dfrac{\pi}{4}\right). \]
\begin{enumerate}
\item Justifier que $f$ admet des primitives sur $]0\,;+\infty[$.
\item \begin{enumerate}
\item Donner une primitive de $x \mapsto 4x^3$ sur $]0\,;+\infty[$.
\item Donner une primitive de $x \mapsto -\dfrac{6}{x^2}$ sur $]0\,;+\infty[$.
\item Montrer qu'une primitive de $x \mapsto 2\sin\left(3x - \dfrac{\pi}{4}\right)$ sur $\mathbb{R}$ est $x \mapsto -\dfrac{2}{3}\cos\left(3x - \dfrac{\pi}{4}\right)$.
\end{enumerate}
\item En déduire l'ensemble des primitives de $f$ sur $]0\,;+\infty[$.
\item Déterminer la primitive $F$ de $f$ sur $]0\,;+\infty[$ qui prend la valeur $5$ en $x = 1$. On donnera l'expression exacte de $F(x)$, puis une valeur approchée de la constante à $10^{-3}$ près, sachant que $\cos\left(3 - \dfrac{\pi}{4}\right) \approx -0,660$.
\end{enumerate}
\end{exercice}

\begin{exercice}[Type Bac — Linéarisation et primitives trigonométriques]
Soit $f$ la fonction définie sur $\mathbb{R}$ par
\[ f(x) = \cos^2(2x) + \sin^3 x. \]
\begin{enumerate}
\item \begin{enumerate}
\item Rappeler la formule donnant $\cos(2a)$ en fonction de $\cos^2 a$.
\item En déduire que pour tout réel $x$ : $\cos^2(2x) = \dfrac{1 + \cos(4x)}{2}$.
\item Déterminer alors l'ensemble des primitives de $x \mapsto \cos^2(2x)$ sur $\mathbb{R}$.
\end{enumerate}
\item \begin{enumerate}
\item Montrer que pour tout réel $x$ : $\sin^3 x = \sin x \left(1 - \cos^2 x\right)$.
\item En posant $u(x) = \cos x$, reconnaître dans l'expression précédente une combinaison de formes $u' \cdot u^r$.
\item En déduire l'ensemble des primitives de $x \mapsto \sin^3 x$ sur $\mathbb{R}$.
\end{enumerate}
\item Conclure : donner l'ensemble des primitives de $f$ sur $\mathbb{R}$.
\item Déterminer la primitive $F$ de $f$ sur $\mathbb{R}$ vérifiant $F\left(\dfrac{\pi}{2}\right) = \dfrac{\pi}{4}$.
\end{enumerate}
\end{exercice}

\begin{exercice}[Problème de synthèse — Le moto-taxi de Buea]
Un moto-taxi (bendskin) circule sur la route de Molyko à Buea. Sa vitesse instantanée, en mètres par seconde, est modélisée pendant les cinq premières secondes du mouvement par la fonction
\[ v(t) = 3t^2 - 2t + 1, \qquad t \in [0\,;5], \]
où $t$ est le temps écoulé en secondes depuis le départ. On note $x(t)$ la distance parcourue (en mètres) à l'instant $t$ : la fonction $x$ est donc une primitive de $v$ sur $[0\,;5]$. Au départ ($t = 0$), le moto-taxi a déjà parcouru $10$ m depuis son point de stationnement.
\begin{enumerate}
\item Justifier que $v$ admet des primitives sur $[0\,;5]$.
\item Déterminer l'ensemble des primitives de $v$ sur $[0\,;5]$.
\item En utilisant la condition $x(0) = 10$, montrer que $x(t) = t^3 - t^2 + t + 10$.
\item Calculer la distance totale parcourue par le moto-taxi au bout de $5$ secondes.
\item Le tarif d'une course est estimé à $25$ FCFA par mètre parcouru (tarif fictif de l'exercice). Quel serait, selon ce modèle, le montant de la course pour ces cinq premières secondes ?
\item Un second moto-taxi part du même point avec une vitesse $w(t) = 2t + 4$ (en m/s) et une distance initiale nulle. Déterminer la distance $y(t)$ qu'il a parcourue à l'instant $t$, puis résoudre l'équation $x(t) = y(t)$ sur $[0\,;5]$. Interpréter le résultat dans le contexte de la course.
\end{enumerate}
\end{exercice}



\newpage

\coursec{Corrigés des exercices}

\begin{corrige}[Exercice 1 -- QCM]
\begin{enumerate}
\item \textbf{Réponse b).} On a $f(x) = \dfrac{1}{x^2} = x^{-2}$. En appliquant la formule des puissances avec $r = -2 \neq -1$ :
\[ F(x) = \frac{x^{-2+1}}{-2+1} = \frac{x^{-1}}{-1} = -\frac{1}{x}. \]
Vérification : $F'(x) = \dfrac{1}{x^2} = f(x)$. \checkmark
\item \textbf{Réponse c).} Une primitive de $\cos(ax+b)$ est $\dfrac{1}{a}\sin(ax+b)$. Ici $a = 2$, donc $F(x) = \dfrac{1}{2}\sin(2x) + C$.
\item \textbf{Réponse b).} On dérive $F(x) = (x^2+1)^5$ avec la formule $(u^n)' = n\,u'\,u^{n-1}$ :
\[ F'(x) = 5 \times 2x \times (x^2+1)^4 = 10x(x^2+1)^4. \]
\item \textbf{Réponse a).} Une primitive de $3x^2 - 4x + 1$ est $F_0(x) = x^3 - 2x^2 + x$. Or $F_0(1) = 1 - 2 + 1 = 0$ : la condition est déjà satisfaite avec $C = 0$.
\end{enumerate}
\end{corrige}

\begin{corrige}[Exercice 2 -- Vrai ou Faux]
\begin{enumerate}
\item \textbf{Vrai.} $F$ est dérivable sur $\mathbb{R}$ comme composée de fonctions dérivables, et par la formule $(u^n)' = n\,u'\,u^{n-1}$ :
\[ F'(x) = 5(x^2+1)^4 \times 2x = 10x(x^2+1)^4 = f(x) \quad \text{pour tout } x \in \mathbb{R}. \]
\item \textbf{Faux.} Le point crucial est que $\mathbb{R}^* = ]-\infty,0[ \cup ]0,+\infty[$ \textbf{n'est pas un intervalle}. Le théorème « deux primitives d'une même fonction diffèrent d'une constante » n'est valable que sur un intervalle. Ici, les primitives de $x \mapsto \frac{1}{x}$ sont de la forme $\ln|x| + C_1$ sur $]-\infty,0[$ et $\ln|x| + C_2$ sur $]0,+\infty[$, avec $C_1$ et $C_2$ \textbf{indépendantes}. Contre-exemple : $F_1(x) = \ln|x|$ et $F_2(x) = \ln|x| + \begin{cases} 0 & \text{si } x < 0 \\ 5 & \text{si } x > 0 \end{cases}$ sont deux primitives dont la différence n'est pas constante.
\item \textbf{Vrai.} $(F + 2024)' = F' + 0 = f$ sur $I$, donc $F + 2024$ est bien une primitive de $f$.
\item \textbf{Faux.} La fonction $x \mapsto \sqrt{x}$ n'est définie que sur $[0,+\infty[$ : elle ne peut donc pas admettre de primitive \emph{sur $\mathbb{R}$}. En revanche, elle est continue sur $[0,+\infty[$ et y admet des primitives, par exemple $F(x) = \dfrac{2}{3}x^{3/2}$.
\item \textbf{Faux.} C'est le contraire : si $F$ est une primitive de $f$ sur $I$, alors toutes les fonctions $F + C$ ($C \in \mathbb{R}$) sont aussi des primitives de $f$ : il y en a une \textbf{infinité}. L'unicité n'apparaît que lorsqu'on impose une condition initiale $F(a) = y_0$.
\end{enumerate}
\end{corrige}

\begin{corrige}[Exercice 3 -- Définitions et cours]
\begin{enumerate}
\item Soit $I$ un intervalle de $\mathbb{R}$ et $f$ une fonction définie sur $I$. On appelle \textbf{primitive de $f$ sur $I$} toute fonction $F$ \textbf{dérivable sur $I$} telle que $F'(x) = f(x)$ pour tout $x \in I$.
\item \textbf{Théorème d'existence :} toute fonction \textbf{continue} sur un intervalle $I$ admet au moins une primitive sur $I$.
\item Il existe \textbf{une et une seule} primitive $F$ de $f$ sur $I$ vérifiant $F(a) = y_0$. (Existence et unicité.)
\item Tableau complété :
\begin{center}
\begin{tabularx}{\linewidth}{|X|X|X|}
\hline
\rowcolor{popblueL}
\textbf{Fonction $f(x)$} & \textbf{Une primitive $F(x)$} & \textbf{Intervalle(s) de validité} \\
\hline
$x^r$ ($r \in \mathbb{Q}$, $r \neq -1$) & $\dfrac{x^{r+1}}{r+1} + C$ & $\mathbb{R}$ si $r \in \mathbb{N}$ ; $]0,+\infty[$ (et $]-\infty,0[$) si $r \in \mathbb{Q} \setminus \mathbb{N}$ \\
\hline
$\cos(ax+b)$ ($a \neq 0$) & $\dfrac{1}{a}\sin(ax+b) + C$ & $\mathbb{R}$ \\
\hline
$\sin(ax+b)$ ($a \neq 0$) & $-\dfrac{1}{a}\cos(ax+b) + C$ & $\mathbb{R}$ \\
\hline
$u'(x)\,[u(x)]^r$ ($r \neq -1$) & $\dfrac{[u(x)]^{r+1}}{r+1} + C$ & Intervalle où $u$ est dérivable, avec $u > 0$ (ou $r \in \mathbb{Z}$) \\
\hline
\end{tabularx}
\end{center}
\end{enumerate}
\end{corrige}

\begin{corrige}[Exercice 4 -- Calculs directs]
\begin{enumerate}
\item On primitive terme à terme : $5x^4 \to x^5$, $-3x^2 \to -x^3$, $7 \to 7x$.
\[ \boxed{F(x) = x^5 - x^3 + 7x + C, \quad C \in \mathbb{R}} \]
Vérification : $F'(x) = 5x^4 - 3x^2 + 7 = f(x)$. \checkmark
\item $g(x) = 2x^{-3}$, forme $x^r$ avec $r = -3 \neq -1$ :
\[ G(x) = 2 \times \frac{x^{-3+1}}{-3+1} + C = 2 \times \frac{x^{-2}}{-2} + C = \boxed{-\frac{1}{x^2} + C, \quad C \in \mathbb{R}} \]
Vérification : $G'(x) = \dfrac{2}{x^3} = g(x)$. \checkmark
\item Forme $A\cos(ax+b)$ avec $A = 4$, $a = 2$ : une primitive est $\dfrac{A}{a}\sin(ax+b)$.
\[ \boxed{H(x) = 2\sin\left(2x - \frac{\pi}{3}\right) + C, \quad C \in \mathbb{R}} \]
Vérification : $H'(x) = 2 \times 2\cos(2x - \pi/3) = 4\cos(2x-\pi/3) = h(x)$. \checkmark
\item On traite séparément les deux termes sur $]0,+\infty[$ :
\begin{itemize}
\item primitive de $-6\sin(3x)$ : $-6 \times \left(-\dfrac{1}{3}\cos(3x)\right) = 2\cos(3x)$ ;
\item primitive de $\dfrac{1}{2\sqrt{x}} = \dfrac{1}{2}x^{-1/2}$ : $\dfrac{1}{2} \times \dfrac{x^{1/2}}{1/2} = \sqrt{x}$.
\end{itemize}
\[ \boxed{K(x) = 2\cos(3x) + \sqrt{x} + C, \quad C \in \mathbb{R}} \]
Vérification : $K'(x) = -6\sin(3x) + \dfrac{1}{2\sqrt{x}} = k(x)$. \checkmark
\end{enumerate}
\end{corrige}

\begin{corrige}[Exercice 5 -- Reconnaître les formes]
\begin{enumerate}
\item[a)] Posons $u(x) = 2x^2 - x + 3$ ; alors $u'(x) = 4x - 1$. On reconnaît exactement :
\[ f(x) = u'(x)\,[u(x)]^3 \quad (r = 3 \neq -1). \]
Le discriminant de $u$ est $\Delta = (-1)^2 - 4 \times 2 \times 3 = 1 - 24 = -23 < 0$, donc $u(x) > 0$ sur $\mathbb{R}$ : la forme est valide sur $\mathbb{R}$ tout entier.
\[ \boxed{F(x) = \frac{(2x^2 - x + 3)^4}{4} + C, \quad C \in \mathbb{R}, \text{ sur } \mathbb{R}} \]
\item[b)] Posons $u(x) = x^3 + 2$ ; alors $u'(x) = 3x^2$, donc $x^2 = \dfrac{1}{3}u'(x)$ et
\[ g(x) = \frac{1}{3}\,u'(x)\,[u(x)]^{-1/2} \quad \left(r = -\frac{1}{2} \neq -1\right). \]
Condition de validité : $u(x) > 0 \iff x^3 > -2 \iff x > -\sqrt[3]{2}$.
\[ G(x) = \frac{1}{3} \times \frac{[u(x)]^{1/2}}{1/2} + C = \boxed{\frac{2}{3}\sqrt{x^3 + 2} + C, \quad C \in \mathbb{R}, \text{ sur } \left]-\sqrt[3]{2}\,;+\infty\right[} \]
\item[c)] Forme $\dfrac{a}{(cx+d)^r}$ avec $a = 3$, $c = 5$, $d = -2$, $r = 4$. Autre méthode : $u(x) = 5x - 2$, $u'(x) = 5$, donc $h(x) = \dfrac{3}{5}u'(x)[u(x)]^{-4}$.
\[ H(x) = \frac{3}{5} \times \frac{[u(x)]^{-3}}{-3} + C = \boxed{-\frac{1}{5(5x-2)^3} + C, \quad C \in \mathbb{R}} \]
valide sur tout intervalle ne contenant pas $\dfrac{2}{5}$, c'est-à-dire sur $\left]-\infty\,;\dfrac{2}{5}\right[$ ou sur $\left]\dfrac{2}{5}\,;+\infty\right[$. \textbf{Vigilance :} la constante $C$ peut être différente sur chacun de ces deux intervalles.
\item[d)] Posons $u(x) = \sin x$ ; alors $u'(x) = \cos x$ et
\[ k(x) = u'(x)\,[u(x)]^4 \quad (r = 4 \in \mathbb{Z}, \text{ donc valide sur } \mathbb{R}). \]
\[ \boxed{K(x) = \frac{\sin^5 x}{5} + C, \quad C \in \mathbb{R}, \text{ sur } \mathbb{R}} \]
Vérification : $K'(x) = \dfrac{5\sin^4 x \cos x}{5} = \cos x \sin^4 x = k(x)$. \checkmark
\end{enumerate}
\end{corrige}

\begin{corrige}[Exercice 6 -- Primitives avec condition initiale]
\begin{enumerate}
\item \textbf{Étape 1 :} $f(x) = 6x^2 - 4x^{-2}$, donc une primitive est $F_0(x) = 2x^3 - 4 \times \dfrac{x^{-1}}{-1} = 2x^3 + \dfrac{4}{x}$.
\textbf{Étape 2 :} $F(x) = 2x^3 + \dfrac{4}{x} + C$.
\textbf{Étape 3 :} $F(1) = 0 \iff 2 + 4 + C = 0 \iff C = -6$.
\[ \boxed{F(x) = 2x^3 + \frac{4}{x} - 6 \quad \text{sur } ]0\,;+\infty[} \]
Vérification : $F(1) = 2 + 4 - 6 = 0$. \checkmark
\item \textbf{Étape 1 :} une primitive de $\sin(3x + \pi/6)$ est $-\dfrac{1}{3}\cos(3x + \pi/6)$.
\textbf{Étape 2 :} $G(x) = -\dfrac{1}{3}\cos\left(3x + \dfrac{\pi}{6}\right) + C$.
\textbf{Étape 3 :} $G(0) = \dfrac{1}{2} \iff -\dfrac{1}{3}\cos\dfrac{\pi}{6} + C = \dfrac{1}{2} \iff -\dfrac{\sqrt{3}}{6} + C = \dfrac{1}{2} \iff C = \dfrac{3}{6} + \dfrac{\sqrt{3}}{6} = \dfrac{3 + \sqrt{3}}{6}$.
\[ \boxed{G(x) = -\frac{1}{3}\cos\left(3x + \frac{\pi}{6}\right) + \frac{3 + \sqrt{3}}{6} \quad \text{sur } \mathbb{R}} \]
\item \textbf{Étape 1 :} posons $u(x) = x^2 + 1$ ; $u'(x) = 2x$, donc $h(x) = u'(x)[u(x)]^{-2}$ ($r = -2 \neq -1$). Une primitive est $H_0(x) = \dfrac{[u(x)]^{-1}}{-1} = -\dfrac{1}{x^2+1}$.
\textbf{Étape 2 :} $H(x) = -\dfrac{1}{x^2+1} + C$.
\textbf{Étape 3 :} $H(0) = 3 \iff -1 + C = 3 \iff C = 4$.
\[ \boxed{H(x) = 4 - \frac{1}{x^2+1} \quad \text{sur } \mathbb{R}} \]
Vérification : $H(0) = 4 - 1 = 3$. \checkmark
\end{enumerate}
\end{corrige}

\begin{corrige}[Exercice 7 -- Retrouver une fonction à partir de sa primitive]
\begin{enumerate}
\item $F$ est dérivable sur $\mathbb{R}$ comme produit et somme de fonctions dérivables. On dérive le premier terme avec la formule du produit $(uv)' = u'v + uv'$, avec $u(x) = \dfrac{x}{2}$ et $v(x) = e^{2x}$ :
\begin{align*}
F'(x) &= \frac{1}{2}e^{2x} + \frac{x}{2} \times 2e^{2x} - \frac{2e^{2x}}{4} \\
&= \frac{e^{2x}}{2} + x\,e^{2x} - \frac{e^{2x}}{2} \\
&= x\,e^{2x}.
\end{align*}
\item On obtient exactement $f(x) = x\,e^{2x}$ : la fonction $f$ est bien retrouvée par dérivation de sa primitive, ce qui illustre la définition même d'une primitive.
\item Puisque $F$ est une primitive de $f$ sur l'intervalle $\mathbb{R}$, l'ensemble des primitives de $f$ sur $\mathbb{R}$ est :
\[ \left\{ x \mapsto \frac{x\,e^{2x}}{2} - \frac{e^{2x}}{4} + C \;\middle|\; C \in \mathbb{R} \right\}. \]
\item On cherche $C$ tel que la primitive s'annule en $0$ :
\[ \frac{0 \times e^{0}}{2} - \frac{e^{0}}{4} + C = 0 \iff -\frac{1}{4} + C = 0 \iff C = \frac{1}{4}. \]
La primitive cherchée est :
\[ \boxed{F(x) = \frac{x\,e^{2x}}{2} - \frac{e^{2x}}{4} + \frac{1}{4} = \frac{e^{2x}(2x - 1) + 1}{4}} \]
Vérification : $F(0) = \dfrac{1 \times (-1) + 1}{4} = 0$. \checkmark
\end{enumerate}
\end{corrige}

\begin{corrige}[Exercice 8 -- Forme $u'/u$]
\begin{enumerate}
\item Le dénominateur est $D(x) = 3x^2 + 2x + 1$. Son discriminant est
\[ \Delta = 2^2 - 4 \times 3 \times 1 = 4 - 12 = -8 < 0. \]
Comme $\Delta < 0$ et que le coefficient de $x^2$ est $3 > 0$, on a $D(x) > 0$ pour tout $x \in \mathbb{R}$ : le dénominateur ne s'annule jamais. $f$ est donc définie sur $\mathbb{R}$, et continue sur $\mathbb{R}$ comme quotient de fonctions polynomiales dont le dénominateur ne s'annule pas.
\item Posons $u(x) = 3x^2 + 2x + 1$. La fonction $u$ est dérivable sur $\mathbb{R}$ et $u'(x) = 6x + 2$, qui est exactement le numérateur de $f$. De plus $u(x) = D(x) > 0$ sur $\mathbb{R}$ d'après la question 1. Donc :
\[ f(x) = \frac{u'(x)}{u(x)} \quad \text{avec } u > 0 \text{ sur } \mathbb{R}. \]
\item D'après la propriété du cours (forme $\dfrac{u'}{u}$ avec $u > 0$), une primitive de $f$ sur $\mathbb{R}$ est $\ln(u(x))$. L'ensemble des primitives de $f$ sur $\mathbb{R}$ est donc :
\[ \boxed{F(x) = \ln\left(3x^2 + 2x + 1\right) + C, \quad C \in \mathbb{R}} \]
\item $F(0) = 2 \iff \ln(1) + C = 2 \iff 0 + C = 2 \iff C = 2$.
\[ \boxed{F(x) = \ln\left(3x^2 + 2x + 1\right) + 2} \]
\textbf{Point de vigilance :} la condition $u > 0$ est \emph{indispensable} pour écrire $\ln(u(x))$ ; c'est pourquoi la question 1 (étude du discriminant) doit être rédigée avec soin.
\end{enumerate}
\end{corrige}

\begin{corrige}[Exercice 9 -- Type Bac : primitives et condition initiale]
\textbf{Barème indicatif (sur 5 pts) :} question 1 : 0,5 pt ; question 2 : 1,5 pt (0,5 pt par primitive) ; question 3 : 1 pt ; question 4 : 2 pts (expression exacte 1 pt, valeur approchée 1 pt).

\begin{enumerate}
\item Sur $]0\,;+\infty[$ : $x \mapsto 4x^3$ est continue (polynôme) ; $x \mapsto -\dfrac{6}{x^2}$ est continue (fonction rationnelle dont le dénominateur ne s'annule pas sur $]0\,;+\infty[$) ; $x \mapsto 2\sin\left(3x - \dfrac{\pi}{4}\right)$ est continue sur $\mathbb{R}$. La fonction $f$, somme de fonctions continues sur l'\textbf{intervalle} $]0\,;+\infty[$, y est continue : d'après le théorème d'existence, $f$ admet des primitives sur $]0\,;+\infty[$.

\textbf{Point de vigilance :} il faut citer explicitement la \textbf{continuité} et le fait que l'ensemble de définition est un \textbf{intervalle} : ce sont les deux hypothèses du théorème.
\item \begin{enumerate}
\item Une primitive de $x \mapsto 4x^3$ est $x \mapsto x^4$.
\item Une primitive de $x \mapsto -\dfrac{6}{x^2}$ est $x \mapsto \dfrac{6}{x}$, car $\left(\dfrac{6}{x}\right)' = -\dfrac{6}{x^2}$.
\item Soit $\varphi(x) = -\dfrac{2}{3}\cos\left(3x - \dfrac{\pi}{4}\right)$. $\varphi$ est dérivable sur $\mathbb{R}$ et, par la formule $(\cos u)' = -u'\sin u$ :
\[ \varphi'(x) = -\frac{2}{3} \times \left(-\sin\left(3x - \frac{\pi}{4}\right)\right) \times 3 = 2\sin\left(3x - \frac{\pi}{4}\right). \]
Donc $\varphi$ est bien une primitive de $x \mapsto 2\sin\left(3x - \dfrac{\pi}{4}\right)$ sur $\mathbb{R}$.
\end{enumerate}
\item Par linéarité (la somme de primitives est une primitive de la somme), une primitive de $f$ est $F_0(x) = x^4 + \dfrac{6}{x} - \dfrac{2}{3}\cos\left(3x - \dfrac{\pi}{4}\right)$. L'ensemble des primitives de $f$ sur $]0\,;+\infty[$ est :
\[ \boxed{F(x) = x^4 + \frac{6}{x} - \frac{2}{3}\cos\left(3x - \frac{\pi}{4}\right) + C, \quad C \in \mathbb{R}} \]
\item On impose $F(1) = 5$ :
\begin{align*}
F(1) = 5 &\iff 1 + 6 - \frac{2}{3}\cos\left(3 - \frac{\pi}{4}\right) + C = 5 \\
&\iff C = -2 + \frac{2}{3}\cos\left(3 - \frac{\pi}{4}\right).
\end{align*}
Expression exacte :
\[ \boxed{F(x) = x^4 + \frac{6}{x} - \frac{2}{3}\cos\left(3x - \frac{\pi}{4}\right) - 2 + \frac{2}{3}\cos\left(3 - \frac{\pi}{4}\right)} \]
Valeur approchée : $\cos\left(3 - \dfrac{\pi}{4}\right) \approx -0,602$ (car $3 - \pi/4 \approx 2,215$ rad) :
\[ C \approx -2 + \frac{2}{3} \times (-0,602) \approx -2 - 0,401 = -2,401. \]
Donc $F(x) \approx x^4 + \dfrac{6}{x} - \dfrac{2}{3}\cos\left(3x - \dfrac{\pi}{4}\right) - 2,401$.

\textbf{Points de vigilance :} ne pas oublier le facteur $\dfrac{1}{a} = \dfrac{1}{3}$ en primitivant $\sin(3x - \pi/4)$ ; conserver la valeur \emph{exacte} de $C$ avant toute approximation ; vérifier le signe : $\cos(3 - \pi/4)$ est négatif, donc $C < -2$.
\end{enumerate}
\end{corrige}

\begin{corrige}[Exercice 10 -- Type Bac : linéarisation et primitives trigonométriques]
\textbf{Barème indicatif (sur 6 pts) :} question 1 : 1,5 pt ; question 2 : 2 pts ; question 3 : 1 pt ; question 4 : 1,5 pt.

\begin{enumerate}
\item \begin{enumerate}
\item Formule de duplication : $\cos(2a) = 2\cos^2 a - 1$, d'où $\cos^2 a = \dfrac{1 + \cos(2a)}{2}$.
\item En appliquant cette formule avec $a = 2x$ (donc $2a = 4x$) :
\[ \cos^2(2x) = \frac{1 + \cos(4x)}{2} \quad \text{pour tout réel } x. \]
\item On primitive la forme linéarisée :
\[ \frac{1 + \cos(4x)}{2} = \frac{1}{2} + \frac{1}{2}\cos(4x). \]
Une primitive de $\dfrac{1}{2}$ est $\dfrac{x}{2}$ ; une primitive de $\dfrac{1}{2}\cos(4x)$ est $\dfrac{1}{2} \times \dfrac{1}{4}\sin(4x) = \dfrac{1}{8}\sin(4x)$. L'ensemble des primitives de $x \mapsto \cos^2(2x)$ sur $\mathbb{R}$ est :
\[ \boxed{x \mapsto \frac{x}{2} + \frac{1}{8}\sin(4x) + C, \quad C \in \mathbb{R}} \]
\end{enumerate}
\item \begin{enumerate}
\item Pour tout réel $x$ : $\sin^3 x = \sin x \times \sin^2 x = \sin x \left(1 - \cos^2 x\right)$, car $\sin^2 x + \cos^2 x = 1$.
\item Posons $u(x) = \cos x$, donc $u'(x) = -\sin x$, c'est-à-dire $\sin x = -u'(x)$. Alors :
\[ \sin^3 x = -u'(x)\left(1 - [u(x)]^2\right) = -u'(x) + u'(x)\,[u(x)]^2. \]
On reconnaît une combinaison des formes $u'$ (primitive : $u$) et $u' \cdot u^2$ (primitive : $\dfrac{u^3}{3}$).
\item Une primitive de $-u'$ est $-u = -\cos x$ ; une primitive de $u'u^2$ est $\dfrac{u^3}{3} = \dfrac{\cos^3 x}{3}$. L'ensemble des primitives de $x \mapsto \sin^3 x$ sur $\mathbb{R}$ est :
\[ \boxed{x \mapsto -\cos x + \frac{\cos^3 x}{3} + C, \quad C \in \mathbb{R}} \]
Vérification : $\left(-\cos x + \dfrac{\cos^3 x}{3}\right)' = \sin x + \dfrac{3\cos^2 x \times (-\sin x)}{3} = \sin x (1 - \cos^2 x) = \sin^3 x$. \checkmark
\end{enumerate}
\item Par somme, l'ensemble des primitives de $f(x) = \cos^2(2x) + \sin^3 x$ sur $\mathbb{R}$ est :
\[ \boxed{F(x) = \frac{x}{2} + \frac{1}{8}\sin(4x) - \cos x + \frac{\cos^3 x}{3} + C, \quad C \in \mathbb{R}} \]
\item On évalue en $x = \dfrac{\pi}{2}$ : $\sin\left(4 \times \dfrac{\pi}{2}\right) = \sin(2\pi) = 0$ ; $\cos\dfrac{\pi}{2} = 0$ ; $\cos^3\dfrac{\pi}{2} = 0$. Donc :
\[ F\left(\frac{\pi}{2}\right) = \frac{\pi}{4} + 0 - 0 + 0 + C = \frac{\pi}{4} + C. \]
La condition $F\left(\dfrac{\pi}{2}\right) = \dfrac{\pi}{4}$ donne $C = 0$. La primitive cherchée est :
\[ \boxed{F(x) = \frac{x}{2} + \frac{1}{8}\sin(4x) - \cos x + \frac{\cos^3 x}{3}} \]
\textbf{Points de vigilance :} pour un exposant \emph{pair}, il faut \textbf{linéariser} avant de primitiver (on ne peut pas écrire directement $\dfrac{\cos^3(2x)}{3}$ !) ; pour un exposant \emph{impair}, on factorise pour faire apparaître une forme $u'u^r$ ; ne pas oublier le coefficient $\dfrac{1}{4}$ devant $\sin(4x)$ issu de la primitive de $\cos(4x)$.
\end{enumerate}
\end{corrige}

\begin{corrige}[Exercice 11 -- Problème de synthèse : le moto-taxi de Buea]
\begin{enumerate}
\item $v$ est une fonction polynôme, donc continue sur $\mathbb{R}$, et en particulier sur l'intervalle $[0\,;5]$. D'après le théorème d'existence, toute fonction continue sur un intervalle y admet des primitives : $v$ admet donc des primitives sur $[0\,;5]$.
\item On primitive terme à terme : $3t^2 \to t^3$, $-2t \to -t^2$, $1 \to t$. L'ensemble des primitives de $v$ sur $[0\,;5]$ est :
\[ \left\{ t \mapsto t^3 - t^2 + t + C \;\middle|\; C \in \mathbb{R} \right\}. \]
\item La distance $x$ est une primitive de $v$, donc $x(t) = t^3 - t^2 + t + C$. La condition initiale $x(0) = 10$ donne $0 - 0 + 0 + C = 10$, soit $C = 10$. Par conséquent :
\[ \boxed{x(t) = t^3 - t^2 + t + 10} \]
\item Au bout de $5$ secondes :
\[ x(5) = 5^3 - 5^2 + 5 + 10 = 125 - 25 + 5 + 10 = 115. \]
Le moto-taxi a parcouru au total $\boxed{115 \text{ mètres}}$ depuis son point de stationnement.
\item Montant de la course : $115 \times 25 = 2875$. Selon ce modèle, la course coûterait $\boxed{2875 \text{ FCFA}}$.
\item \textbf{Distance du second moto-taxi :} $w(t) = 2t + 4$ est continue sur $[0\,;5]$ ; une primitive est $t^2 + 4t$. La condition $y(0) = 0$ donne une constante nulle :
\[ y(t) = t^2 + 4t. \]
\textbf{Résolution de $x(t) = y(t)$ :}
\[ t^3 - t^2 + t + 10 = t^2 + 4t \iff t^3 - 2t^2 - 3t + 10 = 0. \]
Posons $P(t) = t^3 - 2t^2 - 3t + 10$ et étudions son signe sur $[0\,;5]$. $P$ est dérivable et $P'(t) = 3t^2 - 4t - 3$. Le discriminant de $P'$ est $\Delta = 16 + 36 = 52$, d'où les racines $t_1 = \dfrac{4 - \sqrt{52}}{6} \approx -0,54$ (hors de $[0,5]$) et $t_2 = \dfrac{4 + \sqrt{52}}{6} \approx 1,87$. Sur $[0\,;5]$, $P$ est donc décroissante sur $[0\,;t_2]$ puis croissante sur $[t_2\,;5]$. Son minimum sur $[0\,;5]$ est atteint en $t_2 \approx 1,87$ :
\[ P(1,87) \approx (1,87)^3 - 2(1,87)^2 - 3(1,87) + 10 \approx 6,54 - 6,99 - 5,61 + 10 \approx 3,94 > 0. \]
Le minimum de $P$ sur $[0\,;5]$ est strictement positif, donc $P(t) > 0$ pour tout $t \in [0\,;5]$ : l'équation $x(t) = y(t)$ \textbf{n'admet aucune solution} sur $[0\,;5]$.

\textbf{Interprétation :} comme $x(t) - y(t) = P(t) > 0$ sur tout l'intervalle, le premier moto-taxi — qui avait $10$ m d'avance au départ — reste \textbf{strictement devant} le second pendant toutes les cinq premières secondes. L'écart minimal entre les deux est d'environ $3,9$ m, atteint vers $t \approx 1,9$ s : le second se rapproche, mais ne rattrape jamais le premier sur cette durée.
\end{enumerate}
\end{corrige}

\newpage

\coursec{Fiche bilan}

\begin{popfigure}
\begin{tikzpicture}[mindmap, grow cyclic, every node/.style={concept, text=white, font=\small\bfseries},
root concept/.append style={concept color=popdark, font=\normalsize\bfseries},
level 1/.append style={level distance=3.4cm, sibling angle=60},
level 2/.append style={level distance=2.1cm, sibling angle=45, font=\scriptsize\bfseries}]
\node[root concept]{Primitives}
  child[concept color=popblue]{ node {Définition \& Existence}
    child { node {$F'=f$ sur $I$} }
    child { node {Théorème : $f$ continue $\Rightarrow$ primitives} }
  }
  child[concept color=poporange]{ node {Ensemble des primitives}
    child { node {$F+k,\; k\in\mathbb{R}$} }
    child { node {Condition initiale $\Rightarrow$ unicité} }
  }
  child[concept color=popgreen]{ node {Primitives usuelles}
    child { node {$x^r\;(r\neq -1)$, $\frac{1}{x}\to\ln|x|$} }
    child { node {$\cos(ax+b)$, $\sin(ax+b)$} }
    child { node {$\mathrm{e}^x$, $\frac{1}{\cos^2 x}$} }
  }
  child[concept color=poppurple]{ node {Formes composées}
    child { node {$u'\,u^r\;(r\neq -1)$} }
    child { node {$\frac{a}{(cx+d)^r}\;(r\neq 1)$} }
    child { node {$\frac{u'}{u}\to\ln|u|$} }
    child { node {$u'\mathrm{e}^u\to\mathrm{e}^u$} }
  }
  child[concept color=poppink]{ node {Puissances de cos, sin}
    child { node {$n$ impair : substitution $u=\sin x$ ou $\cos x$} }
    child { node {$n$ pair : linéarisation} }
  }
  child[concept color=popgold]{ node {Stratégie Bac}
    child { node {Linéarité, ajustement constante} }
    child { node {Vérification par dérivation} }
  };
\end{tikzpicture}
\legende{Carte mentale de la leçon M05 : les six idées forces sur les primitives.}
\end{popfigure}

\coursec{Notion de primitive et existence}

\begin{definition}[Primitive d'une fonction sur un intervalle]
Soit $I$ un intervalle de $\mathbb{R}$ et $f : I \to \mathbb{R}$ une fonction.
On dit qu'une fonction $F : I \to \mathbb{R}$ est une \cle{primitive} de $f$ sur $I$ si $F$ est dérivable sur $I$ et si, pour tout $x \in I$,
\[
F'(x) = f(x).
\]
\end{definition}

\begin{attention}[L'intervalle est indispensable]
La définition d'une primitive exige que $f$ et $F$ soient définies sur un \textbf{même intervalle} $I$.
Si $f$ n'est pas définée sur un intervalle (par exemple $f(x)=\frac{1}{x}$ sur $\mathbb{R}^*$), on ne parle pas de primitive \emph{sur $\mathbb{R}^*$}, mais de primitives sur \emph{chacun des intervalles} composant le domaine ($]-\infty,0[$ et $]0,+\infty[$).
Au Bac, précisez toujours l'intervalle considéré.
\end{attention}

\begin{propriete}[Existence des primitives (Théorème fondamental)]
Toute fonction \cle{continue} sur un intervalle $I$ admet au moins une primitive sur $I$.
Plus précisément, pour tout $a \in I$, la fonction $F(x) = \int_a^x f(t)\,dt$ est une primitive de $f$ sur $I$.
\end{propriete}

\begin{exemplebox}[Exemples et contre-exemples]
\begin{enumerate}
  \item $f(x)=x^2$ est continue sur $\mathbb{R}$. $F(x)=\frac{x^3}{3}$ est une primitive sur $\mathbb{R}$ car $F'(x)=x^2$.
  \item $f(x)=\frac{1}{x}$ est continue sur $]-\infty,0[$ et sur $]0,+\infty[$. $F(x)=\ln|x|$ est une primitive sur \emph{chacun} de ces deux intervalles.
  \item $f(x)=\lfloor x \rfloor$ (partie entière) n'est pas continue sur $\mathbb{R}$ (discontinuités en les entiers). Elle n'admet \textbf{pas} de primitive sur $\mathbb{R}$ (une fonction dérivable est nécessairement continue).
\end{enumerate}
\end{exemplebox}

\coursec{Ensemble des primitives ; condition initiale}

\begin{propriete}[Ensemble des primitives]
Soit $f$ une fonction continue sur un intervalle $I$ et $F_0$ une primitive de $f$ sur $I$.
L'ensemble de \textbf{toutes} les primitives de $f$ sur $I$ est :
\[
\{ F : I \to \mathbb{R} \mid \exists k \in \mathbb{R},\; \forall x \in I,\; F(x) = F_0(x) + k \}.
\]
On note souvent cet ensemble $\int f(x)\,dx = F_0(x) + k,\quad k \in \mathbb{R}$.
\end{propriete}

\begin{propriete}[Primitive avec condition initiale]
Soit $f$ continue sur $I$, $x_0 \in I$ et $y_0 \in \mathbb{R}$.
Il existe une \textbf{unique} primitive $F$ de $f$ sur $I$ telle que $F(x_0) = y_0$.
Elle s'obtient en prenant $k = y_0 - F_0(x_0)$ où $F_0$ est une primitive quelconque de $f$.
\end{propriete}

\begin{methode}[Déterminer la primitive avec condition initiale]
\begin{enumerate}
  \item Calculer une primitive $F_0$ de $f$ sur l'intervalle $I$ (omettre la constante $k$ dans un premier temps).
  \item Écrire la forme générale : $F(x) = F_0(x) + k,\; k \in \mathbb{R}$.
  \item Utiliser la condition $F(x_0) = y_0$ : résoudre $F_0(x_0) + k = y_0$ pour trouver $k$.
  \item Écrire la primitive finale $F(x) = F_0(x) + k$ en précisant l'intervalle $I$.
\end{enumerate}
\end{methode}

\begin{exoresolu}[Primitive avec condition initiale]
Déterminer la primitive $F$ de $f(x) = 3x^2 - 2\sin x$ sur $\mathbb{R}$ telle que $F(0) = 5$.
\tcblower
\textbf{Solution.}
\begin{enumerate}
  \item Une primitive de $f$ sur $\mathbb{R}$ est $F_0(x) = x^3 + 2\cos x$ (car $(x^3)'=3x^2$ et $(2\cos x)'=-2\sin x$).
  \item Forme générale : $F(x) = x^3 + 2\cos x + k,\; k \in \mathbb{R}$.
  \item Condition : $F(0) = 0^3 + 2\cos 0 + k = 2 + k = 5 \;\Rightarrow\; k = 3$.
  \item La primitive cherchée est $F(x) = x^3 + 2\cos x + 3$ sur $\mathbb{R}$.
\end{enumerate}
\textbf{Vérification :} $F'(x) = 3x^2 - 2\sin x = f(x)$ et $F(0)=3+2=5$. \checkmark
\end{exoresolu}

\coursec{Primitives des fonctions usuelles}

\begin{propriete}[Linéarité de la primitive]
Si $f$ et $g$ admettent des primitives sur $I$ et $\lambda, \mu \in \mathbb{R}$, alors $\lambda f + \mu g$ admet des primitives sur $I$ et
\[
\int (\lambda f(x) + \mu g(x))\,dx = \lambda \int f(x)\,dx + \mu \int g(x)\,dx.
\]
\end{propriete}

\begin{aretenir}[Tableau des primitives usuelles (à connaître par cœur)]
\begin{center}
\rowcolors{1}{popblueL}{white}
\begin{tabularx}{\linewidth}{|l|X|l|}\hline
\rowcolor{popblue!40} \textbf{Fonction $f(x)$} & \textbf{Primitive $F(x)$} & \textbf{Intervalle de validité} \\ \hline
$k$ (constante) & $kx$ & $\mathbb{R}$ \\ \hline
$x^r$ \quad ($r \in \mathbb{Q},\; r \neq -1$) & $\dfrac{x^{r+1}}{r+1}$ & Dépend de $r$ (ex: $r=-2 \Rightarrow ]-\infty,0[\cup]0,+\infty[$) \\ \hline
$\dfrac{1}{x}$ & $\ln|x|$ & $]-\infty,0[$ ou $]0,+\infty[$ \\ \hline
$\mathrm{e}^{x}$ & $\mathrm{e}^{x}$ & $\mathbb{R}$ \\ \hline
$\cos(ax+b)$ \quad ($a \neq 0$) & $\dfrac{1}{a}\sin(ax+b)$ & $\mathbb{R}$ \\ \hline
$\sin(ax+b)$ \quad ($a \neq 0$) & $-\dfrac{1}{a}\cos(ax+b)$ & $\mathbb{R}$ \\ \hline
$\dfrac{1}{\cos^2(ax+b)} = 1+\tan^2(ax+b)$ \quad ($a \neq 0$) & $\dfrac{1}{a}\tan(ax+b)$ & Intervalle où $\cos(ax+b)\neq 0$ \\ \hline
\end{tabularx}
\end{center}
\end{aretenir}

\begin{attention}[Pièges classiques sur les usuelles]
\begin{itemize}
  \item \textbf{Signe pour le sinus :} $\int \sin(ax+b)\,dx = -\frac{1}{a}\cos(ax+b) + k$. Le signe $-$ est systématiquement oublié.
  \item \textbf{Facteur $\frac{1}{a}$ :} Pour $\cos(ax+b)$ et $\sin(ax+b)$, la dérivation de $ax+b$ donne $a$, donc la primitive divise par $a$.
  \item \textbf{Domaine pour $x^r$ :} Si $r = -\frac{1}{2}$ ($\frac{1}{\sqrt{x}}$), l'intervalle est $]0,+\infty[$. Si $r=-2$ ($\frac{1}{x^2}$), ce sont $]-\infty,0[$ et $]0,+\infty[$.
  \item \textbf{$\frac{1}{x}$ :} La primitive est $\ln|x|$ (valeur absolue), pas $\ln x$. Valable sur $]-\infty,0[$ \textbf{et} $]0,+\infty[$ séparément.
\end{itemize}
\end{attention}

\begin{exemplebox}[Applications de la linéarité]
\begin{enumerate}
  \item $\displaystyle \int (4x^3 - 5\sqrt{x} + \frac{2}{x^2})\,dx$ sur $]0,+\infty[$.
  $\sqrt{x}=x^{1/2}$, $\frac{2}{x^2}=2x^{-2}$.
  $F(x) = 4\frac{x^4}{4} - 5\frac{x^{3/2}}{3/2} + 2\frac{x^{-1}}{-1} + k = x^4 - \frac{10}{3}x^{3/2} - \frac{2}{x} + k$.
  \item $\displaystyle \int (3\cos(2x) - 4\sin(x-\pi/3))\,dx$ sur $\mathbb{R}$.
  $F(x) = 3\cdot\frac{1}{2}\sin(2x) - 4\cdot(-\cos(x-\pi/3)) + k = \frac{3}{2}\sin(2x) + 4\cos(x-\pi/3) + k$.
\end{enumerate}
\end{exemplebox}

\coursec{Primitives des formes composées : $u'\,u^r$ et $\frac{a}{(cx+d)^r}$}

\begin{propriete}[Forme $u'\,u^r$]
Soit $u$ une fonction dérivable sur un intervalle $I$ et $r \in \mathbb{Q}$, $r \neq -1$.
Si $u(x) > 0$ sur $I$ (ou $r$ entier relatif), alors une primitive de $u'(x)\,u(x)^r$ sur $I$ est
\[
\frac{u(x)^{r+1}}{r+1}.
\]
\end{propriete}

\begin{propriete}[Forme $\frac{a}{(cx+d)^r}$]
Soient $a,c,d \in \mathbb{R}$, $c \neq 0$, $r \in \mathbb{Q}$, $r \neq 1$.
Sur tout intervalle $I$ où $cx+d \neq 0$ (et $cx+d > 0$ si $r \notin \mathbb{Z}$), une primitive de $\frac{a}{(cx+d)^r}$ est
\[
\frac{a}{c(1-r)}\,(cx+d)^{1-r} = \frac{a}{c(1-r)(cx+d)^{r-1}}.
\]
\end{propriete}

\begin{methode}[Reconnaître et traiter la forme $u'\,u^r$]
\begin{enumerate}
  \item Identifier une fonction $u(x)$ « candidate » (souvent ce qui est entre parenthèses, au dénominateur, ou sous racine).
  \item Calculer $u'(x)$.
  \item Vérifier que l'expression à primitiver s'écrit $C \cdot u'(x) \cdot u(x)^r$ où $C$ est une constante.
  \item La primitive est $C \cdot \frac{u(x)^{r+1}}{r+1} + k$.
  \item \textbf{Vérifier en dérivant.}
\end{enumerate}
\end{methode}

\begin{exoresolu}[Formes composées : $u'\,u^r$ et $\frac{a}{(cx+d)^r}$]
\begin{enumerate}
  \item $\displaystyle \int \frac{2x}{(x^2+3)^4}\,dx$ sur $\mathbb{R}$.
  \item $\displaystyle \int \frac{5}{(2x-1)^3}\,dx$ sur $]\frac{1}{2},+\infty[$.
  \item $\displaystyle \int (3x^2-1)\sqrt{x^3-x}\,dx$ sur un intervalle où $x^3-x \ge 0$.
\end{enumerate}
\tcblower
\textbf{Solutions.}
\begin{enumerate}
  \item $u(x)=x^2+3$, $u'(x)=2x$, $r=-4$. Forme $u'\,u^{-4}$.
  Primitive : $\frac{u^{-3}}{-3} = -\frac{1}{3(x^2+3)^3} + k$.
  \item Forme $\frac{a}{(cx+d)^r}$ avec $a=5$, $c=2$, $d=-1$, $r=3$.
  Primitive : $\frac{5}{2(1-3)}(2x-1)^{-2} = -\frac{5}{4(2x-1)^2} + k$.
  \item $u(x)=x^3-x$, $u'(x)=3x^2-1$, $r=\frac{1}{2}$. Forme $u'\,u^{1/2}$.
  Primitive : $\frac{u^{3/2}}{3/2} = \frac{2}{3}(x^3-x)^{3/2} + k$.
\end{enumerate}
\end{exoresolu}

\coursec{Primitives de $\cos^n x$ et $\sin^n x$}

\begin{propriete}[Formules de linéarisation (rappels)]
Pour tout $x \in \mathbb{R}$ :
\[
\cos^2 x = \frac{1+\cos 2x}{2}, \qquad \sin^2 x = \frac{1-\cos 2x}{2}.
\]
Plus généralement, on utilise les formules d'Euler ou de Moivre pour linéariser $\cos^n x$ et $\sin^n x$ ($n$ pair).
\end{propriete}

\begin{methode}[Primitive de $\cos^n x$ ou $\sin^n x$ ($n \in \mathbb{N}^*$)]
\begin{enumerate}
  \item \textbf{Si $n$ est impair} ($n=2p+1$) :
  \begin{itemize}
    \item Pour $\cos^{2p+1} x$ : écrire $\cos^{2p+1} x = \cos^{2p} x \cdot \cos x = (1-\sin^2 x)^p \cdot \cos x$.
    Poser $u = \sin x$, $du = \cos x\,dx$. On obtient $\int (1-u^2)^p\,du$ (polynôme en $u$).
    \item Pour $\sin^{2p+1} x$ : écrire $\sin^{2p+1} x = \sin^{2p} x \cdot \sin x = (1-\cos^2 x)^p \cdot \sin x$.
    Poser $u = \cos x$, $du = -\sin x\,dx$. On obtient $-\int (1-u^2)^p\,du$.
  \end{itemize}
  \item \textbf{Si $n$ est pair} ($n=2p$) :
  Linéariser $\cos^{2p} x$ ou $\sin^{2p} x$ en combinaisons de $\cos(kx)$ ($k$ pair) à l'aide des formules de Moivre/Euler, puis primitiver terme à terme.
\end{enumerate}
\end{methode}

\begin{exoresolu}[Puissances de cosinus et sinus]
\begin{enumerate}
  \item Calculer $\int \cos^3 x\,dx$ sur $\mathbb{R}$.
  \item Calculer $\int \sin^4 x\,dx$ sur $\mathbb{R}$.
\end{enumerate}
\tcblower
\textbf{Solutions.}
\begin{enumerate}
  \item $n=3$ impair. $\cos^3 x = \cos^2 x \cdot \cos x = (1-\sin^2 x)\cos x$.
  Poser $u=\sin x$, $du=\cos x\,dx$.
  $\int (1-u^2)\,du = u - \frac{u^3}{3} + k = \sin x - \frac{\sin^3 x}{3} + k$.
  \item $n=4$ pair. Linéarisation :
  $\sin^4 x = (\sin^2 x)^2 = \left(\frac{1-\cos 2x}{2}\right)^2 = \frac{1}{4}(1 - 2\cos 2x + \cos^2 2x)$.
  $\cos^2 2x = \frac{1+\cos 4x}{2}$.
  $\sin^4 x = \frac{1}{4}\left(1 - 2\cos 2x + \frac{1}{2} + \frac{1}{2}\cos 4x\right) = \frac{3}{8} - \frac{1}{2}\cos 2x + \frac{1}{8}\cos 4x$.
  Primitive : $\frac{3}{8}x - \frac{1}{2}\cdot\frac{1}{2}\sin 2x + \frac{1}{8}\cdot\frac{1}{4}\sin 4x + k = \frac{3}{8}x - \frac{1}{4}\sin 2x + \frac{1}{32}\sin 4x + k$.
\end{enumerate}
\end{exoresolu}

\coursec{Compléments utiles au Bac : $\frac{u'}{u}$, $u'\mathrm{e}^u$ ; synthèse méthodologique}

\begin{propriete}[Formes composées supplémentaires (Programme Tle D, E, TI ; Tle C : $\ln$ vu en cours)]
Soit $u$ une fonction dérivable et strictement positive (ou non nulle pour $\ln|u|$) sur un intervalle $I$.
\begin{itemize}
  \item $\displaystyle \int \frac{u'(x)}{u(x)}\,dx = \ln|u(x)| + k$.
  \item $\displaystyle \int u'(x)\,\mathrm{e}^{u(x)}\,dx = \mathrm{e}^{u(x)} + k$.
  \item \textbf{Cas particulier :} $\int \frac{u'(x)}{\sqrt{u(x)}}\,dx = 2\sqrt{u(x)} + k$ (forme $u'\,u^{-1/2}$).
\end{itemize}
\end{propriete}

\begin{attention}[Domaine pour $\ln|u|$]
La primitive $\ln|u(x)|$ est définie sur les intervalles où $u(x) \neq 0$.
Si $u$ change de signe, on doit travailler sur des intervalles où $u$ garde un signe constant.
Au Bac, on écrit $\ln|u(x)|$ pour couvrir les deux cas ($u>0$ et $u<0$).
\end{attention}

\begin{exoresolu}[Formes $\frac{u'}{u}$ et $u'\mathrm{e}^u$]
\begin{enumerate}
  \item $\displaystyle \int \frac{2x}{x^2+1}\,dx$ sur $\mathbb{R}$.
  \item $\displaystyle \int (2x+3)\,\mathrm{e}^{x^2+3x}\,dx$ sur $\mathbb{R}$.
  \item $\displaystyle \int \frac{\cos x}{\sin x}\,dx$ sur $]0,\pi[$.
\end{enumerate}
\tcblower
\textbf{Solutions.}
\begin{enumerate}
  \item $u(x)=x^2+1$, $u'(x)=2x$. Forme $\frac{u'}{u}$.
  Primitive : $\ln|x^2+1| + k = \ln(x^2+1) + k$ (car $x^2+1>0$).
  \item $u(x)=x^2+3x$, $u'(x)=2x+3$. Forme $u'\,\mathrm{e}^u$.
  Primitive : $\mathrm{e}^{x^2+3x} + k$.
  \item $u(x)=\sin x$, $u'(x)=\cos x$. Sur $]0,\pi[$, $\sin x > 0$.
  Primitive : $\ln(\sin x) + k$.
\end{enumerate}
\end{exoresolu}

\begin{methode}[Synthèse : Stratégie complète pour primitiver au Bac]
\begin{enumerate}
  \item \textbf{Identifier l'intervalle $I$} (domaine de la fonction, continuité).
  \item \textbf{Simplifier / Développer} l'expression (racines $\to$ puissances rationnelles, fractions $\to$ sommes, trigonométrie $\to$ linéarisation si puissances paires).
  \item \textbf{Appliquer la linéarité} : primitiver terme à terme.
  \item \textbf{Pour chaque terme}, tester dans l'ordre :
  \begin{enumerate}
    \item Est-ce une \textbf{usuelle} directe ? ($x^r$, $\cos(ax+b)$, $\mathrm{e}^x$, $\frac{1}{x}$, etc.)
    \item Est-ce de la forme $\frac{a}{(cx+d)^r}$ ? $\to$ facteur $\frac{1}{c}$.
    \item Est-ce de la forme $u'\,u^r$ ? $\to$ repérer $u$, vérifier $u'$, ajuster constante.
    \item Est-ce de la forme $\frac{u'}{u}$ ? $\to \ln|u|$.
    \item Est-ce de la forme $u'\,\mathrm{e}^u$ ? $\to \mathrm{e}^u$.
    \item Puissances de $\cos/\sin$ : impair $\to$ substitution ; pair $\to$ linéarisation.
  \end{enumerate}
  \item \textbf{Réunir les primitives} en une seule expression $F(x) + k$.
  \item \textbf{VÉRIFIER SYSTÉMATIQUEMENT} en dérivant $F(x)$ : on doit retrouver exactement $f(x)$.
  \item Si condition initiale : calculer $k$ et écrire la primitive \textbf{unique} finale.
\end{enumerate}
\end{methode}

\begin{exoresolu}[Synthèse : Exercice type Bac]
Soit $f(x) = \frac{3x^2}{\sqrt{x^3+1}} + 2\cos(2x) - \frac{4}{x}$ définie sur $]0,+\infty[$.
Déterminer la primitive $F$ de $f$ sur $]0,+\infty[$ telle que $F(1) = 0$.
\tcblower
\textbf{Solution.}
\begin{enumerate}
  \item \textbf{Intervalle :} $]0,+\infty[$ ($x>0$ pour $\frac{4}{x}$ et $\sqrt{x^3+1}$ toujours défini).
  \item \textbf{Terme 1 :} $\frac{3x^2}{\sqrt{x^3+1}} = 3x^2 (x^3+1)^{-1/2}$.
  $u(x)=x^3+1$, $u'(x)=3x^2$, $r=-\frac{1}{2}$. Forme $u'\,u^{-1/2}$.
  Le facteur $3$ est \textbf{déjà} dans $u'(x)=3x^2$ : on ne le remet pas une seconde fois !
  Primitive : $\dfrac{u^{1/2}}{1/2} = 2\sqrt{u} = 2\sqrt{x^3+1}$.
  \item \textbf{Terme 2 :} $2\cos(2x)$. Usuelle : $2 \cdot \frac{1}{2}\sin(2x) = \sin(2x)$.
  \item \textbf{Terme 3 :} $-\frac{4}{x} = -4 \cdot \frac{1}{x}$. Usuelle : $-4\ln|x| = -4\ln x$ (car $x>0$).
  \item \textbf{Forme générale :} $F(x) = 2\sqrt{x^3+1} + \sin(2x) - 4\ln x + k$, $k\in\mathbb{R}$.
  \item \textbf{Condition initiale :} $F(1) = 2\sqrt{2} + \sin 2 - 4\ln 1 + k = 2\sqrt{2} + \sin 2 + k = 0$, donc $k = -2\sqrt{2} - \sin 2$.
  \item \textbf{Primitive cherchée :}
  \[
  F(x) = 2\sqrt{x^3+1} + \sin(2x) - 4\ln x - 2\sqrt{2} - \sin 2 \quad \text{sur } ]0,+\infty[.
  \]
  \textbf{Vérification finale :} $F'(x) = \frac{3x^2}{\sqrt{x^3+1}} + 2\cos(2x) - \frac{4}{x} = f(x)$. \checkmark
\end{enumerate}
\end{exoresolu}

\coursec{Exercices d'entraînement}

\begin{exercice}[Exercice 1 : Primitives usuelles et composées]
Calculer les primitives sur les intervalles indiqués :
\begin{enumerate}
  \item $f(x) = 5x^4 - \frac{3}{x^2} + 2\sqrt{x}$ sur $]0,+\infty[$.
  \item $g(x) = \frac{2}{(3x-1)^2}$ sur $]\frac{1}{3},+\infty[$.
  \item $h(x) = (2x+1)\sqrt{x^2+x}$ sur un intervalle où $x^2+x \ge 0$.
  \item $p(x) = 3\cos(2x+\frac{\pi}{4}) - 4\sin(x-\frac{\pi}{6})$ sur $\mathbb{R}$.
\end{enumerate}
\end{exercice}

\begin{exercice}[Exercice 2 : Formes $\frac{u'}{u}$ et $u'\mathrm{e}^u$]
Déterminer une primitive sur l'intervalle convenable :
\begin{enumerate}
  \item $f(x) = \frac{2x+1}{x^2+x+1}$ sur $\mathbb{R}$.
  \item $g(x) = (3x^2-2x)\,\mathrm{e}^{x^3-x^2}$ sur $\mathbb{R}$.
  \item $h(x) = \frac{\cos x}{2+\sin x}$ sur $\mathbb{R}$.
  \item $k(x) = \frac{\mathrm{e}^x}{\mathrm{e}^x+1}$ sur $\mathbb{R}$.
\end{enumerate}
\end{exercice}

\begin{exercice}[Exercice 3 : Puissances de cosinus et sinus]
Calculer les primitives sur $\mathbb{R}$ :
\begin{enumerate}
  \item $\cos^5 x$
  \item $\sin^2 x \cos^2 x$
  \item $\cos^3 x \sin^2 x$
\end{enumerate}
\end{exercice}

\begin{exercice}[Exercice 4 : Condition initiale]
Trouver la primitive $F$ de $f(x) = \frac{1}{x^2} + \sin x$ sur $]0,+\infty[$ telle que $F(\frac{\pi}{2}) = 1$.
\end{exercice}

\begin{corrige}[Exercice 1]
\begin{enumerate}
  \item $F(x) = x^5 + \frac{3}{x} + \frac{4}{3}x^{3/2} + k$.
  \item $G(x) = -\frac{2}{3(3x-1)} + k$ (car $\int \frac{2}{(3x-1)^2}dx = \frac{2}{3}\int \frac{3}{(3x-1)^2}dx = \frac{2}{3}\frac{(3x-1)^{-1}}{-1}$).
  \item $u=x^2+x$, $u'=2x+1$, $r=1/2$. $H(x) = \frac{2}{3}(x^2+x)^{3/2} + k$ sur $]-\infty,-1]$ ou $[0,+\infty[$.
  \item $P(x) = \frac{3}{2}\sin(2x+\frac{\pi}{4}) + 4\cos(x-\frac{\pi}{6}) + k$.
\end{enumerate}
\end{corrige}

\begin{corrige}[Exercice 2]
\begin{enumerate}
  \item $u=x^2+x+1>0$, $u'=2x+1$. $F(x) = \ln(x^2+x+1) + k$.
  \item $u=x^3-x^2$, $u'=3x^2-2x$. $G(x) = \mathrm{e}^{x^3-x^2} + k$.
  \item $u=2+\sin x > 0$, $u'=\cos x$. $H(x) = \ln(2+\sin x) + k$.
  \item $u=\mathrm{e}^x+1 > 0$, $u'=\mathrm{e}^x$. $K(x) = \ln(\mathrm{e}^x+1) + k$.
\end{enumerate}
\end{corrige}

\begin{corrige}[Exercice 3]
\begin{enumerate}
  \item $\cos^5 x = (1-\sin^2 x)^2 \cos x$. $u=\sin x$. $\int (1-u^2)^2 du = \int (1-2u^2+u^4)du = u - \frac{2}{3}u^3 + \frac{1}{5}u^5 + k = \sin x - \frac{2}{3}\sin^3 x + \frac{1}{5}\sin^5 x + k$.
  \item $\sin^2 x \cos^2 x = \frac{1}{4}\sin^2 2x = \frac{1}{4}\cdot\frac{1-\cos 4x}{2} = \frac{1}{8} - \frac{1}{8}\cos 4x$.
  Primitive : $\frac{1}{8}x - \frac{1}{32}\sin 4x + k$.
  \item $\cos^3 x \sin^2 x = \cos x (1-\sin^2 x)\sin^2 x = \cos x (\sin^2 x - \sin^4 x)$. $u=\sin x$.
  $\int (u^2 - u^4)du = \frac{1}{3}u^3 - \frac{1}{5}u^5 + k = \frac{1}{3}\sin^3 x - \frac{1}{5}\sin^5 x + k$.
\end{enumerate}
\end{corrige}

\begin{corrige}[Exercice 4]
$f(x) = x^{-2} + \sin x$ sur $]0,+\infty[$.
Primitive générale : $F(x) = -\frac{1}{x} - \cos x + k$.
Condition : $F(\frac{\pi}{2}) = -\frac{2}{\pi} - \cos\frac{\pi}{2} + k = -\frac{2}{\pi} + k = 1 \Rightarrow k = 1 + \frac{2}{\pi}$.
Réponse : $F(x) = -\frac{1}{x} - \cos x + 1 + \frac{2}{\pi}$ sur $]0,+\infty[$.
\end{corrige}

\begin{bilan}

\textbf{Définitions et théorèmes clés}
\begin{itemize}
  \item $F$ est une \cle{primitive} de $f$ sur l'intervalle $I$ $\iff$ $F$ dérivable sur $I$ et $F'(x)=f(x),\ \forall x\in I$.
  \item \textbf{Théorème d'existence :} Toute fonction \cle{continue} sur un intervalle $I$ admet des primitives sur $I$.
  \item Si $F_0$ est une primitive de $f$ sur $I$, \textbf{toutes} les primitives sont $F(x) = F_0(x) + k,\ k\in\mathbb{R}$.
  \item Donnée une condition initiale $F(x_0)=y_0$ ($x_0\in I$), il existe une \textbf{unique} primitive : $k = y_0 - F_0(x_0)$.
\end{itemize}

\textbf{Tableau des primitives usuelles (à connaître par cœur)}
\begin{center}
\rowcolors{1}{popblueL}{white}
\begin{tabularx}{\linewidth}{|l|X|l|}\hline
\rowcolor{popblue!40} \textbf{Fonction $f(x)$} & \textbf{Primitive $F(x)$} & \textbf{Intervalle} \\ \hline
$k$ (constante) & $kx$ & $\mathbb{R}$ \\ \hline
$x^r$ \quad ($r\in\mathbb{Q},\ r\neq -1$) & $\dfrac{x^{r+1}}{r+1}$ & Selon $r$ ($x>0$ si $r\notin\mathbb{Z}$) \\ \hline
$\dfrac{1}{x}$ & $\ln|x|$ & $]-\infty,0[$ ou $]0,+\infty[$ \\ \hline
$\mathrm{e}^{x}$ & $\mathrm{e}^{x}$ & $\mathbb{R}$ \\ \hline
$\cos(ax+b)$ \quad ($a\neq 0$) & $\dfrac{1}{a}\sin(ax+b)$ & $\mathbb{R}$ \\ \hline
$\sin(ax+b)$ \quad ($a\neq 0$) & $-\dfrac{1}{a}\cos(ax+b)$ & $\mathbb{R}$ \\ \hline
$\dfrac{1}{\cos^2(ax+b)}$ \quad ($a\neq 0$) & $\dfrac{1}{a}\tan(ax+b)$ & $\cos(ax+b)\neq 0$ \\ \hline
\end{tabularx}
\end{center}

\textbf{Formes composées (reconnaissance immédiate)}
\begin{center}
\rowcolors{1}{popgreenL}{white}
\begin{tabularx}{\linewidth}{|X|X|}\hline
\rowcolor{popgreen!40} \textbf{Forme} & \textbf{Primitive} \\ \hline
$u'\,u^{r}$ \quad ($r\neq -1$) & $\dfrac{u^{r+1}}{r+1}$ \\ \hline
$\dfrac{a}{(cx+d)^r}$ \quad ($r\neq 1$) & $\dfrac{a}{c(1-r)(cx+d)^{r-1}}$ \\ \hline
$\dfrac{u'}{u}$ & $\ln|u|$ \\ \hline
$u'\,\mathrm{e}^{u}$ & $\mathrm{e}^{u}$ \\ \hline
$\dfrac{u'}{\sqrt{u}}$ \quad ($u>0$) & $2\sqrt{u}$ \\ \hline
\end{tabularx}
\end{center}

\textbf{Méthodes express}
\begin{itemize}
  \item \cle{Linéarité} : primitiver terme à terme.
  \item \cle{Forme $u'\,u^r$} : 1) Repérer $u$ ; 2) Vérifier $u'$ (à constante près) ; 3) Ajuster la constante ; 4) Appliquer $\frac{u^{r+1}}{r+1}$.
  \item \cle{$\cos^n x$ ou $\sin^n x$} :
  \begin{itemize}
    \item $n$ \textbf{impair} : factoriser un $\cos x$ (resp. $\sin x$), exprimer le reste en $\sin x$ (resp. $\cos x$), substitution $u=\sin x$ (resp. $u=\cos x$).
    \item $n$ \textbf{pair} : linéariser avec $\cos^2 x=\frac{1+\cos 2x}{2}$, $\sin^2 x=\frac{1-\cos 2x}{2}$ (ou formules de Moivre).
  \end{itemize}
  \item \cle{Condition initiale} : $F(x_0)=y_0 \Rightarrow k = y_0 - F_0(x_0)$.
  \item \textbf{Vérification ultime} : dériver le résultat $\to$ doit retrouver $f(x)$ exactement.
\end{itemize}
\end{bilan}

\begin{aretenir}[Checklist finale avant le Bac]
Avant l'épreuve, assure-toi de savoir :
\begin{itemize}
  \item $\square$ Énoncer la définition exacte d'une primitive ($F'=f$ sur un \emph{intervalle}).
  \item $\square$ Citer le théorème : « Toute fonction continue sur un intervalle admet des primitives ».
  \item $\square$ Écrire l'ensemble des primitives sous la forme $F+k$, $k\in\mathbb{R}$, sans oublier la constante $k$.
  \item $\square$ Déterminer la constante $k$ grâce à une condition initiale $F(x_0)=y_0$.
  \item $\square$ Citer sans hésiter les primitives de $x^r$ ($r\in\mathbb{Q}$, $r\neq -1$), $\frac{1}{x^2}$, $\frac{1}{\sqrt{x}}$, $\frac{1}{x}\to\ln|x|$.
  \item $\square$ Primitiver $\cos(ax+b)$ et $\sin(ax+b)$ : facteur $\frac{1}{a}$ et signe $-$ pour le sinus.
  \item $\square$ Reconnaître la forme $u'\,u^r$ et ajuster la constante multiplicative.
  \item $\square$ Primitiver $\frac{a}{(cx+d)^r}$ en pensant au facteur $\frac{1}{c}$.
  \item $\square$ Traiter $\cos^n x$ et $\sin^n x$ : substitution $u=\sin x$ ou $\cos x$ si $n$ impair ; linéarisation si $n$ pair.
  \item $\square$ Utiliser les formes $\frac{u'}{u}\to\ln|u|$ et $u'\mathrm{e}^u\to\mathrm{e}^u$.
  \item $\square$ Vérifier \textbf{systématiquement} le résultat en dérivant.
  \item $\square$ Soigner la rédaction : préciser l'intervalle $I$, la constante $k\in\mathbb{R}$, et conclure par une phrase.
\end{itemize}
\end{aretenir}

\begin{savaistu}[Le saviez-vous ? — Histoire et applications]
\begin{itemize}
  \item Le symbole $\int$ (intégrale) a été introduit par Leibniz en 1675, comme un « S » allongé pour \emph{Summa} (somme). La notation $dx$ rappelle l'infinitésimal.
  \item Le lien entre primitive et aire sous la courbe (Théorème fondamental de l'analyse) a été établi indépendamment par Newton et Leibniz. C'est la pierre angulaire du calcul intégral.
  \item Au Cameroun, les primitives sont utilisées en \textbf{économie} (calcul du coût total à partir du coût marginal $C'(x)$), en \textbf{physique} (position à partir de la vitesse $v(t)$), en \textbf{télécommunications} (décibels, intégrales de signaux), et en \textbf{géomatique} (calcul d'aires sur cartes).
  \item La fonction $\ln|x|$ comme primitive de $\frac{1}{x}$ illustre l'unification des branches $x>0$ et $x<0$ : la dérivée de $\ln(-x)$ pour $x<0$ est bien $\frac{1}{x}$.
\end{itemize}
\end{savaistu}

\findecours

\end{document}
